\documentclass[11pt]{amsart}
\usepackage[margin=1.15in]{geometry}
\usepackage{amsmath,amssymb,amsthm}
\usepackage{booktabs}
\usepackage[hidelinks]{hyperref}

% ---------------------------------------------------------------------
% Theorem environments.  Environment names follow the companion paper
% (paper/weighted-races.tex); numbering is by section in this article.
% ---------------------------------------------------------------------
\newtheorem{letterthm}{Theorem}
\renewcommand{\theletterthm}{\Alph{letterthm}}
\newtheorem*{genthm}{General Theorem}
\newtheorem*{qthm}{Theorem for $q=3,4$ (computer-assisted)}
\newtheorem*{nonmonothm}{Non-monotonicity Theorem (computer-assisted)}
\newtheorem{theorem}{Theorem}[section]
\newtheorem{proposition}[theorem]{Proposition}
\newtheorem{lemma}[theorem]{Lemma}
\newtheorem{corollary}[theorem]{Corollary}
\theoremstyle{definition}
\newtheorem{remark}[theorem]{Remark}
\newtheorem{definition}[theorem]{Definition}
\numberwithin{equation}{section}
\raggedbottom

\newcommand{\eps}{\varepsilon}
\newcommand{\loglog}{\log\log}
\newcommand{\PP}{\mathbb P}
\newcommand{\EE}{\mathbb E}
\newcommand{\cE}{\mathcal E}
\newcommand{\hd}{\widehat\delta}
\newcommand{\re}{\operatorname{Re}}
\newcommand{\im}{\operatorname{Im}}
% visible marker for bibliographic details that still have to be checked
\newcommand{\todocite}[1]{\textsf{[TODO-cite: #1]}}

% ---------------------------------------------------------------------
% Outcome of the independent audit of the non-monotonicity certificate
% (monotonicity/nonmonotone-audit/AUDIT.md), kept as a macro so that it
% can be revised in one place.  There is no placeholder in this file.
% All certificate numbers quoted in Sections 10 and 11 are copied from
% the reports and result files named there; none is chosen here.
% ---------------------------------------------------------------------
\newcommand{\NonmonoAuditOutcome}{%
The certificate was audited by an independent AI agent, which
re-derived every inequality used and re-computed the results with code
of its own (report \texttt{AUDIT.md} in the directory
\texttt{monotonicity/nonmonotone-audit/}). Its findings: no break of
rigour; one gap, namely that the programs computing the enclosures
accept any zero file, so that the completeness of the list of zeros is
tied to them only by convention (deleting the first ordinate from the
file produces a false ``enclosure'' without any error); this is closed
by making the identity of the zero files a hypothesis, which is the
digest clause of \textup{(N5)}; and six minor items, none affecting a
statement made here. For the two files with the digests listed in
\textup{(N5)} the audit re-certified every window, reproduced the
counts $266$ and $228$ with a different contour and a different
criterion, and regenerated the file for conductor $232$ byte for byte
from scratch. Its own rigorous enclosures (different code, a different
integrator, a one-sided tail model, a different truncation point and
envelope constant, and also with only $100$, respectively $80$, zeros)
contain all twelve enclosures of Theorem~\ref{thm:nonmono} and on their
own give $\hd(\tfrac65)-\hd(\tfrac35)\ge5.705\times10^{-3}$ for
conductor $163$ and $\hd(1)-\hd(\tfrac{18}{25})\ge7.92\times10^{-4}$
for conductor $232$, as well as every other comparison of the theorem.
The audit covers the files pinned in its list \texttt{SHA256SUMS}.
After it, the programs were revised as it recommended: they now
record the digest of the zero file in every output line and refuse a
zero file that is not bound to a passing certificate record, and the
result file was regenerated. The regenerated enclosures are identical,
digit for digit, to the audited ones quoted in
Theorem~\ref{thm:nonmono}; the revised programs were not re-audited.
The audit is internal and AI-conducted, like all refereeing of this
article (Section~\ref{sec:gaps}).}

\begin{document}

\title[Chebyshev's bias as a function of the weight]{Monotonicity and
non-monotonicity of Chebyshev's bias in the weight}

\author{Raghav Sharma}
\email{raghavsharma1729@gmail.com}
\thanks{Draft of November 2026, not yet submitted. This work was
carried out with substantial AI assistance (Anthropic's Claude): the
derivations, the proofs, the numerical checks and the literature
review were AI-produced; what has and what has not been independently
verified is recorded in Sections~\ref{sec:numerics}
and~\ref{sec:gaps}; the author takes responsibility for the content.
Comments are very welcome. Contact: raghavsharma1729@gmail.com. Code
and data: \url{https://github.com/raghavs1729/weighted-prime-races}.}
\date{November 2026}
\subjclass[2020]{Primary 11N13; Secondary 11M26, 11Y35, 60E15}
\keywords{Chebyshev's bias, prime number races, weighted prime sums,
logarithmic density, coupling, Riesz means, explicit formula,
computer-assisted proof}

\begin{abstract}
Let $\chi$ be a real primitive Dirichlet character, $m$ a real
number, $M=m+\frac12$, and
$D_m(x)=-\sum_{p\le x}\chi(p)p^m$ (for $\chi=\chi_4$: the race of
primes $\equiv3$ against primes $\equiv1$ mod $4$ with weights
$p^m$). A companion paper showed that under GRH,
$L(\frac12,\chi)\ne0$ and LI the logarithmic density $\delta_\chi(m)$
of $\{x:D_m(x)>0\}$ exists for $m>-\frac12$, and conjectured that
$\delta_3$ and $\delta_4$ are strictly decreasing. Under the same
hypotheses we prove the following. (1) For every real primitive
$\chi$, $\delta_\chi$ is $C^1$; it is strictly decreasing for
$M\le M_0(\chi)$ (by a pathwise coupling; $M_0$ is the root of an
explicit function of the zeros) and for $M\ge40$ (by an envelope
criterion for the derivative, with an explicit lower bound and no zero
data); and $M_0(\chi)<40$. (2) For $q=3,4$, in a computer-assisted
sense, $\delta_q$ is strictly decreasing on all of
$(-\frac12,\infty)$: the range in between is covered by a
ball-arithmetic certificate of the envelope criterion on
$1\le M\le2000$. A second certificate by a different method and
internal audits confirm it on overlapping ranges; the inputs taken on
trust are listed. (3) Again computer-assisted, monotonicity
\emph{fails} in the intermediate range for the characters of conductor
$163$ and $232$. Independently: (A) under GRH alone,
$D_m(x)-x^{-(m'-m)}D_{m'}(x)$ is a Riesz-type mean of $D_{m'}$ with an
absolutely convergent explicit formula, and when an explicit margin
is positive every large $x$ with $D_{m'}(x)\ge0$ has $D_m(x)>0$; for
$q\in\{3,4\}$ and pairs with
$4M'\frac{\Lambda'}{\Lambda}(M'+\frac12,\chi)\le2$ its eventual
positivity is equivalent to GRH$(\chi)$; (C) weights $p^m(\log p)^b$
give the same limiting distribution for every real $b$. All
refereeing so far is internal and AI-conducted. This work was carried
out with substantial AI assistance (see the first footnote and
Section~\ref{sec:gaps}).
\end{abstract}

\maketitle

\section{Introduction}\label{sec:intro}

\noindent\textbf{Disclosure.} This article was produced with
substantial AI assistance (Anthropic's Claude): the derivations,
proofs, numerical checks and literature review are AI-produced.
Section~\ref{sec:numerics} records the numerical checks that were
actually run, and Section~\ref{sec:gaps} lists every place where a
statement is weaker than one might hope, depends on a computation, or
rests on a quoted result that was not re-derived here.

\subsection{Setting}\label{subsec:setting}

Throughout, $\chi$ is a real primitive Dirichlet character modulo
$q\ge3$ and $\chi^*:=-\chi$, so that $\chi^*(p)=+1$ on quadratic
nonresidues and $-1$ on residues. For real $m$ and $b$ put
\[
D_{m,b}(x)\;:=\;\sum_{p\le x}\chi^*(p)\,p^{m}(\log p)^{b},
\qquad
D_m:=D_{m,0},\qquad R_m:=D_{m,1}.
\]
For $\chi=\chi_4$, $D_m(x)=\sum_{p\le x,\,p\equiv3(4)}p^m
-\sum_{p\le x,\,p\equiv1(4)}p^m$; $m=0$ is Chebyshev's classical race.
We always write
\[
M=m+\tfrac12,\qquad M'=m'+\tfrac12,\qquad \kappa=m'-m=M'-M ,
\]
and restrict to $m>-\tfrac12$, i.e.\ $M>0$. The normalised races are
\begin{equation}\label{eq:normalised}
\cE_{m,b}(x)\;:=\;2M\,x^{-M}(\log x)^{1-b}\,D_{m,b}(x)\qquad(x>1),
\end{equation}
so that $\cE_{m,1}$ and $\cE_{m,0}$ are the functions $E_m$ and
$E_m'$ of the companion paper~\cite{WR}.

The nontrivial zeros of $L(s,\chi)$ are written $\rho=\beta+i\gamma$
and are always listed with multiplicity. GRH$(\chi)$ is the statement
that $\beta=\tfrac12$ for all of them. We put
\[
n_0=n_0(\chi):=\operatorname{ord}_{s=1/2}L(s,\chi)\ \ge0 ,
\]
and NCZ$(\chi)$ (``no central zero'') is the statement $n_0=0$; it is
unconditional for $q\in\{3,4\}$ (see \cite[\S2]{WR}). Under
GRH$(\chi)$, sums $\sum_{\gamma>0}$ run over the positive ordinates
with multiplicity, and LI$(\chi)$ is the linear independence over
$\mathbb Q$ of the multiset of positive ordinates (as in \cite{WR}, LI
does not include NCZ). Under GRH$(\chi)$ define, for $u,M,M'>0$,
\begin{align}
\sigma_M^2&:=\sum_{\gamma>0}\frac{8M^2}{M^2+\gamma^2},
&Q_\chi(u)&:=\sum_{\gamma>0}\frac{4u^3}{(u^2+\gamma^2)^{3/2}},
\label{eq:defQ}\\
\lambda_\chi(M,M')&:=\frac{1+2n_0}{MM'}
-\sum_{\gamma>0}\frac{4}{|M+i\gamma|\,|M'+i\gamma|},
&\eta_\chi(M,M')&:=\int_M^{M'}\frac{1-Q_\chi(u)}{u^2}\,du .
\label{eq:deflambda}
\end{align}
All four series converge because $\sum_{\gamma>0}\gamma^{-2}<\infty$.
The function $Q_\chi$ is continuous and strictly increasing from $0$
to $\infty$ (Lemma~\ref{lem:Q}); $M_0(\chi)$ denotes its unique root
of $Q_\chi(M_0)=1$.

Let $(\theta_\gamma)_{\gamma>0}$ be independent random variables,
uniform on $[0,2\pi)$, one for each positive ordinate, and
\begin{equation}\label{eq:XM}
X_M\;:=\;1+\sum_{\gamma>0}c_\gamma(M)\cos\theta_\gamma,
\qquad
c_\gamma(M):=\frac{4M}{\sqrt{M^2+\gamma^2}} .
\end{equation}
(In \cite{WR} this variable is written $X_m$ and $c_\gamma=2a_\gamma$;
its variance is $\sigma_M^2$.) Under GRH$(\chi)$ and NCZ$(\chi)$ put
\[
\delta_\chi(m)\;:=\;\PP\bigl(X_{m+1/2}>0\bigr),\qquad
\delta_4:=\delta_{\chi_4},\quad\delta_3:=\delta_{\chi_3}.
\]
Without LI$(\chi)$ this is a quantity attached to the \emph{model} law
of $X_M$ (one independent phase for each positive ordinate, counted
with multiplicity); it is the logarithmic density of the race only
under LI$(\chi)$, which in particular forces the zeros to be simple.
Statements below that assume GRH$(\chi)$ and NCZ$(\chi)$ but not
LI$(\chi)$ are statements about this model quantity. It is often
convenient to use $M$ rather than $m$ as the variable, and we write
\[
\hd_\chi(M)\;:=\;\delta_\chi\bigl(M-\tfrac12\bigr)\;=\;\PP(X_M>0)
\qquad(M>0),\qquad \hd_q:=\hd_{\chi_q}\quad(q=3,4);
\]
derivatives of $\hd_\chi$ are taken with respect to $M$.
The companion paper proves \cite[Theorem~1, Proposition~1]{WR}: under
GRH$(\chi)$ and NCZ$(\chi)$ the functions $u\mapsto\cE_{m,1}(e^u)$ and
$u\mapsto\cE_{m,0}(e^u)$ have a common limiting distribution $\nu_m$
for each $m>-\tfrac12$; under LI$(\chi)$ in addition, $\nu_m$ is the
law of $X_M$, it is absolutely continuous with support $\mathbb R$,
and the logarithmic density of $\{x\ge2:D_m(x)>0\}$ exists and equals
$\delta_\chi(m)\in(\tfrac12,1)$, a continuous function of $m$ with
$\delta_\chi(m)\to1$ as $m\to-\tfrac12^+$. It also proves
$\delta_\chi(m)\to\tfrac12$ as $m\to\infty$ with an asymptotic formula,
proves that $\delta_\chi$ is strictly decreasing for all sufficiently
large $m$ \cite[Theorem~4]{WR}, shows --- in a computer-assisted
sense described there \cite[Corollary~7 and Remark~13]{WR} --- that
$\delta_4$ is strictly decreasing for $m\ge16.82$ and $\delta_3$ for
$m\ge18.62$, and conjectures \cite[Conjecture~1]{WR} that $\delta_4$ and $\delta_3$ are
strictly decreasing on all of $(-\tfrac12,\infty)$. Remark~(c) after
that conjecture explains why the compact initial range resisted: the
family $X_M$ ``is not a dilation'', and peakedness comparisons fail
for the bimodal arcsine summands.

\subsection{Results}

\begin{letterthm}[inheritance of leads; GRH only]\label{thm:A}
Let $\chi$ be a real primitive character and assume GRH$(\chi)$
(neither LI$(\chi)$ nor NCZ$(\chi)$ is assumed). Let
$-\tfrac12<m<m'$, $\kappa=m'-m$, $M=m+\tfrac12$, $M'=m'+\tfrac12$.
\begin{enumerate}
\item[(i)] For all $x\ge2$,
\begin{align*}
D_m(x)-x^{-\kappa}D_{m'}(x)
&=\kappa\int_2^x t^{-\kappa-1}D_{m'}(t)\,dt\\
&=\frac{\kappa\,x^{M}}{\log x}\Bigl[\frac{1+2n_0}{2MM'}
+2\,\re\sum_{\gamma>0}\frac{x^{i\gamma}}{(M+i\gamma)(M'+i\gamma)}\Bigr]
+O\Bigl(\frac{x^{M}}{\log^2x}\Bigr),
\end{align*}
where the first equality is an exact identity, the series converges
absolutely, and the implied constant depends only on $\chi$, $m$ and
$m'$.
\item[(ii)] If $\lambda_\chi(M,M')>0$, then there is
$x_0=x_0(\chi,m,m')$ such that for all $x\ge x_0$
\[
D_m(x)-x^{-(m'-m)}D_{m'}(x)
\;=\;\kappa\int_2^x t^{-\kappa-1}D_{m'}(t)\,dt
\;\ge\;\frac{\kappa\,\lambda_\chi(M,M')}{4}\cdot\frac{x^{M}}{\log x}\;>\;0 .
\]
In particular, for every $x\ge x_0$: \ $D_{m'}(x)\ge0$ implies
$D_m(x)>0$, and $D_m(x)\le0$ implies $D_{m'}(x)<0$.
\item[(iii)] $\lambda_\chi(M,M')\ge
\bigl(2+4n_0-\sigma_M\sigma_{M'}\bigr)/(2MM')$. In particular
$\lambda_\chi(M,M')>0$ whenever $\sigma_M\sigma_{M'}<2+4n_0$, hence
whenever $\sigma_{M'}^2\le2$.
\end{enumerate}
\end{letterthm}

In words: for a pair of weights with positive margin, a lead of the
nonresidue class in the heavier race forces a lead in the lighter
race at the same $x$, and a lead (or tie) of the residue class in the
lighter race forces a lead in the heavier one. For $q=4$ and
$(m,m')=(0,1)$ one has $\sigma_{1/2}\sigma_{3/2}=0.4614<2$, so under
GRH$(\chi_4)$ every sufficiently large $x$ at which
$\sum_{p\le x,\,p\equiv3(4)}p\ge\sum_{p\le x,\,p\equiv1(4)}p$ satisfies
$\pi(x;4,3)>\pi(x;4,1)$. Corollaries~\ref{cor:densities},
\ref{cor:sharp}, \ref{cor:landau} and~\ref{cor:equiv} give the
consequences for upper and lower densities without LI, the sharpness
of the condition $\lambda_\chi\ge0$ under LI, a Landau-type converse,
and the equivalence with GRH$(\chi)$ for $q\in\{3,4\}$.

\begin{letterthm}[coupling monotonicity]\label{thm:B}
Let $\chi$ be a real primitive character and assume GRH$(\chi)$ and
NCZ$(\chi)$. Let $Q_\chi$, $M_0(\chi)$, $\eta_\chi$ be as in
\eqref{eq:defQ}--\eqref{eq:deflambda} and
$\delta_\chi(m)=\PP(X_{m+1/2}>0)$.
\begin{enumerate}
\item[(i)] For all $-\tfrac12<m<m'$ with $\eta:=\eta_\chi(M,M')\ge0$,
\[
\delta_\chi(m)-\delta_\chi(m')\;\ge\;
\PP\bigl(-\eta\,M'<X_{M'}\le0\bigr),
\]
and the right-hand side is strictly positive if $\eta>0$.
\item[(ii)] If $M<M'\le M_0(\chi)$ then $\eta_\chi(M,M')>0$.
Consequently $m\mapsto\delta_\chi(m)$ is strictly decreasing on
$\bigl(-\tfrac12,\,M_0(\chi)-\tfrac12\bigr]$.
\item[(iii)] Under LI$(\chi)$ in addition, $\delta_\chi(m)$ is the
logarithmic density of $\{x\ge2:D_{m,b}(x)>0\}$ for every real $b$
(in particular of $\{D_m>0\}$ and of $\{R_m>0\}$), so (i) and (ii)
are statements about the logarithmic densities of the races. (They
are not statements about nesting of the sets $\{x:D_m(x)>0\}$
themselves; for that see Theorem~\ref{thm:A} and
Remark~\ref{rem:Arange}(3).)
\item[(iv)] $M_0(\chi)$ satisfies
\[
\Bigl(\frac{\gamma_1}{2C(\chi)}\Bigr)^{1/3}
\;\le\;\Bigl(4\sum_{\gamma>0}\gamma^{-3}\Bigr)^{-1/3}
\;<\;M_0(\chi)\;<\;\kappa_0\,\gamma_1,
\qquad \kappa_0:=(4^{2/3}-1)^{-1/2}=0.8111\ldots,
\]
where $\gamma_1$ is the least positive ordinate and
$C(\chi)=\sum_{\gamma>0}2\gamma^{-2}=(\log\Lambda)''(\tfrac12,\chi)$;
moreover $Q_\chi(M)<\sigma_M^2/2$ for every $M>0$, so
$M_0(\chi)$ exceeds the root of $\sigma_M^2=2$.
\end{enumerate}
\end{letterthm}

Numerically (Section~\ref{sec:numerics}; floating-point evaluation
with a rigorous tail bound, not interval arithmetic)
$M_0(\chi_4)\in[3.688,3.690]$ and $M_0(\chi_3)\in[4.483,4.485]$, so
Theorem~\ref{thm:B} gives strict monotonicity of $\delta_4$ on
$-\tfrac12<m\le3.18$ and of $\delta_3$ on $-\tfrac12<m\le3.98$; the
weaker closed bound $(\gamma_1/2C(\chi))^{1/3}$ equals $2.68$ resp.\
$3.28$. In particular the whole interval $(-\tfrac12,-0.2)$, on which
\cite{WR} had no control beyond a Chernoff envelope, is covered.

\begin{letterthm}[logarithmic weights]\label{thm:C}
Let $\chi$ be a real primitive character, assume GRH$(\chi)$ and
NCZ$(\chi)$, and fix $m>-\tfrac12$ and $b\in\mathbb R$. Then:
\begin{enumerate}
\item[(i)] $\cE_{m,b}(x)=\cE_{m,1}(x)+O(1/\log x)$ for $x\ge3$, the
implied constant depending only on $\chi$, $m$ and $b$.
\item[(ii)] $u\mapsto\cE_{m,b}(e^u)$ has a limiting distribution, and
it is the measure $\nu_m$ of \cite[Theorem~1]{WR}; in particular it
does not depend on $b$.
\item[(iii)] Under LI$(\chi)$ in addition, the logarithmic density of
$\{x\ge2:D_{m,b}(x)>0\}$ exists and equals $\delta_\chi(m)$ for every
real $b$.
\end{enumerate}
\end{letterthm}

The cases $b=1$ and $b=0$ of (ii)--(iii) are \cite[Theorem~1,
Proposition~1]{WR}. Theorem~\ref{thm:C} answers, for $m>-\tfrac12$
and under the stated hypotheses, a question raised by G.~Martin in
correspondence with the author (August 2026): he suggested the weights
$p^a(\log p)^b$ and suspected that for fixed $a$ the density does not
depend on $b$.

The remaining results concern the function $\hd_\chi$ on the whole
half-line $M>0$. The first holds for every real primitive character.

\begin{genthm}
Let $\chi$ be a real primitive Dirichlet character and assume
GRH$(\chi)$ and $L(\tfrac12,\chi)\ne0$. Then:
\begin{enumerate}
\item[(i)] $\hd_\chi$ is continuously differentiable on $(0,\infty)$;
\item[(ii)] $\hd_\chi$ is strictly decreasing on $(0,M_0(\chi)]$;
\item[(iii)] for every $M\ge40$,
\[
-M\,\hd_\chi'(M)\;\ge\;\frac{0.3219}{\sqrt{2\pi}\,\sigma_M}\;>\;0 ,
\]
so $\hd_\chi$ is strictly decreasing on $[40,\infty)$;
\item[(iv)] $M_0(\chi)<40$.
\end{enumerate}
Under LI$(\chi)$ in addition, $\hd_\chi(M)$ is the logarithmic density
of $\{x\ge2:D_{m,b}(x)>0\}$, $m=M-\tfrac12$, for every real $b$.
\end{genthm}

Part (ii) is Theorem~\ref{thm:B}(ii). Parts (i) and (iii) are
Lemma~\ref{lem:diff} and Theorem~\ref{thm:E}; they rest on the
identity $-M\hd_\chi'(M)=f_M(1)+R(M)$, in which $f_M$ is the density
of $X_M-1$ and $R$ is a sign-indefinite term carrying small weights,
and on a criterion (Theorem~\ref{thm:envelope}) that bounds both terms
by integrals of explicit positive envelopes, with no cancellation to
track. The proof of (iii) uses no ordinate of any zero; it reduces to
four inequalities between elementary functions at explicit points,
which were evaluated in ball arithmetic (Section~\ref{sec:largeweight}).
The precise statement, with a sharper function of $\sigma_M$ in place
of the constant $0.3219$, is Theorem~\ref{thm:general}. Nothing is
asserted on the intermediate range $M_0(\chi)<M<40$, and the third
theorem below shows that nothing can be asserted there in general.

\begin{qthm}
Let $q\in\{3,4\}$ and assume GRH$(\chi_q)$. Then $\hd_q$ is strictly
decreasing on $(0,\infty)$, that is, $m\mapsto\delta_q(m)$ is strictly
decreasing on $(-\tfrac12,\infty)$ and \cite[Conjecture~1]{WR} holds
for this $q$, \textbf{provided} the following are correct:
\begin{itemize}
\item[\textup{(T1)}] the two files of ordinates pinned by their
SHA-256 digests in Section~\ref{subsec:inputs} ($511$ ordinates of
$L(s,\chi_4)$, $537$ of $L(s,\chi_3)$), together with the certificate
that each listed value is within $10^{-20}$ of an ordinate, that the
zeros are simple, and that the lists are complete up to their last
entry;
\item[\textup{(T2)}] the ball-arithmetic library FLINT/Arb and the
certifying programs that call it;
\item[\textup{(T3)}] the variance identity of the companion paper
\cite[eq.~(8)]{WR} (Proposition~\ref{prop:WRvar}).
\end{itemize}
Under LI$(\chi_q)$ in addition, the logarithmic density of
$\{x\ge2:D_{m,b}(x)>0\}$ is, for every real $b$, a strictly decreasing
function of $m\in(-\tfrac12,\infty)$.
\end{qthm}

The proof (Theorem~\ref{thm:q34}) has three overlapping pieces:
Theorem~\ref{thm:B} on $(0,M_0(\chi_q)]$, where $M_0(\chi_q)>1$; a
ball-arithmetic certificate of the envelope criterion on $765$
contiguous cells covering $1\le M\le2000$, for each modulus; and part
(iii) of the General Theorem on $M\ge40$. Two independent
confirmations are available on parts of the range and are \emph{not}
links of this chain (Section~\ref{subsec:confirm}): a second
certificate, by a different method (direct enclosure of the
oscillatory inversion integral), which has been audited and
independently re-proved on $3.6\le M\le17.4$ ($q=4$) and
$4.4\le M\le19.2$ ($q=3$) and extends, unaudited, to $M=320$ resp.\
$350$; and the analytic large-weight lemma of the companion paper for
$M\ge300$ \cite[Lemma~24]{WR}, which has been audited as well. In
particular \cite[Theorem~4 and Corollary~7]{WR} are not needed. The
word ``computer-assisted'' is meant literally: the certificates
establish the conclusion given (T1)--(T3); the programs were read,
re-run and mutation-tested by internal referees, but nothing in this
part is formally verified, and all refereeing so far is internal and
AI-conducted (Section~\ref{sec:gaps}).

\begin{nonmonothm}
Let $\chi_{-163}$ and $\chi_{-232}$ be the real primitive (odd)
characters of conductor $163$ and $232$, i.e.\ the Kronecker symbols
$(\frac{-163}{\cdot})$ and $(\frac{-232}{\cdot})$.
\begin{enumerate}
\item[(i)] Assume GRH$(\chi_{-163})$. Then
\[
\hd(\tfrac3{10})>\hd(\tfrac{57}{100})<\hd(\tfrac35)<\hd(\tfrac{117}{100})
>\hd(\tfrac65)>\hd(2)\qquad(\hd=\hd_{\chi_{-163}}),
\]
with $\hd(\tfrac65)-\hd(\tfrac35)\in[0.005705485,\,0.005705571]$.
\item[(ii)] Assume GRH$(\chi_{-232})$. Then
\[
\hd(\tfrac25)>\hd(\tfrac{71}{100})<\hd(\tfrac{18}{25})<\hd(1)<\hd(\tfrac{101}{100})
>\hd(\tfrac32)\qquad(\hd=\hd_{\chi_{-232}}),
\]
with $\hd(1)-\hd(\tfrac{18}{25})\in[0.000792990,\,0.000793135]$.
\item[(iii)] Consequently, under GRH$(\chi)$ and LI$(\chi)$ for
$\chi=\chi_{-163}$ or $\chi_{-232}$, the logarithmic density
$\delta_\chi(m)$ of $\{x\ge2:D_{m,b}(x)>0\}$ (any real $b$) is not a
monotone function of $m\in(-\tfrac12,\infty)$: for $\chi_{-163}$,
$\delta_\chi(-\tfrac15)>\delta_\chi(\tfrac1{10})<\delta_\chi(\tfrac7{10})
>\delta_\chi(\tfrac32)$.
\end{enumerate}
These statements are certified in ball arithmetic; they hold provided
the hypotheses and trusted inputs \textup{(N1)--(N5)} of
Section~\ref{sec:nonmono} are correct. (The only unproved mathematical
hypothesis is GRH for the character in question. The completeness of
the lists of zeros used is certified. FLINT/Arb, the certifying
programs and the identity of the two zero files, pinned by digests,
are trusted; the variance identity and the limiting-distribution
theorem of \cite{WR} are imported.)
\end{nonmonothm}

The precise statement, with all certified enclosures, is
Theorem~\ref{thm:nonmono}. Both increases occur above $M_0(\chi)$
($M_0(\chi_{-163})\approx0.1644$, $M_0(\chi_{-232})\approx0.2511$) and
far below $40$, so the intermediate range of the General Theorem
genuinely matters, and the theorem for $q=3,4$ is not an instance of a
general fact about real characters. An internal audit of this
certificate found no break of rigour and confirmed all enclosures with
independently written rigorous code; its one substantive finding, that
the identity of the zero files must be a stated hypothesis, is
incorporated (Section~\ref{sec:nonmono}). The mechanism is already visible in a model with two zeros
(Proposition~\ref{prop:twozero}).

\subsection{What is new and what is classical}

\emph{Theorem~\ref{thm:B}.} The mechanism is elementary and, as far as
we could determine, new in this context. Dividing by the weight,
\[
\frac{X_M}{M}\;=\;\frac1M+\sum_{\gamma>0}
\frac{4\cos\theta_\gamma}{\sqrt{M^2+\gamma^2}} .
\]
The bias term $1/M$ decreases in $M$ at rate $M^{-2}$; each
oscillating coefficient also decreases in $M$, at rate
$4M(M^2+\gamma^2)^{-3/2}=O(\gamma^{-3})$, and these rates are
\emph{absolutely summable} over the zeros although the coefficients
themselves are not. Hence for a \emph{fixed} phase vector the
difference $X_M/M-X_{M'}/M'$ is bounded below by the deterministic
quantity $\eta_\chi(M,M')$ (Lemma~\ref{lem:det}), which is positive
as long as $Q_\chi<1$. Thus the family \emph{is} a dilation up to an
absolutely summable perturbation that is dominated by the gain in
bias, and the events $\{X_M>0\}$ are nested almost surely; no unimodality or
peakedness is needed (Remark~\ref{rem:dilation}). We did not find a
monotonicity statement for a Rubinstein--Sarnak density \cite{RS94}
in a continuous parameter, or a coupling argument of this kind, in the
literature we searched (the search is recorded in the repository;
it is of course not a proof of absence).

\emph{Theorem~\ref{thm:A}.} The method is classical: after a
smoothing, a positive bias term is compared with an absolutely
convergent sum over zeros, for \emph{all} large $x$. This is the
mechanism of Hardy--Littlewood and Landau \cite{HL18,Lan18} for the
Abel means $\sum_p\chi_4(p)e^{-p/x}$ (their divergence to $-\infty$ is
equivalent to GRH$(\chi_4)$), of Fujii \cite{Fuj88} and
Platt--Trudgian \cite{PT19} for the kernels $e^{-(p/x)^\alpha}$, of
Johnston \cite{Joh23} for $\int_2^x(\pi(t)-\mathrm{li}(t))\,dt$, of
Washington \cite[Theorem~3]{Was24} for integrated sums of powers of
primes, and of Zhao \cite{Zhao25} for the error terms in Mertens'
theorems; see Remark~\ref{rem:literature}. The identity
$D_m(x)-x^{-\kappa}D_{m'}(x)=\kappa\int_2^xt^{-\kappa-1}D_{m'}(t)\,dt$
is one line of partial summation, and its left-hand side is literally
a Riesz typical mean of order one. What we did not find in the
literature searched is its \emph{interpretation} as a link between two
races of different weights --- the weight shift itself is the
smoothing --- and the resulting statement that leads are inherited
across weights at the same $x$, with its consequence for upper and
lower densities without LI (Corollary~\ref{cor:densities}). The
equivalence with GRH (Corollary~\ref{cor:equiv}) is a Riesz-mean
cousin of the Hardy--Littlewood--Landau equivalence, of modest
novelty.

Sums of $k$-th powers of primes with $k>0$ have been studied as
approximations to $\pi(x^{k+1})$ by Gerard--Washington \cite{GW18},
Washington \cite{Was24} and, in arithmetic progressions, by Boran,
Byun, Li, Miller and Reyes \cite{BBLMR24}; these works compare one
weighted sum with an unweighted count and do not race two residue
classes against each other with a common weight, nor compute
logarithmic densities. For decaying weights see
\cite{AK23,Sheth24,SK25,Hay25,Aym22} and the discussion in \cite{WR}.

\emph{The General Theorem, the theorem for $q=3,4$ and the
non-monotonicity theorem.} The derivative identity of
Lemma~\ref{lem:diff} is the rigorous form of a formal computation
(Remark~\ref{rem:infinitesimal}); identities of this kind are
standard for sums of independent variables, and the companion paper
differentiates the inversion integral for large weights
\cite[Lemma~21]{WR}. What we did not find elsewhere is the envelope
criterion of Theorem~\ref{thm:envelope}, in which the main term is the
density itself and is bounded below by a Gaussian-type integral, the
remainder is bounded above with the same envelopes, and the ratio of
the two bounds stays below $\tfrac12$ in the limit --- so that the
criterion has room on the whole range $1\le M\le2000$ for $q=3,4$ and
needs only the twelve lowest ordinates. The large-weight part
(Theorem~\ref{thm:E}) replaces, by an independent and uniform
argument, the large-weight results of \cite{WR}, which are for
unspecified or character-dependent thresholds. We did not find, in
the literature searched, an instance in which a Rubinstein--Sarnak-type
density is shown to be non-monotone in a continuous parameter, as
Theorem~\ref{thm:nonmono} shows conditionally and with computer
assistance for conductors $163$ and $232$; but we did not search
specifically for one (Section~\ref{sec:gaps}). These three results were obtained after the
first version of this article; their proofs are transcribed from
working notes in the repository, where they were refereed internally
(Appendix~\ref{sec:review}).

\emph{Theorem~\ref{thm:C}} is a short partial summation once a local
mean-square bound (Lemma~\ref{lem:localms}) is available; the cases
$b\in\{0,1\}$ are in \cite{WR}, and Bhattacharya--Martin--Simpson
\cite{BMS26} treat correlations of error terms for reciprocal-type
weights. We did not find the statement for general real $b$ in the
literature searched.

\emph{Classical inputs} used without proof: the truncated explicit
formula for $\psi(x,\chi)$ in the form \cite[Lemma~1]{WR} (cf.\
\cite[Ch.~19]{Dav00}); the prime number theorem with error
$O(x/\log^2x)$; Landau's theorem on Mellin transforms of functions of
constant sign \cite[Lemma~15.1]{MV07}; Kronecker's theorem; the
limiting-distribution theorems of Akbary--Ng--Shahabi \cite{ANS14} in
the form in which they are applied in \cite{WR}; and standard facts on
weak convergence; and, for Sections~\ref{sec:envelope}
and~\ref{sec:largeweight}, the facts about the Bessel function $J_0$
collected in Section~\ref{subsec:bessel}, three of which are quoted
from \cite{WR}.

\subsection{Organisation}
Section~\ref{sec:prelim} collects notation and the quoted inputs.
Section~\ref{sec:det} contains the deterministic lemma and the
properties of $Q$. Section~\ref{sec:thmB} proves
Theorem~\ref{thm:B}, its sharpness, and the bounds for $M_0$.
Section~\ref{sec:formula} proves the asymptotic formula of
Theorem~\ref{thm:A}(i); Section~\ref{sec:thmA} proves the rest of
Theorem~\ref{thm:A} and its corollaries. Section~\ref{sec:thmC}
proves Theorem~\ref{thm:C}. Section~\ref{sec:envelope} proves the
derivative identity and the envelope criterion;
Section~\ref{sec:largeweight} proves the uniform large-weight theorem
and the General Theorem. Section~\ref{sec:main} states and proves the
theorem for $q=3,4$, with the certificate, the trusted inputs and the
independent confirmations. Section~\ref{sec:nonmono} contains the
two-zero model and the non-monotonicity theorem.
Section~\ref{sec:numerics} reports numerical sanity checks,
Section~\ref{sec:gaps} lists gaps and caveats, and
Appendix~\ref{sec:review} records the review history.

\section{Preliminaries and quoted inputs}\label{sec:prelim}

Implied constants may depend on $\chi$ (equivalently on $q$) and on
the fixed exponents $m,m',b$, never on $x$, $t$, $u$, $T$ or on the
zero $\rho$ under consideration. We write
$\psi(x,\chi)=\sum_{n\le x}\Lambda(n)\chi(n)$ and
\[
\psi_1(x,\chi):=\int_0^x\psi(t,\chi)\,dt
=\sum_{n\le x}\Lambda(n)\chi(n)(x-n) .
\]

\subsection{Zeros}\label{subsec:zeros}
We use the following standard facts \cite[\S\S12, 16]{Dav00}. The
nontrivial zeros satisfy $0<\beta<1$ (so $\rho\ne0$), they are
symmetric under $\rho\mapsto\bar\rho$ because $\chi$ is real, and
\begin{equation}\label{eq:zerosum}
\sum_\rho\frac{1}{|\rho|^2}<\infty,
\qquad
N(T+1)-N(T)\ll\log(T+2)\quad(T\ge0),
\qquad
N(T)\sim\frac{T}{2\pi}\log T ,
\end{equation}
where $N(T)$ counts zeros with $0<\gamma\le T$. Consequently, under
GRH$(\chi)$,
\begin{equation}\label{eq:gammasums}
\sum_{\gamma>0}\gamma^{-2}<\infty,\qquad
\sum_{\gamma>0}\gamma^{-2}\log(\gamma+3)<\infty,\qquad
\sum_{\gamma>0}\gamma^{-1}=\infty ,
\end{equation}
and there are infinitely many distinct positive ordinates.

\subsection{Inputs quoted from the companion paper}
We quote four results of \cite{WR}; the theorem numbers refer to the
August 2026 version, and we also give the name each result carries
there.

\begin{lemma}[{\cite[Lemma~1]{WR}, ``truncated explicit formula, all
truncations''}]\label{lem:WRclassical}
Unconditionally, for all $x\ge2$ and $T\ge2$,
\[
\psi(x,\chi)=-\sum_{|\gamma|\le T}\frac{x^\rho}{\rho}
+O\Bigl(\frac{x\log^2(xT)}{T}+\log x\Bigr),
\]
with no inequality assumed between $T$ and $x$.
\end{lemma}

\begin{lemma}[{\cite[proofs of Theorem~1(i) and of
Proposition~1]{WR}}]\label{lem:WRexp}
Assume GRH$(\chi)$ and NCZ$(\chi)$ and fix $m>-\tfrac12$. Put
$r_\gamma(M):=4M/(M+i\gamma)$, so that $|r_\gamma(M)|=c_\gamma(M)$.
For $b\in\{0,1\}$, all $u\ge2$ and all $T\ge2$,
\[
\cE_{m,b}(e^u)\;=\;1+\re\sum_{0<\gamma\le T}r_\gamma(M)\,e^{i\gamma u}
+\cE^{(b)}(u,T),
\]
where, with some $c>0$ and $\kappa_m:=\min(M,\tfrac12)$,
\[
\bigl|\cE^{(b)}(u,T)\bigr|\;\ll\;
\frac{u\,e^{u/2}\log^2(e^uT)}{T}+\frac{1-b}{u}+u^2e^{-\kappa_mu}
+u\,e^{-u/6}+u\,e^{-c\sqrt u}.
\]
In particular
$\lim_{Y\to\infty}\frac1Y\int_2^Y|\cE^{(b)}(u,e^Y)|^2\,du=0$.
\end{lemma}

(The displays in \cite{WR} are slightly sharper for $b=1$; the common
bound above is all we use. The last assertion is verified in
\cite{WR}: at $T=e^Y$ and $2\le u\le Y$ the first term is
$\ll Y^3e^{u/2-Y}$, and the remaining terms are fixed functions of $u$
whose squares have Ces\`aro mean tending to $0$.)

\begin{theorem}[{\cite[Theorem~1 and Proposition~1]{WR}}]
\label{thm:WRinterp}
Assume GRH$(\chi)$ and NCZ$(\chi)$ and fix $m>-\tfrac12$. For
$b\in\{0,1\}$ the function $u\mapsto\cE_{m,b}(e^u)$ has a limiting
distribution $\nu_m$, the same for both values of $b$. Under
LI$(\chi)$, $\nu_m$ is the law of $X_M$; it is absolutely continuous
with support $\mathbb R$; the logarithmic densities of
$\{x\ge2:D_m(x)>0\}$ and of $\{x\ge2:R_m(x)>0\}$ exist and equal
$\delta_\chi(m)=\PP(X_M>0)\in(\tfrac12,1)$; and $\delta_\chi$ is
continuous on $(-\tfrac12,\infty)$.
\end{theorem}

The proofs in \cite{WR} obtain existence from
\cite[Corollary~1.3(a)]{ANS14} applied to the expansion of
Lemma~\ref{lem:WRexp}, and the identification under LI from
\cite[Theorem~1.9]{ANS14}. We use those two results of \cite{ANS14}
in exactly the same way once more, in Lemma~\ref{lem:joint}. (A
referee of this article checked the statements of
\cite[Corollary~1.3(a), Theorems~1.4 and~1.9]{ANS14} against the
arXiv version 1306.1657 of that paper and confirmed that they are as
used here and in \cite{WR}; for complex coefficients $r_n$ the Bessel
factors of Theorem~1.9 carry $|r_n|$.)

\begin{proposition}[{\cite[eq.~(8) and \S2]{WR}}, variance
identity]\label{prop:WRvar}
Assume GRH$(\chi)$ and NCZ$(\chi)$, let $\chi(-1)=(-1)^a$ and
$\Lambda(s,\chi)=(q/\pi)^{(s+a)/2}\Gamma\bigl(\tfrac{s+a}2\bigr)
L(s,\chi)$. Then for every $M>0$
\[
\sigma_M^2=4M\,\frac{\Lambda'}{\Lambda}\bigl(M+\tfrac12,\chi\bigr),
\qquad
C(\chi):=\sum_{\gamma>0}\frac{2}{\gamma^2}
=(\log\Lambda)''\bigl(\tfrac12,\chi\bigr).
\]
\end{proposition}

(In \cite{WR}, Proposition~2 states and proves the first identity for
$\chi_4$; the identity for a general real primitive $\chi$ is eq.~(8)
there, in the proof of Proposition~3, where it is observed that the
proof of Proposition~2 applies verbatim. The closed form of $C(\chi)$
is in \cite[\S2]{WR}.)

\subsection{Densities and limiting distributions}
For a measurable set $S\subset[1,\infty)$ the upper and lower
logarithmic densities are the $\limsup$ and $\liminf$, as
$X\to\infty$, of $(\log X)^{-1}\int_{S\cap[1,X]}dt/t$; the upper and
lower natural densities are those of $X^{-1}\int_{S\cap[1,X]}dt$. When
upper and lower values agree we speak of the logarithmic (natural)
density. A locally bounded measurable function
$\phi:[0,\infty)\to\mathbb R^d$ \emph{has the limiting distribution}
$\mu$ (a Borel probability measure on $\mathbb R^d$) if
$\frac1Y\int_0^YF(\phi(u))\,du\to\int F\,d\mu$ as $Y\to\infty$ for
every bounded continuous $F:\mathbb R^d\to\mathbb R$; equivalently,
if the probability measures $\mu_Y(B):=\frac1Y\,
\mathrm{meas}\{u\in[0,Y]:\phi(u)\in B\}$ converge weakly to $\mu$.

\begin{lemma}\label{lem:stability}
Let $\phi$ have the limiting distribution $\mu$.
\begin{enumerate}
\item[(a)] If $\psi:[0,\infty)\to\mathbb R^d$ is locally bounded and
measurable and $\psi(u)-\phi(u)\to0$ as $u\to\infty$, then $\psi$ has
the limiting distribution $\mu$.
\item[(b)] If $B\subset\mathbb R^d$ is a Borel set with
$\mu(\partial B)=0$, then $\mu_Y(B)\to\mu(B)$; if $B$ is open, then
$\liminf_{Y\to\infty}\mu_Y(B)\ge\mu(B)$.
\end{enumerate}
\end{lemma}

\begin{proof}
(b) is the portmanteau theorem \cite[Theorem~2.1]{Bil99}, applied
along an arbitrary sequence $Y_n\to\infty$. (a) For a bounded
Lipschitz function $F$ with constant $L_F$,
$|F(\psi(u))-F(\phi(u))|\le\min(2\|F\|_\infty,L_F|\psi(u)-\phi(u)|)
\to0$, so $\frac1Y\int_0^YF(\psi)\,du$ and $\frac1Y\int_0^YF(\phi)\,du$
have the same limit $\int F\,d\mu$. Convergence for all bounded
uniformly continuous $F$ is equivalent to weak convergence
\cite[Theorem~2.1]{Bil99}, and bounded Lipschitz functions suffice for
that implication (the proof of the theorem uses only the Lipschitz
functions $x\mapsto\max(0,1-\mathrm{dist}(x,F_0)/\eps)$ for closed
$F_0$).
\end{proof}

\begin{lemma}[full support]\label{lem:fullsupport}
Let $c_n>0$ with $c_n\to0$, $\sum_nc_n=\infty$ and
$\sum_nc_n^2<\infty$, let $\beta_n\in\mathbb R$ be arbitrary, and let
$\theta_n$ be independent and uniform on $[0,2\pi)$. Then
$W:=\sum_nc_n\cos(\theta_n-\beta_n)$ converges almost surely, and
$\PP(W\in J)>0$ for every nonempty open interval $J$.
\end{lemma}

\begin{proof}
The summands are independent, centred, with variances $c_n^2/2$, so
the series converges almost surely by Kolmogorov's convergence
criterion (see e.g.\ \cite[\S VII.8]{Fel71}). Since
$\theta_n-\beta_n$ reduced mod $2\pi$ are again independent and
uniform, we may take $\beta_n=0$; and since $\theta\mapsto\theta+\pi$
preserves the uniform law and flips the sign of $W$, it suffices to
treat $J=(v-\eps,v+\eps)$ with $v\ge0$. This is the argument of
\cite[proof of Theorem~1(ii), ``Full support'']{WR}, which uses only
the three properties of $(c_n)$ assumed here. Choose $N_0$ with
$c_n<\eps/4$ for $n>N_0$, then $N_1\ge N_0$ minimal with
$S:=\sum_{N_0<n\le N_1}c_n\ge v$ (so $v\le S<v+\eps/4$; if $v=0$ take
$N_1=N_0$), then $N\ge N_1$ with $\sum_{n>N}c_n^2/2\le\eps^2/64$. By
Chebyshev's inequality the tail $W_N:=\sum_{n>N}c_n\cos\theta_n$
satisfies $\PP(|W_N|<\eps/4)\ge\tfrac34$. Let
$\delta:=\eps/(4+4\sum_{n\le N}c_n)$ and let $A$ be the event that
$|\cos\theta_n-1|<\delta$ for $N_0<n\le N_1$ and $|\cos\theta_n|<\delta$
for the remaining $n\le N$. Then $\PP(A)>0$, $A$ is independent of
$W_N$, and on $A$ the head $\sum_{n\le N}c_n\cos\theta_n$ lies within
$\eps/4$ of $S$. On $A\cap\{|W_N|<\eps/4\}$, which has probability
$\ge\tfrac34\PP(A)>0$, we get $|W-v|<\tfrac34\eps$.
\end{proof}

By \eqref{eq:gammasums}, Lemma~\ref{lem:fullsupport} applies to
$c_n=c_{\gamma_n}(M)$ for every $M>0$ and to every tail
$(c_{\gamma_n}(M))_{n>J}$: indeed $c_\gamma(M)\le4M/\gamma$ and
$c_\gamma(M)\ge2\sqrt2\,M/\gamma$ for $\gamma\ge M$.

\section{The deterministic lemma}\label{sec:det}

This section contains no number theory and no probability. Throughout
it, $0<\gamma_1\le\gamma_2\le\cdots$ is an arbitrary infinite sequence
of positive real numbers with
\begin{equation}\label{eq:hypgamma}
\sum_{j\ge1}\gamma_j^{-2}<\infty .
\end{equation}
Note that then $\sum_j\gamma_j^{-3}\le\gamma_1^{-1}\sum_j\gamma_j^{-2}
<\infty$ as well.

\begin{lemma}[the function $Q$]\label{lem:Q}
For $u>0$ let
\[
Q(u)\;:=\;\sum_{j\ge1}\frac{4u^3}{(u^2+\gamma_j^2)^{3/2}} .
\]
\begin{enumerate}
\item[(i)] $Q$ is finite, continuous and strictly increasing on
$(0,\infty)$, with $0<Q(u)<4u^3\sum_j\gamma_j^{-3}$; in particular
$Q(u)\to0$ as $u\to0^+$. Moreover $Q(u)\to\infty$ as $u\to\infty$.
\item[(ii)] There is a unique $M_0>0$ with $Q(M_0)=1$; $Q<1$ on
$(0,M_0)$ and $Q>1$ on $(M_0,\infty)$.
\item[(iii)] For $0<M<M'$ put
$\Delta_j(M,M'):=4(M^2+\gamma_j^2)^{-1/2}-4(M'^2+\gamma_j^2)^{-1/2}$.
Then $\Delta_j>0$ and
\[
\sum_{j\ge1}\Delta_j(M,M')\;=\;\int_M^{M'}\frac{Q(u)}{u^2}\,du\;<\;\infty .
\]
\end{enumerate}
\end{lemma}

\begin{proof}
(i) Write $Q(u)=\sum_jg(u/\gamma_j)$ with
$g(t):=4t^3(1+t^2)^{-3/2}=4(1+t^{-2})^{-3/2}$. The function $g$ is
continuous and strictly increasing on $(0,\infty)$, with
$0<g(t)<\min(4,4t^3)$ and $g(t)\to4$ as $t\to\infty$. Hence
$0<g(u/\gamma_j)<4u^3\gamma_j^{-3}$, and on every interval $(0,U]$ the
series is dominated termwise by the convergent series
$4U^3\sum_j\gamma_j^{-3}$; by the Weierstrass test $Q$ is finite and
continuous, and $Q(u)<4u^3\sum_j\gamma_j^{-3}\to0$ as $u\to0^+$. A sum
of strictly increasing functions is strictly increasing. For every
$N$, $Q(u)\ge\sum_{j\le N}g(u/\gamma_j)\to4N$ as $u\to\infty$, so
$Q(u)\to\infty$.

(ii) follows from (i) by the intermediate value theorem and strict
monotonicity.

(iii) Since $\frac{d}{du}\bigl[-4(u^2+\gamma^2)^{-1/2}\bigr]
=4u(u^2+\gamma^2)^{-3/2}>0$,
\[
\Delta_j(M,M')=\int_M^{M'}\frac{4u}{(u^2+\gamma_j^2)^{3/2}}\,du>0 ,
\]
and by Tonelli's theorem
$\sum_j\Delta_j=\int_M^{M'}\sum_j4u(u^2+\gamma_j^2)^{-3/2}\,du
=\int_M^{M'}Q(u)u^{-2}\,du$, which is finite because $Q$ is continuous
on $[M,M']$.
\end{proof}

\begin{lemma}[deterministic coupling inequality]\label{lem:det}
Let $(\gamma_j)$ satisfy \eqref{eq:hypgamma} and let
$x_j\in[-1,1]$ be real numbers such that the series
$\sum_jx_j/\gamma_j$ converges (conditionally, in the given order).
\begin{enumerate}
\item[(i)] For every $M>0$ the series
\[
\varphi(M)\;:=\;\frac1M+\sum_{j\ge1}\frac{4x_j}{\sqrt{M^2+\gamma_j^2}}
\]
converges (in the given order).
\item[(ii)] For $0<M<M'$,
\[
\varphi(M)-\varphi(M')\;=\;\Bigl(\frac1M-\frac1{M'}\Bigr)
+\sum_{j\ge1}x_j\,\Delta_j(M,M'),
\]
the last series being absolutely convergent, and
\begin{equation}\label{eq:detineq}
\int_M^{M'}\frac{1+Q(u)}{u^2}\,du\;\ge\;
\varphi(M)-\varphi(M')\;\ge\;
\eta(M,M'):=\int_M^{M'}\frac{1-Q(u)}{u^2}\,du .
\end{equation}
\item[(iii)] If $0<M<M'\le M_0$, then $\eta(M,M')>0$; hence $\varphi$
is strictly decreasing on $(0,M_0]$. More generally, whenever
$\eta(M,M')\ge0$: \ $\varphi(M')>0$ implies $\varphi(M)>0$, and
$\varphi(M')>-\eta(M,M')$ implies $\varphi(M)>0$.
\end{enumerate}
\end{lemma}

\begin{proof}
(i) For $\gamma>0$ and $s:=\sqrt{M^2+\gamma^2}\ge\gamma$ we have
$s-\gamma=M^2/(s+\gamma)\le M^2/(2\gamma)$, hence
\begin{equation}\label{eq:gapbound}
0\;\le\;\frac1\gamma-\frac1{\sqrt{M^2+\gamma^2}}
=\frac{s-\gamma}{\gamma s}\;\le\;\frac{M^2}{2\gamma^3}.
\end{equation}
Therefore $\sum_jx_j\bigl(4\gamma_j^{-1}-4(M^2+\gamma_j^2)^{-1/2}\bigr)$
converges absolutely, and subtracting it from the convergent series
$\sum_j4x_j/\gamma_j$ gives the convergence of the series defining
$\varphi(M)$.

(ii) Subtracting the two convergent series termwise gives the
identity; the series $\sum_jx_j\Delta_j$ converges absolutely because
$|x_j|\le1$ and $\sum_j\Delta_j<\infty$ by Lemma~\ref{lem:Q}(iii). The
same two facts give
$-\sum_j\Delta_j\le\sum_jx_j\Delta_j\le\sum_j\Delta_j$, and
$\frac1M-\frac1{M'}=\int_M^{M'}u^{-2}\,du$; inserting
$\sum_j\Delta_j=\int_M^{M'}Q(u)u^{-2}\,du$ yields \eqref{eq:detineq}.

(iii) If $M'\le M_0$ then $1-Q(u)>0$ for all $u\in[M,M')$ by
Lemma~\ref{lem:Q}(ii), so $\eta(M,M')>0$ and
$\varphi(M)>\varphi(M')$. The remaining assertions are immediate from
$\varphi(M)\ge\varphi(M')+\eta(M,M')$.
\end{proof}

\begin{remark}\label{rem:detB}
If the constant $1$ in $\varphi$ is replaced by a constant $B>0$
(that is, $\varphi_B(M)=B/M+\sum_j4x_j(M^2+\gamma_j^2)^{-1/2}$), the
same proof gives
$\varphi_B(M)-\varphi_B(M')\ge\int_M^{M'}(B-Q(u))u^{-2}\,du$. For the
races of this paper $B=1+2n_0(\chi)$; we use only $B=1$.
\end{remark}

\section{Coupling monotonicity: proof of Theorem~\ref{thm:B}}
\label{sec:thmB}

Assume GRH$(\chi)$ and NCZ$(\chi)$ throughout this section, let
$0<\gamma_1\le\gamma_2\le\cdots$ be the positive ordinates of
$L(s,\chi)$ with multiplicity, and write $\theta_j$ for
$\theta_{\gamma_j}$. By \eqref{eq:gammasums} the sequence satisfies
\eqref{eq:hypgamma}, and the functions $Q$, $\eta$, $M_0$ of
Section~\ref{sec:det} are the $Q_\chi$, $\eta_\chi$, $M_0(\chi)$ of
the introduction.

\subsection{The coupling}
Let
\[
\Omega_0\;:=\;\Bigl\{\sum_{j\ge1}\frac{\cos\theta_j}{\gamma_j}
\ \text{converges}\Bigr\}.
\]
The summands are independent and centred with variances
$1/(2\gamma_j^2)$, which are summable, so $\PP(\Omega_0)=1$ by
Kolmogorov's criterion (as in Lemma~\ref{lem:fullsupport}). On
$\Omega_0$, Lemma~\ref{lem:det}(i) with $x_j=\cos\theta_j$ shows that
the series \eqref{eq:XM} converges \emph{for every $M>0$
simultaneously}; we use this to define all the variables $X_M$,
$M>0$, on the one event $\Omega_0$ (and, say, $X_M:=1$ off
$\Omega_0$). For each fixed $M$ this is the variable $X_M$ of
\eqref{eq:XM}, and $X_M=M\varphi(M)$ on $\Omega_0$ in the notation of
Lemma~\ref{lem:det}.

\begin{proposition}[pathwise nesting]\label{prop:nest}
On $\Omega_0$, for all $0<M<M'$,
\begin{equation}\label{eq:pathwise}
\frac{X_M}{M}-\frac{X_{M'}}{M'}\;\ge\;\eta_\chi(M,M') .
\end{equation}
Consequently, if $\eta:=\eta_\chi(M,M')\ge0$, then on $\Omega_0$
\[
\{X_{M'}>0\}\ \cup\ \{-\eta M'<X_{M'}\le0\}\ \subset\ \{X_M>0\},
\]
the union being disjoint, and
\begin{equation}\label{eq:gapB}
\PP(X_M>0)-\PP(X_{M'}>0)\;\ge\;\PP\bigl(-\eta M'<X_{M'}\le0\bigr).
\end{equation}
\end{proposition}

\begin{proof}
\eqref{eq:pathwise} is \eqref{eq:detineq} with $x_j=\cos\theta_j$. If
$X_{M'}>-\eta M'$ then $X_M/M\ge X_{M'}/M'+\eta>0$. The two events on
the left are disjoint and both contained in $\{X_M>0\}$ (up to the
null set $\Omega_0^c$), which gives \eqref{eq:gapB}.
\end{proof}

\begin{proof}[Proof of Theorem~\ref{thm:B}(i)--(iii)]
(i) is \eqref{eq:gapB}. If $\eta>0$, the interval $(-\eta M',0)$ is
nonempty and open, and $X_{M'}-1=\sum_jc_{\gamma_j}(M')\cos\theta_j$
has full support by Lemma~\ref{lem:fullsupport}; hence
$\PP(-\eta M'<X_{M'}\le0)>0$.

(ii) is Lemma~\ref{lem:det}(iii) together with (i).

(iii) Under LI$(\chi)$, Theorem~\ref{thm:WRinterp} identifies
$\delta_\chi(m)$ with the logarithmic density of $\{D_m>0\}$ and of
$\{R_m>0\}$, and Theorem~\ref{thm:C}(iii) (proved in
Section~\ref{sec:thmC}, independently of the present section) extends
this to $\{D_{m,b}>0\}$ for every real $b$.
\end{proof}

\begin{remark}[why ``the family is not a dilation'' is circumvented]
\label{rem:dilation}
Remark~(c) after \cite[Conjecture~1]{WR} observes that the centred
variables $V_M=X_M-1$ do not form a dilation family --- the ratio
$c_\gamma(M')/c_\gamma(M)$ depends on $\gamma$ --- and that the
classical peakedness comparisons of Birnbaum and Proschan
\cite{Bir48,Pro65} need unimodal summands, which the arcsine
variables $c_\gamma\cos\theta_\gamma$ are not.
Proposition~\ref{prop:nest} avoids both obstacles. First, it compares
$X_M$, bias included, rather than $V_M$: by \eqref{eq:pathwise}, on
$\Omega_0$,
\[
X_M\;=\;\frac{M}{M'}\,X_{M'}\;+\;M\,\Bigl(\frac{X_M}{M}
-\frac{X_{M'}}{M'}\Bigr),
\qquad
M\,\Bigl(\frac{X_M}{M}-\frac{X_{M'}}{M'}\Bigr)\;\ge\;M\,\eta_\chi(M,M'),
\]
so $X_M$ \emph{is} a dilation of $X_{M'}$ by the factor $M/M'<1$, up
to a remainder which is an absolutely convergent series in the zeros
(its coefficients $M\Delta_\gamma$ are $O(\gamma^{-3})$) and whose
worst case is dominated by the gain $M(1/M-1/M')$ in bias as long as
$\eta_\chi\ge0$. The $\gamma$-dependence of the ratios
$c_\gamma(M')/c_\gamma(M)$ is entirely contained in this remainder.
Second, it is a statement about sample paths under one particular
coupling of the two laws (the same phases for both weights), not a
comparison of distribution functions; since $\delta_\chi$ depends
only on the marginal laws, any coupling may be used, and nesting of
the events $\{X_M>0\}$ requires no shape information about the
summands.
\end{remark}

\begin{remark}[infinitesimal form; formal]\label{rem:infinitesimal}
Formally differentiating $S_M:=X_M-1$ in $M$ gives
$\partial_MS_M=(S_M-Z_M)/M$ with
$Z_M:=\sum_\gamma\frac{M^2}{M^2+\gamma^2}c_\gamma(M)\cos\theta_\gamma$,
an absolutely convergent series with $|Z_M|\le Q_\chi(M)$, and hence,
formally, $\delta_\chi'(m)=-M^{-1}f_M(-1)\,\EE[1+Z_M\mid S_M=-1]$,
where $f_M$ is the density of $S_M$. When $Q_\chi(M)<1$ the
conditional expectation is positive, which is the infinitesimal
version of Theorem~\ref{thm:B}(ii). The identity is not used in this
section; it is made rigorous, for every $M>0$, in
Lemma~\ref{lem:diff} (see \eqref{eq:Did} and the remark after it).
\end{remark}

\subsection{Sharpness of the condition \texorpdfstring{$Q_\chi\le1$}{Q<=1}
for almost sure nesting under the equal-phase coupling}

\begin{proposition}\label{prop:sharpB}
Assume GRH$(\chi)$ and NCZ$(\chi)$ and let $0<M<M'$. Then
\[
\PP\bigl(X_{M'}>0\ \text{and}\ X_M\le0\bigr)
\;\begin{cases}=0&\text{if }\eta_\chi(M,M')\ge0,\\[2pt]
>0&\text{if }\eta_\chi(M,M')<0.\end{cases}
\]
In particular, a nondegenerate interval $I\subset(0,\infty)$ has the
property that
almost surely $\{X_{M'}>0\}\subset\{X_M>0\}$ for all $M<M'$ in $I$ if
and only if $I\subset(0,M_0(\chi)]$.
\end{proposition}

\begin{proof}
The case $\eta_\chi\ge0$ is Proposition~\ref{prop:nest}. Let
$\eta:=\eta_\chi(M,M')<0$ and write $\Delta_j=\Delta_j(M,M')$,
$c_j'=c_{\gamma_j}(M')$. Choose $J$ with
$\sum_{j>J}\Delta_j\le|\eta|/8$ (Lemma~\ref{lem:Q}(iii)) and then
$\eps'\in(0,1)$ with
$\eps'\sum_{j\le J}\Delta_j\le|\eta|/4$ and
$\eps'\sum_{j\le J}c_j'\le M'|\eta|/4$. Let $A$ be the event that
$\cos\theta_j<-1+\eps'$ for all $j\le J$; then $\PP(A)>0$. On
$A\cap\Omega_0$, by Lemma~\ref{lem:det}(ii),
\[
\frac{X_M}{M}-\frac{X_{M'}}{M'}
\le\Bigl(\frac1M-\frac1{M'}\Bigr)-(1-\eps')\sum_{j\le J}\Delta_j
+\sum_{j>J}\Delta_j
=\eta+\eps'\sum_{j\le J}\Delta_j+2\sum_{j>J}\Delta_j
\le\frac{\eta}{2} ,
\]
and the head $H':=1+\sum_{j\le J}c_j'\cos\theta_j$ satisfies
$h_0\le H'\le h_0+M'|\eta|/4$ with $h_0:=1-\sum_{j\le J}c_j'$. The
tail $W':=\sum_{j>J}c_j'\cos\theta_j$ is independent of $A$ and has
full support by Lemma~\ref{lem:fullsupport} (applied to the tail
sequence), so the event
$B:=\{-h_0<W'<-h_0+M'|\eta|/4\}$ has positive probability and
$\PP(A\cap B)=\PP(A)\PP(B)>0$. On $A\cap B\cap\Omega_0$ we have
$0<X_{M'}=H'+W'<M'|\eta|/2$, and therefore
$X_M/M\le X_{M'}/M'+\eta/2<|\eta|/2-|\eta|/2=0$.

For the last assertion: if $I\subset(0,M_0]$ then $\eta_\chi(M,M')>0$
for all $M<M'$ in $I$, and the nesting holds on $\Omega_0$ for all
pairs at once. Otherwise $I$ contains two points $M_0\le M<M'$; then $Q_\chi>1$ on $(M,M')$, so $\eta_\chi(M,M')<0$
and nesting fails with positive probability for this pair.
\end{proof}

Proposition~\ref{prop:sharpB} is a statement about the
\emph{equal-phase} coupling of Section~\ref{sec:thmB} only (the same
phase vector for every weight): that coupling ceases to nest the
events beyond $M_0$. It does \emph{not} say that $\delta_\chi$ ceases
to be monotone there. Other couplings of the same marginal laws may
nest the events on a larger range; the quantile coupling
$X_M=F_M^{-1}(U)$ nests $\{X_{M'}>0\}\subset\{X_M>0\}$ exactly when
$\delta_\chi(m)\ge\delta_\chi(m')$, so it does so exactly where
$\delta_\chi$ is monotone, but it gives no means of proving
monotonicity. Note also that
$\eta_\chi(M,M')\ge0$ can hold for pairs with $M<M_0<M'$ (the
positive part of the integral over $[M,M_0]$ may dominate), so
Theorem~\ref{thm:B}(i) compares some pairs of weights that are not
both in the range of Theorem~\ref{thm:B}(ii); for instance, for
$\chi_4$ the numerical values of Section~\ref{sec:numerics} give
$\eta_{\chi_4}(1,5)\approx0.55$ and $\eta_{\chi_4}(2,5)\approx0.09$,
whence $\delta_4(\tfrac12)>\delta_4(\tfrac92)$ and
$\delta_4(\tfrac32)>\delta_4(\tfrac92)$.

\subsection{Bounds for \texorpdfstring{$M_0(\chi)$}{M0(chi)}}

\begin{proposition}\label{prop:M0}
Assume GRH$(\chi)$ and NCZ$(\chi)$, and let $\gamma_1$ be the least
positive ordinate of $L(s,\chi)$.
\begin{enumerate}
\item[(a)] $\displaystyle
\Bigl(4\sum_{\gamma>0}\gamma^{-3}\Bigr)^{-1/3}<M_0(\chi)<\kappa_0\gamma_1$
with $\kappa_0=(4^{2/3}-1)^{-1/2}=0.8111\ldots$
\item[(b)] $\displaystyle M_0(\chi)>
\Bigl(\frac{\gamma_1}{2C(\chi)}\Bigr)^{1/3}$, where
$C(\chi)=(\log\Lambda)''(\tfrac12,\chi)$ is given in closed form by
Proposition~\ref{prop:WRvar}.
\item[(c)] $Q_\chi(M)<\sigma_M^2/2$ for all $M>0$. Hence
$M_0(\chi)>M_\sigma(\chi)$, the unique root of $\sigma_M^2=2$.
\item[(d)] $M_0(\chi)\ge\min(1,\gamma_1)\,(c_1\log q)^{-1/3}$ with
$c_1:=12.8\,C_R'$, where $C_R'$ is the absolute constant of
\cite[Lemma~9(b)]{WR} (there $C_R'=2+2C_R$, and $C_R=2.8$ is
admissible by \cite[Corollary~1.2]{BMOR21}, so $c_1\le98$).
\end{enumerate}
\end{proposition}

\begin{proof}
(a) By Lemma~\ref{lem:Q}(i), $Q_\chi(M)<4M^3\sum_\gamma\gamma^{-3}$,
so $Q_\chi<1$ at $M=(4\sum_\gamma\gamma^{-3})^{-1/3}$ and $M_0$ is
larger. For the upper bound, $Q_\chi(M)>g(M/\gamma_1)$ with $g$ as in
the proof of Lemma~\ref{lem:Q} (there are further zeros), and
$g(t)=1$ if and only if $4^{2/3}t^2=1+t^2$, i.e.\ $t=\kappa_0$; hence
$Q_\chi(\kappa_0\gamma_1)>1$.

(b) $\sum_\gamma\gamma^{-3}\le\gamma_1^{-1}\sum_\gamma\gamma^{-2}
=C(\chi)/(2\gamma_1)$; insert into (a).

(c) $\dfrac{4M^3}{(M^2+\gamma^2)^{3/2}}
=\dfrac{4M^2}{M^2+\gamma^2}\cdot\dfrac{M}{\sqrt{M^2+\gamma^2}}
<\dfrac{4M^2}{M^2+\gamma^2}$, and
$\sum_\gamma4M^2/(M^2+\gamma^2)=\sigma_M^2/2$. Each term of
$\sigma_M^2$ is continuous and strictly increasing in $M$, from $0$
to $8$, and the series converges locally uniformly, so $\sigma_M^2$
increases continuously from $0$ to $\infty$ and $M_\sigma$ is well
defined; $Q_\chi(M_\sigma)<1$ gives $M_0>M_\sigma$.

(d) Let $N_\chi(T)$ be the number of ordinates in $(0,T]$. By
\cite[Lemma~9(b)]{WR} (which assumes GRH$(\chi)$ and NCZ$(\chi)$),
$N_\chi(T)\le C_R'\,T\log(qT)$ for $T\ge1$. Partial summation gives
\[
\sum_{\gamma>1}\gamma^{-3}=-N_\chi(1)+3\int_1^\infty N_\chi(t)t^{-4}\,dt
\le3C_R'\int_1^\infty\frac{\log q+\log t}{t^3}\,dt
=C_R'\bigl(\tfrac32\log q+\tfrac34\bigr),
\]
and $\sum_{\gamma\le1}\gamma^{-3}\le N_\chi(1)\gamma_1^{-3}\le
C_R'\gamma_1^{-3}\log q$. Since $\tfrac34\le0.69\log q$ for $q\ge3$,
\[
4\sum_{\gamma>0}\gamma^{-3}\;\le\;
4C_R'\log q\,\bigl(\gamma_1^{-3}+2.19\bigr)
\;\le\;12.8\,C_R'\log q\cdot\min(1,\gamma_1)^{-3},
\]
and (a) gives the claim.
\end{proof}

\begin{remark}[numerical values]\label{rem:M0values}
For $q=4$: $\gamma_1=6.0209\ldots$, $C(\chi_4)=0.1560296\ldots$, and
the bounds of Proposition~\ref{prop:M0} are
\[
2.682<3.298<M_0(\chi_4)<4.884,\qquad M_\sigma(\chi_4)=1.8239\ldots,
\qquad M_0(\chi_4)\in[3.688,\,3.690].
\]
For $q=3$: $\gamma_1=8.0397\ldots$, $C(\chi_3)=0.1134021\ldots$,
\[
3.284<4.103<M_0(\chi_3)<6.522,\qquad M_\sigma(\chi_3)=2.1277\ldots,
\qquad M_0(\chi_3)\in[4.483,\,4.485].
\]
The first number in each chain is (b); the second is (a), evaluated
from the certified zero lists of \cite[\S7]{WR} with the unlisted
zeros included through the bound
$\sum_{\gamma>\Gamma}\gamma^{-3}\le\Gamma^{-1}\sum_{\gamma>\Gamma}\gamma^{-2}$
(the listed zeros alone would give the larger numbers $3.2986$ and
$4.1043$, which are not lower bounds for $M_0$). The brackets for
$M_0$ are obtained from the two-sided bound
\begin{equation}\label{eq:Qenclose}
Q_\Gamma(M)\;\le\;Q_\chi(M)\;\le\;Q_\Gamma(M)
+\frac{2M^3}{\Gamma}\Bigl(C(\chi)-\sum_{\gamma\le\Gamma}\frac{2}{\gamma^2}\Bigr),
\qquad Q_\Gamma(M):=\sum_{\gamma\le\Gamma}\frac{4M^3}{(M^2+\gamma^2)^{3/2}},
\end{equation}
valid for any $\Gamma$ up to which the list of ordinates is complete
(the tail is at most $4M^3\sum_{\gamma>\Gamma}\gamma^{-3}\le
4M^3\Gamma^{-1}\sum_{\gamma>\Gamma}\gamma^{-2}$), with
$\Gamma=639.9\ldots$ ($511$ zeros) for $\chi_4$ and
$\Gamma=699.7\ldots$ ($537$ zeros) for $\chi_3$. These evaluations
are in multiprecision floating point, not in interval arithmetic.
Ball-arithmetic enclosures of $M_0(\chi_4)$ and $M_0(\chi_3)$,
obtained with a sharper tail bound and contained in the brackets
above, are part of the direct certificate described in
Section~\ref{subsec:confirm}(a); see \eqref{eq:M0cert}. See also Section~\ref{sec:numerics}.
\end{remark}

\begin{remark}[the size of $M_0(\chi)$ for large $q$]\label{rem:corner}
In the low-lying normalisation $y_j:=\gamma_j\log q/(2\pi)$ of
\cite[Remark~10]{WR} and with $\lambda:=M\log q$, one has
\emph{exactly}
\[
Q_\chi\Bigl(\frac{\lambda}{\log q}\Bigr)
=\sum_{j\ge1}g\Bigl(\frac{\lambda}{2\pi y_j}\Bigr),
\qquad g(t)=\frac{4t^3}{(1+t^2)^{3/2}} ,
\]
so $\lambda_0(\chi):=M_0(\chi)\log q$ is the root of a function of the
rescaled ordinates alone, and Proposition~\ref{prop:M0}(a) reads
$\lambda_0(\chi)<2\pi\kappa_0\,y_1=\kappa_0\,\gamma_1\log q$. These are
theorems. The range of Theorem~\ref{thm:B} is thus tied to the lowest
rescaled zeros, i.e.\ to the ``critical corner'' of
\cite[Remarks~9 and~10]{WR}.

\emph{A heuristic that should not be trusted.} Replacing the sum by
an integral against the mean density one of the rescaled ordinates
gives $\int_0^\infty g(\lambda/(2\pi y))\,dy=2\lambda/\pi$ and would
suggest $M_0(\chi)\approx\pi/(2\log q)$. This is not supported by
computation, and we do not assert it: the mean density is a poor
description of the lowest zeros of quadratic characters (in the
Katz--Sarnak symplectic model \cite{KS99} the one-level density
vanishes at the origin), and $\lambda_0(\chi)$ depends on the first
few rescaled zeros of the individual character.

\emph{Numerical evidence (quoted, not certified).} The repository
note \texttt{EXTENSIONS.md} (in the directory \texttt{monotonicity/};
see its \S4.2) reports floating-point computations of $M_0(\chi)$ for eleven real primitive characters with
$5\le q\le10007$ (low zeros located numerically, not certified; we
did not re-run them). There $M_0(\chi)\log q$ takes values between
about $3.1$ and $6.7$, shows no sign of converging, and follows
$\gamma_1\log q$: the smallest value, $3.14$, occurs at $q=1013$,
where $\gamma_1\log q=4.11$ is also smallest, and the largest,
$6.74$, at $q=1009$, where $\gamma_1\log q=11.89$ is largest. For
$q=3,4$ the values of Remark~\ref{rem:M0values} give
$M_0\log q=4.93$ and $5.11$. All of these are consistent with the
proved inequality $\lambda_0(\chi)<0.8112\,\gamma_1\log q$.

\emph{Bounds in terms of $q$.} Proposition~\ref{prop:M0}(d) is proved
here. The same note derives sharper asymptotic bounds from an
explicit zero count of Bennett, Martin, O'Bryant and Rechnitzer
\cite{BMOR21} with leading error constant $b=0.22737$: as
$q\to\infty$, uniformly over real primitive $\chi$ mod $q$ satisfying
GRH$(\chi)$ and NCZ$(\chi)$,
\[
\frac{\gamma_1}{\bigl((A_0+A_1\gamma_1^3)\log q\bigr)^{1/3}}\,(1-o(1))
\;\le\;M_0(\chi)\;\le\;
\Bigl(\frac{\pi^3b^2}{\log q}\Bigr)^{1/3}(1+o(1))
=\frac{1.1703}{(\log q)^{1/3}}\,(1+o(1)),
\]
with $A_0=2\bigl(\tfrac5{7\pi}+b\bigr)$ and
$A_1=6\bigl(\tfrac1{2\pi(5/7)^2}+\tfrac{b}{3(5/7)^3}\bigr)$. We quote
these bounds from that note. We re-derived the passage from the zero
count to the bounds (Abel summation,
$Q_\chi(u)=\int_0^\infty N_\chi(t)\,|f_u'(t)|\,dt$ with
$f_u(t)=4u^3(u^2+t^2)^{-3/2}$ and $|f_u'(t)|\le12u^3t^{-4}$), but we
did not check the form of the zero count against \cite{BMOR21}, so
they are not claimed as theorems of this article. Neither bound
reaches the scale $1/\log q$: the upper bound is of order
$(\log q)^{-1/3}$, and for $\gamma_1\asymp1/\log q$ the lower bound is
of order $(\log q)^{-4/3}$. Whether $M_0(\chi)\log q$ is bounded above
and below along all real primitive characters is open; a zero count
with error $O(\log q)$ cannot decide it, because at height
$\asymp1/\log q$ the error is as large as the count.
\end{remark}

\section{An absolutely convergent explicit formula for the weight shift}
\label{sec:formula}

In this section we prove Theorem~\ref{thm:A}(i). Only GRH$(\chi)$ is
assumed; a zero at $s=\tfrac12$ is allowed.

\subsection{Two lemmas}

\begin{lemma}\label{lem:intlog}
Let $a>0$ and $k\ge0$. For all $x\ge2$,
\[
\int_{3/2}^xt^{a-1}(\log t)^{-k}\,dt\;\le\;C(a,k)\,\frac{x^a}{(\log x)^k}.
\]
Moreover, for complex $s$ with $\re s=a$ and all $x\ge2$,
\begin{equation}\label{eq:intlog-s}
\int_{3/2}^x\frac{t^{s-1}}{\log t}\,dt\;=\;\frac{x^s}{s\log x}
+O_a\Bigl(\frac{x^a}{\log^2x}\Bigr),
\end{equation}
with an implied constant depending on $a$ only (not on $\im s$).
\end{lemma}

\begin{proof}
For $x\ge\tfrac94$ split at $\sqrt x\ge\tfrac32$: on
$[\tfrac32,\sqrt x]$ the integrand is at most
$(\log\tfrac32)^{-k}t^{a-1}$, contributing at most
$(\log\tfrac32)^{-k}x^{a/2}/a$, and on $[\sqrt x,x]$ it is at most
$2^k(\log x)^{-k}t^{a-1}$, contributing at most
$2^k(\log x)^{-k}x^a/a$; finally
$x^{a/2}\ll_{a,k}x^a(\log x)^{-k}$. For $2\le x<\tfrac94$ the claim is
trivial. For \eqref{eq:intlog-s} integrate by parts:
\[
\int_{3/2}^x\frac{t^{s-1}}{\log t}\,dt
=\frac{x^s}{s\log x}-\frac{(3/2)^s}{s\log(3/2)}
+\frac1s\int_{3/2}^x\frac{t^{s-1}}{\log^2t}\,dt ,
\]
and use $|s|\ge a$, $|t^{s-1}|=t^{a-1}$ and the first part with $k=2$.
\end{proof}

\begin{lemma}[explicit formula for $\psi_1$]\label{lem:psi1}
Unconditionally, for all $t\ge1$,
\[
\psi_1(t,\chi)\;=\;-\sum_\rho\frac{t^{\rho+1}}{\rho(\rho+1)}+r_1(t),
\qquad |r_1(t)|\le C_\chi\,t\log(2t),
\]
where the sum over all nontrivial zeros converges absolutely:
$\sum_\rho|\rho(\rho+1)|^{-1}\le\sum_\rho|\rho|^{-2}<\infty$.
\end{lemma}

\begin{proof}
Since $0<\beta<1$ we have $|\rho+1|\ge|\rho|>0$, so
$|\rho(\rho+1)|^{-1}\le|\rho|^{-2}$, which is summable by
\eqref{eq:zerosum}; also $|t^{\rho+1}|\le t^2$ for $t\ge1$. Let
$t\ge2$. Since $\psi(u,\chi)=0$ for $u<2$,
$\psi_1(t,\chi)=\int_2^t\psi(u,\chi)\,du$. Fix $T\ge2$ and integrate
the formula of Lemma~\ref{lem:WRclassical}, whose error term is
uniform in $u\in[2,t]$, over $u\in[2,t]$:
\[
\psi_1(t,\chi)=-\sum_{|\gamma|\le T}
\frac{t^{\rho+1}-2^{\rho+1}}{\rho(\rho+1)}
+O\Bigl(\frac{t^2\log^2(tT)}{T}+t\log t\Bigr),
\]
with an implied constant independent of $t$ and $T$. Let
$T\to\infty$ with $t$ fixed: the first error term tends to $0$ and
the sum converges absolutely to the sum over all zeros, so
\[
\psi_1(t,\chi)=-\sum_\rho\frac{t^{\rho+1}}{\rho(\rho+1)}
+\sum_\rho\frac{2^{\rho+1}}{\rho(\rho+1)}+O(t\log t),
\]
and the middle term is bounded by $4\sum_\rho|\rho|^{-2}$. For
$1\le t<2$, $\psi_1(t,\chi)=0$ and
$|\sum_\rho t^{\rho+1}/(\rho(\rho+1))|\le4\sum_\rho|\rho|^{-2}$.
\end{proof}

\subsection{The formula}

\begin{proposition}\label{prop:formula}
Assume GRH$(\chi)$. Let $-\tfrac12<m<m'$, $\kappa=m'-m$,
$M=m+\tfrac12$, $M'=m'+\tfrac12$, and for $x\ge2$ put
\begin{equation}\label{eq:defI}
I(x)\;=\;I_{m,m'}(x)\;:=\;D_m(x)-x^{-\kappa}D_{m'}(x) .
\end{equation}
\begin{enumerate}
\item[(a)] (Exact identities.) For all $x\ge2$,
\[
I(x)=\sum_{p\le x}\chi^*(p)\,p^m\Bigl(1-\Bigl(\frac px\Bigr)^{\kappa}\Bigr)
=\kappa\int_2^xt^{-\kappa-1}D_{m'}(t)\,dt .
\]
\item[(b)] For all $x\ge2$,
\begin{equation}\label{eq:formula}
I(x)=\frac{\kappa\,x^{M}}{\log x}\Bigl[\frac{1+2n_0}{2MM'}
+2\,\re\sum_{\gamma>0}\frac{x^{i\gamma}}{(M+i\gamma)(M'+i\gamma)}\Bigr]
+O\Bigl(\frac{x^{M}}{\log^2x}\Bigr),
\end{equation}
where $\sum_{\gamma>0}|(M+i\gamma)(M'+i\gamma)|^{-1}
\le\sum_{\gamma>0}\gamma^{-2}<\infty$ and the implied constant depends
only on $\chi$, $m$, $m'$.
\end{enumerate}
\end{proposition}

\begin{proof}
\emph{(a).} The first equality is the definition. For the second,
$D_m(x)=\int_{3/2}^xt^{-\kappa}\,dD_{m'}(t)
=x^{-\kappa}D_{m'}(x)+\kappa\int_{3/2}^xt^{-\kappa-1}D_{m'}(t)\,dt$ by
Riemann--Stieltjes partial summation ($D_{m'}(\tfrac32)=0$), and
$D_{m'}=0$ on $[\tfrac32,2)$.

\emph{(b).} All quantities in \eqref{eq:formula} are bounded for $x$
in a bounded subinterval of $[2,\infty)$, so we may assume $x\ge4$.
Put
\[
g(t):=t^m-x^{-\kappa}t^{m'}=t^m\bigl(1-(t/x)^\kappa\bigr),
\qquad
h(t):=\frac{g(t)}{\log t}\qquad(\tfrac32\le t\le x).
\]
Thus $I(x)=\sum_{p\le x}\chi^*(p)g(p)=-\sum_{p\le x}\chi(p)(\log p)h(p)$,
and $g(x)=h(x)=0$. For $\tfrac32\le t\le x$ we have $(t/x)^\kappa\le1$,
hence
\begin{equation}\label{eq:gbounds}
0\le g(t)\le t^m,\qquad |g'(t)|\le(|m|+|m'|)\,t^{m-1},\qquad
g'(x)=-\kappa x^{m-1},
\end{equation}
and $g''(t)=m(m-1)t^{m-2}-m'(m'-1)x^{-\kappa}t^{m'-2}$, so
$|g''(t)|\ll t^{m-2}$. Writing $L=\log t$,
\begin{equation}\label{eq:hderivs}
h'=\frac{g'}{L}-\frac{g}{tL^2},\qquad
h''=\frac{g''}{L}+k,\qquad
k:=-\frac{2g'}{tL^2}+\frac{g}{t^2L^2}+\frac{2g}{t^2L^3},
\end{equation}
so that, by \eqref{eq:gbounds} and $L\ge\log\tfrac32$,
\begin{equation}\label{eq:hbounds}
|k(t)|\ll\frac{t^{m-2}}{\log^2t},\qquad
|h''(t)|\ll\frac{t^{m-2}}{\log t},\qquad
h'(x)=-\frac{\kappa\,x^{m-1}}{\log x}.
\end{equation}

\emph{Step 1 (prime powers).} Let
$S_\Lambda(x):=\sum_{n\le x}\Lambda(n)\chi(n)h(n)$. Sorting by prime
powers $n=p^k$ and using $h(p^k)\log p=g(p^k)/k$,
\begin{equation}\label{eq:SLsplit}
S_\Lambda(x)=-I(x)+\frac12\sum_{\substack{p\le\sqrt x\\ p\nmid q}}g(p^2)
+\sum_{k\ge3}\frac1k\sum_{p^k\le x}\chi(p)^kg(p^k).
\end{equation}
By \eqref{eq:gbounds}, the last double sum is at most
$\sum_{k\ge3}\sum_{p^k\le x}p^{km}
\le(\log2)^{-1}\sum_{k\ge3}\sum_{p^k\le x}p^{km}\log p$, and the
latter is $O(x^{m+1/3}\log x+\log^2x)$ by \cite[proof of Lemma~3,
``the $k\ge3$ terms'']{WR} (Chebyshev's bound $\theta(y)\ll y$ and
Mertens' theorem; unconditional). Since $M>0$ and $M>m+\tfrac13$, this
is $O(x^M/\log^2x)$. The primes $p\mid q$ contribute
$O(1)$ to the middle sum, as $0\le g(p^2)\le\max(1,q^{2m})$.

\emph{Step 2 (the bias).} Let $y=\sqrt x\ge2$ and
$f(t):=g(t^2)=t^{2m}\bigl(1-(t/y)^{2\kappa}\bigr)$ on $[\tfrac32,y]$,
so $f(y)=0$ and $|f'(t)|\ll t^{2m-1}$. With
$\ell(t):=\int_{3/2}^t du/\log u$, the prime number theorem gives
$\pi(t)=\ell(t)+O(t/\log^2t)$ for $t\ge\tfrac32$
\cite[\S18]{Dav00}. Partial summation, using $f(y)=0$ and
$\pi(\tfrac32)=\ell(\tfrac32)=0$, gives
\[
\sum_{p\le y}f(p)=-\int_{3/2}^yf'(t)\pi(t)\,dt
=\int_{3/2}^y\frac{f(t)}{\log t}\,dt
+O\Bigl(\int_{3/2}^y\frac{t^{2m}}{\log^2t}\,dt\Bigr).
\]
By Lemma~\ref{lem:intlog} with $a=2M>0$ the error is
$O(y^{2M}/\log^2y)$, and by \eqref{eq:intlog-s} with the real
exponents $s=2M$ and $s=2M'$,
\[
\int_{3/2}^y\frac{f(t)}{\log t}\,dt
=\frac{y^{2M}}{2M\log y}-y^{-2\kappa}\frac{y^{2M'}}{2M'\log y}
+O\Bigl(\frac{y^{2M}}{\log^2y}\Bigr)
=\frac{y^{2M}}{\log y}\cdot\frac{\kappa}{2MM'}
+O\Bigl(\frac{y^{2M}}{\log^2y}\Bigr).
\]
Since $y^{2M}=x^M$ and $\log y=\tfrac12\log x$,
\begin{equation}\label{eq:biasterm}
\frac12\sum_{\substack{p\le\sqrt x\\ p\nmid q}}g(p^2)
=\frac{\kappa\,x^M}{2MM'\log x}+O\Bigl(\frac{x^M}{\log^2x}\Bigr).
\end{equation}

\emph{Step 3 (two integrations by parts).} The function
$\psi(\cdot,\chi)$ is a step function vanishing on $[0,2)$, and
$h\in C^2[\tfrac32,x]$ with $h(x)=0$. Hence
\[
S_\Lambda(x)=\int_{3/2}^xh\,d\psi=-\int_{3/2}^xh'(t)\psi(t,\chi)\,dt
=-h'(x)\psi_1(x,\chi)+\int_{3/2}^xh''(t)\psi_1(t,\chi)\,dt ,
\]
where we used $\psi(\tfrac32,\chi)=\psi_1(\tfrac32,\chi)=0$ and that
$\psi_1(\cdot,\chi)$ is absolutely continuous with derivative
$\psi(\cdot,\chi)$ almost everywhere. By \eqref{eq:hbounds},
\begin{equation}\label{eq:SLpsi1}
S_\Lambda(x)=\frac{\kappa\,x^{m-1}}{\log x}\,\psi_1(x,\chi)
+\int_{3/2}^xh''(t)\,\psi_1(t,\chi)\,dt .
\end{equation}

\emph{Step 4 (inserting the zeros).} Insert Lemma~\ref{lem:psi1}.
Under GRH$(\chi)$, $|t^{\rho+1}|=t^{3/2}$, so the series
$\sum_\rho t^{\rho+1}/(\rho(\rho+1))$ converges absolutely and
uniformly on $[\tfrac32,x]$ and may be integrated termwise against
the bounded function $h''$. Thus
\begin{equation}\label{eq:SLzeros}
S_\Lambda(x)=\sum_\rho c_\rho(x)+R(x),
\end{equation}
\[
c_\rho(x):=-\frac{1}{\rho(\rho+1)}\Bigl[\frac{\kappa\,x^{m+\rho}}{\log x}
+\int_{3/2}^xh''(t)\,t^{\rho+1}\,dt\Bigr],
\qquad
R(x):=\frac{\kappa\,x^{m-1}}{\log x}\,r_1(x)+\int_{3/2}^xh''(t)\,r_1(t)\,dt .
\]
By Lemma~\ref{lem:psi1} and \eqref{eq:hbounds}, and since
$\log(2t)/\log t$ is bounded on $[\tfrac32,\infty)$,
\[
|R(x)|\ll x^m+\int_{3/2}^xt^{m-1}\,dt\ll(1+x^m)\log x
\ll\frac{x^M}{\log^2x},
\]
because $M>0$ and $M>m$.

\emph{Step 5 (one zero).} Fix a zero $\rho=\tfrac12+i\gamma$,
$\gamma\in\mathbb R$, and write $s=m+\rho=M+i\gamma$,
$s'=m'+\rho=M'+i\gamma$. By \eqref{eq:hderivs},
\[
\int_{3/2}^xh''(t)\,t^{\rho+1}\,dt
=m(m-1)\int_{3/2}^x\frac{t^{s-1}}{\log t}\,dt
-m'(m'-1)\,x^{-\kappa}\int_{3/2}^x\frac{t^{s'-1}}{\log t}\,dt
+\int_{3/2}^xk(t)\,t^{\rho+1}\,dt .
\]
The last integral is $O(\int_{3/2}^xt^{M-1}\log^{-2}t\,dt)
=O(x^M/\log^2x)$ by \eqref{eq:hbounds} and Lemma~\ref{lem:intlog}.
Applying \eqref{eq:intlog-s} with $a=M$ and with $a=M'$ (note
$x^{-\kappa}x^{M'}=x^M$ and $x^{-\kappa}x^{s'}=x^{s}$),
\[
\int_{3/2}^xh''(t)\,t^{\rho+1}\,dt
=\Bigl[\frac{m(m-1)}{m+\rho}-\frac{m'(m'-1)}{m'+\rho}\Bigr]
\frac{x^{m+\rho}}{\log x}+O\Bigl(\frac{x^M}{\log^2x}\Bigr),
\]
with an implied constant independent of $\gamma$. Now
$m(m-1)=(m+\rho)(m-1-\rho)+\rho(\rho+1)$, and likewise for $m'$, so
\begin{equation}\label{eq:bracket}
\kappa+\frac{m(m-1)}{m+\rho}-\frac{m'(m'-1)}{m'+\rho}
=\rho(\rho+1)\Bigl[\frac1{m+\rho}-\frac1{m'+\rho}\Bigr]
=\frac{\kappa\,\rho(\rho+1)}{(m+\rho)(m'+\rho)} .
\end{equation}
Therefore
\begin{equation}\label{eq:crho}
c_\rho(x)=-\frac{\kappa\,x^{M+i\gamma}}{(M+i\gamma)(M'+i\gamma)\log x}
+O\Bigl(\frac{1}{|\rho|^2}\cdot\frac{x^M}{\log^2x}\Bigr),
\end{equation}
uniformly in $\rho$.

\emph{Step 6 (assembling).} Summing \eqref{eq:crho} over all zeros
(the error terms are summable by \eqref{eq:zerosum}) and inserting
into \eqref{eq:SLzeros},
\[
S_\Lambda(x)=-\frac{\kappa\,x^M}{\log x}
\sum_\rho\frac{x^{i\gamma}}{(M+i\gamma)(M'+i\gamma)}
+O\Bigl(\frac{x^M}{\log^2x}\Bigr).
\]
The zeros are $\tfrac12$ (with multiplicity $n_0$) and the conjugate
pairs $\tfrac12\pm i\gamma$, $\gamma>0$; hence
\[
\sum_\rho\frac{x^{i\gamma}}{(M+i\gamma)(M'+i\gamma)}
=\frac{n_0}{MM'}+2\,\re\sum_{\gamma>0}\frac{x^{i\gamma}}{(M+i\gamma)(M'+i\gamma)}.
\]
Combining with \eqref{eq:SLsplit}, Step~1 and \eqref{eq:biasterm},
\[
I(x)=\frac{\kappa\,x^M}{2MM'\log x}-S_\Lambda(x)
+O\Bigl(\frac{x^M}{\log^2x}\Bigr),
\]
which is \eqref{eq:formula}.
\end{proof}

\begin{remark}[signs and conventions]\label{rem:signs}
With $\chi^*=-\chi$ the bias term of $I$ is \emph{positive} (it comes
from the squares of primes, $\chi(p)^2=1$), a central zero increases
it, and the oscillating sum enters with a plus sign and the kernel
$1/((M+i\gamma)(M'+i\gamma))$. In terms of the normalised races
\eqref{eq:normalised}, \eqref{eq:formula} reads
\begin{equation}\label{eq:formula-normalised}
\frac{\cE_{m,0}(x)}{M}-\frac{\cE_{m',0}(x)}{M'}
=\kappa\Bigl[\frac{1+2n_0}{MM'}
+4\,\re\sum_{\gamma>0}\frac{x^{i\gamma}}{(M+i\gamma)(M'+i\gamma)}\Bigr]
+O\Bigl(\frac1{\log x}\Bigr),
\end{equation}
which is what one obtains formally by subtracting the (conditionally
convergent) expansions of Lemma~\ref{lem:WRexp} for the two weights,
since $r_\gamma(M)/M-r_\gamma(M')/M'=4\kappa/((M+i\gamma)(M'+i\gamma))$.
The point of Proposition~\ref{prop:formula} is that the difference is
controlled \emph{pointwise}, with an absolutely convergent zero sum
and an error $O(1/\log x)$, whereas each of the two expansions is
controlled only in mean square. Formula \eqref{eq:formula} was checked
numerically against a sieve (Section~\ref{sec:numerics}).
\end{remark}

\section{Inheritance of leads: Theorem~\ref{thm:A} and its corollaries}
\label{sec:thmA}

\subsection{Proof of Theorem~\ref{thm:A}}

\begin{proof}[Proof of Theorem~\ref{thm:A}]
(i) is Proposition~\ref{prop:formula}.

(ii) Put $B:=\sum_{\gamma>0}4/(|M+i\gamma|\,|M'+i\gamma|)$, so that
$\lambda_\chi(M,M')=(1+2n_0)/(MM')-B$. For every $x$,
\begin{equation}\label{eq:bracketbounds}
\begin{gathered}
\Bigl|2\,\re\sum_{\gamma>0}\frac{x^{i\gamma}}{(M+i\gamma)(M'+i\gamma)}
\Bigr|\le\frac B2,
\qquad\text{so}\\
\frac{\lambda_\chi(M,M')}{2}\;\le\;
\frac{1+2n_0}{2MM'}
+2\,\re\sum_{\gamma>0}\frac{x^{i\gamma}}{(M+i\gamma)(M'+i\gamma)}
\;\le\;\frac{1+2n_0}{MM'}-\frac{\lambda_\chi(M,M')}{2}.
\end{gathered}
\end{equation}
Let $C_0$ be the implied constant in \eqref{eq:formula}. Then
$I(x)\ge\frac{\kappa x^M}{\log x}\bigl(\frac{\lambda_\chi}{2}
-\frac{C_0}{\kappa\log x}\bigr)\ge\frac{\kappa\lambda_\chi}{4}\cdot
\frac{x^M}{\log x}$ as soon as $\log x\ge4C_0/(\kappa\lambda_\chi)$.
If $x\ge x_0$ and $D_{m'}(x)\ge0$, then
$D_m(x)=x^{-\kappa}D_{m'}(x)+I(x)\ge I(x)>0$; the last assertion is
the contrapositive, with $D_{m'}(x)\le-x^{\kappa}I(x)<0$ when
$D_m(x)\le0$.

(iii) By the Cauchy--Schwarz inequality,
\[
B\le\Bigl(\sum_{\gamma>0}\frac{4}{M^2+\gamma^2}\Bigr)^{1/2}
\Bigl(\sum_{\gamma>0}\frac{4}{M'^2+\gamma^2}\Bigr)^{1/2}
=\frac{\sigma_M}{\sqrt2\,M}\cdot\frac{\sigma_{M'}}{\sqrt2\,M'}
=\frac{\sigma_M\sigma_{M'}}{2MM'} ,
\]
which gives the inequality. If $\sigma_{M'}^2\le2$ then, since
$\sigma_M^2$ is strictly increasing in $M$
(Proposition~\ref{prop:M0}(c)), $\sigma_M\sigma_{M'}<\sigma_{M'}^2\le2
\le2+4n_0$.
\end{proof}

\begin{remark}[range of validity; examples]\label{rem:Arange}
(1) Let $M_\sigma(\chi)$ be the root of $\sigma_M^2=2$. By
Theorem~\ref{thm:A}(iii), $\lambda_\chi(M,M')>0$ for \emph{all}
$0<M<M'\le M_\sigma(\chi)$. Under NCZ$(\chi)$ the variance is given in
closed form by Proposition~\ref{prop:WRvar}; numerically
$M_\sigma(\chi_4)=1.8239\ldots$ and $M_\sigma(\chi_3)=2.1277\ldots$,
i.e.\ $m'\le1.32$ resp.\ $m'\le1.62$.

(2) The Cauchy--Schwarz criterion is far from necessary when $M$ is
small. Evaluating $\lambda_\chi$ from the zero lists of \cite{WR}, with
the tail $\sum_{\gamma>\Gamma}4\gamma^{-2}=2C(\chi)-\sum_{\gamma\le\Gamma}
4\gamma^{-2}$ as a rigorous upper bound for the omitted part
(floating-point evaluation; Section~\ref{sec:numerics}):
\[
\begin{array}{llll}
q=4:\quad&\lambda(\tfrac12,\tfrac32)\in[1.025,1.033],
&\lambda(\tfrac12,\tfrac52)\in[0.498,0.506],
&\lambda(\tfrac12,\tfrac{15}2)\in[0.011,0.018],\\[2pt]
&\lambda(\tfrac32,\tfrac52)\in[-0.032,-0.024];&&\\[2pt]
q=3:\quad&\lambda(\tfrac12,\tfrac32)\in[1.108,1.115],
&\lambda(\tfrac32,\tfrac52)\in[0.045,0.052].&
\end{array}
\]
Thus, under GRH$(\chi_4)$, a lead of the class $3$ mod $4$ in the race
weighted by $p^{7}$ at a large $x$ forces a lead in Chebyshev's
classical race at the same $x$; the pair $(m,m')=(1,2)$ has positive
margin for $q=3$ and negative margin for $q=4$.

(3) Comparison with Theorem~\ref{thm:B}. Under NCZ$(\chi)$, for all
$0<M<M'$,
\begin{equation}\label{eq:etalambda}
\eta_\chi(M,M')\;\ge\;\kappa\,\lambda_\chi(M,M') ,
\end{equation}
because $0<|M+i\gamma|^{-1}-|M'+i\gamma|^{-1}
\le|(M+i\gamma)^{-1}-(M'+i\gamma)^{-1}|
=\kappa/(|M+i\gamma|\,|M'+i\gamma|)$. So every pair covered by
Theorem~\ref{thm:A}(ii) is covered by Theorem~\ref{thm:B}(i) (under
its stronger hypotheses), and the range $\sigma_M^2\le2$ is strictly
inside the range $Q_\chi\le1$ (Proposition~\ref{prop:M0}(c)). The
loss has a transparent origin. In the actual race the phase of the
zero $\gamma$ at weight $M$ is $\gamma\log x-\arg(M+i\gamma)$, which
rotates with $M$; the joint behaviour of the two races at the same
$x$ is described by the ``natural'' coupling of
Lemma~\ref{lem:joint} below, with coefficients
$r_\gamma(M)=4M/(M+i\gamma)$. Theorem~\ref{thm:B} is free to use the
more favourable equal-phase coupling (the same phases for both
weights, coefficients $|r_\gamma(M)|$) because it concerns the
marginal laws only. The pair $(m,m')=(1,2)$ for $\chi_4$ shows the
difference concretely: $\lambda_{\chi_4}(\tfrac32,\tfrac52)<0$, so
under LI the set of $x$ with $D_2(x)>0\ge D_1(x)$ has positive
logarithmic density (Corollary~\ref{cor:sharp}(b)), although
$M'=\tfrac52<M_0(\chi_4)$ and the equal-phase coupling nests
$\{X_{5/2}>0\}\subset\{X_{3/2}>0\}$ almost surely. Theorem~\ref{thm:B}
orders the \emph{densities} of the two lead sets; it does not nest
the lead sets.
\end{remark}

\subsection{Densities without LI}

\begin{corollary}\label{cor:densities}
Assume GRH$(\chi)$ and let $-\tfrac12<m<m'$ with
$\lambda_\chi(M,M')>0$. Let $d$ denote any one of: upper logarithmic
density, lower logarithmic density, upper natural density, lower
natural density. Then
\[
d\{x:D_m(x)>0\}\;\ge\;d\{x:D_{m'}(x)\ge0\}\;\ge\;d\{x:D_{m'}(x)>0\},
\]
and likewise $d\{x:D_{m'}(x)<0\}\ge d\{x:D_m(x)\le0\}\ge d\{x:D_m(x)<0\}$.
In particular, on the interval $(-\tfrac12,M_\sigma(\chi)-\tfrac12]$
the upper and the lower logarithmic (and natural) densities of
$\{x:D_m(x)>0\}$ are non-increasing functions of $m$, under
GRH$(\chi)$ alone.
\end{corollary}

\begin{proof}
By Theorem~\ref{thm:A}(ii), $\{x\ge x_0:D_{m'}(x)\ge0\}
\subset\{x:D_m(x)>0\}$ and $\{x\ge x_0:D_m(x)\le0\}\subset
\{x:D_{m'}(x)<0\}$. Each of the four densities is monotone under
inclusion and insensitive to a bounded set. The last sentence follows
from Remark~\ref{rem:Arange}(1).
\end{proof}

We make no claim about the \emph{existence} of the logarithmic
density without LI, nor about strictness; under LI,
Theorem~\ref{thm:B} gives existence and strict decrease on a larger
range.

\subsection{Sharpness under LI}

\begin{lemma}[joint limiting distribution]\label{lem:joint}
Assume GRH$(\chi)$, NCZ$(\chi)$ and LI$(\chi)$, and let
$-\tfrac12<m<m'$. On the probability space of \eqref{eq:XM} define
\[
Y_M:=1+\re\sum_{\gamma>0}r_\gamma(M)\,e^{i\theta_\gamma},
\qquad
Y_{M'}:=1+\re\sum_{\gamma>0}r_\gamma(M')\,e^{i\theta_\gamma},
\qquad r_\gamma(M)=\frac{4M}{M+i\gamma},
\]
with the \emph{same} phases in both (the series converge almost
surely). Then the $\mathbb R^2$-valued function
$u\mapsto(\cE_{m,0}(e^u),\cE_{m',0}(e^u))$ has the limiting
distribution $\mathrm{law}(Y_M,Y_{M'})$.
The marginals are the laws of $X_M$ and $X_{M'}$, and almost surely
\begin{equation}\label{eq:Ydiff}
\frac{Y_M}{M}-\frac{Y_{M'}}{M'}
=\kappa\Bigl[\frac1{MM'}+\sum_{\gamma>0}b_\gamma\cos(\theta_\gamma-\alpha_\gamma)\Bigr],
\qquad b_\gamma:=\frac{4}{|M+i\gamma|\,|M'+i\gamma|},
\end{equation}
with $\alpha_\gamma:=\arg\bigl((M+i\gamma)(M'+i\gamma)\bigr)$, an
absolutely convergent series.
\end{lemma}

\begin{proof}
The series converge almost surely by Kolmogorov's criterion applied
to real and imaginary parts. Since
$\re(r\,e^{i\theta})=|r|\cos(\theta+\arg r)$ and the shifted phases
are again independent and uniform, $Y_M$ has the law of $X_M$.
Formula \eqref{eq:Ydiff} follows from
$r_\gamma(M)/M-r_\gamma(M')/M'=4\kappa/((M+i\gamma)(M'+i\gamma))$.

Fix $(\xi_1,\xi_2)\in\mathbb R^2$ and consider the real function
$\phi_{\xi}(u):=\xi_1\,\cE_{m,0}(e^u)+\xi_2\,\cE_{m',0}(e^u)$. By
Lemma~\ref{lem:WRexp} (the unlogged case $b=0$) for both weights, for
$u\ge2$ and $T\ge2$,
\[
\phi_{\xi}(u)=(\xi_1+\xi_2)+\re\sum_{0<\gamma\le T}
\bigl(\xi_1r_\gamma(M)+\xi_2r_\gamma(M')\bigr)e^{i\gamma u}
+\cE_{\xi}(u,T),
\]
where $\cE_{\xi}$ is the corresponding combination of the two error
terms, so that
\[
\lim_{Y\to\infty}\frac1Y\int_2^Y|\cE_{\xi}(u,e^Y)|^2\,du=0
\]
(use $(x+y)^2\le2x^2+2y^2$). The coefficients are $O(1/\gamma)$, and
$N(T+1)-N(T)\ll\log(T+2)$. These are the hypotheses of
\cite[Corollary~1.3(a)]{ANS14}; they are verified exactly as in the
proof of \cite[Theorem~1(i)]{WR}. Hence $\phi_{\xi}$ has a limiting
distribution. Under LI$(\chi)$, \cite[Theorem~1.9]{ANS14}
identifies it, exactly as in \cite{WR}, as the law of
\[
(\xi_1+\xi_2)+\sum_{\gamma>0}\bigl|\xi_1r_\gamma(M)+\xi_2r_\gamma(M')\bigr|
\cos\theta_\gamma ,
\]
which is the law of $\xi_1Y_M+\xi_2Y_{M'}$ (shift the phases as
above).

Let $\mu_Y$ be the empirical measures on $\mathbb R^2$ of
$u\mapsto(\cE_{m,0}(e^u),\cE_{m',0}(e^u))$, $u\in[0,Y]$. Applying the
previous paragraph to the bounded continuous functions $\cos$ and
$\sin$ shows that the characteristic function of $\mu_Y$ at
$(\xi_1,\xi_2)$, namely $\frac1Y\int_0^Ye^{i\phi_{\xi}(u)}\,du$,
converges to $\EE\,e^{i(\xi_1Y_M+\xi_2Y_{M'})}$ for every
$(\xi_1,\xi_2)$. The limit is the characteristic function of
$\mathrm{law}(Y_M,Y_{M'})$, so by L\'evy's continuity theorem on
$\mathbb R^2$ \cite[Ch.~XV]{Fel71}, applied along an arbitrary
sequence $Y_n\to\infty$, $\mu_Y$ converges weakly to
$\mathrm{law}(Y_M,Y_{M'})$.
\end{proof}

(Alternatively, Theorems~1.4 and~1.9 of \cite{ANS14} are stated for
functions with values in $\mathbb R^\ell$ and give
Lemma~\ref{lem:joint} directly. We keep the Cram\'er--Wold argument
because it uses only the scalar statements quoted in \cite{WR}.)

\begin{corollary}[sharpness of $\lambda_\chi\ge0$]\label{cor:sharp}
Let $-\tfrac12<m<m'$.
\begin{enumerate}
\item[(a)] Under GRH$(\chi)$,
\[
\liminf_{x\to\infty}\frac{\log x}{\kappa\,x^M}\,I_{m,m'}(x)\;\ge\;
\frac{\lambda_\chi(M,M')}{2},
\]
with equality under LI$(\chi)$. In particular, under GRH$(\chi)$ and
LI$(\chi)$: if $\lambda_\chi(M,M')<0$, then the Riesz-type mean
$\int_2^xt^{-\kappa-1}D_{m'}(t)\,dt$ is negative for arbitrarily
large $x$, whereas it is positive for all large $x$ if
$\lambda_\chi(M,M')>0$.
\item[(b)] Assume GRH$(\chi)$, NCZ$(\chi)$ and LI$(\chi)$. The set of
violations of the inheritance of leads,
\[
\mathcal V\;:=\;\{x\ge2:\ D_{m'}(x)>0\ \text{and}\ D_m(x)\le0\},
\]
has a logarithmic density, equal to $\PP(Y_{M'}>0\ge Y_M)$. This
density is $0$ if $\lambda_\chi(M,M')\ge0$ (and $\mathcal V$ is
bounded if $\lambda_\chi(M,M')>0$), and it is strictly positive if
$\lambda_\chi(M,M')<0$.
\end{enumerate}
\end{corollary}

\begin{proof}
(a) The inequality is \eqref{eq:formula} and
\eqref{eq:bracketbounds}. Under LI$(\chi)$ the positive ordinates are
distinct and linearly independent over $\mathbb Q$. Given $\eps>0$
choose $J$ with $\sum_{j>J}b_{\gamma_j}<\eps$. By Kronecker's theorem
\cite[Ch.~XXIII]{HW} there are arbitrarily large $u$ with
$|\gamma_ju-\alpha_{\gamma_j}-\pi|<\eps$ modulo $2\pi$ for all
$j\le J$. For such $u$ and $x=e^u$,
\[
4\,\re\sum_{\gamma>0}\frac{x^{i\gamma}}{(M+i\gamma)(M'+i\gamma)}
=\sum_{\gamma>0}b_\gamma\cos(\gamma u-\alpha_\gamma)
\le-\cos\eps\sum_{j\le J}b_{\gamma_j}+\eps
\le-B+\eps\,(B+2),
\]
with $B=\sum_{\gamma>0}b_\gamma$ as in the proof of
Theorem~\ref{thm:A} (we used $1-\cos\eps\le\eps$). By
\eqref{eq:formula} the $\liminf$ is at most
$\frac12\bigl(\frac{1+2n_0}{MM'}-B\bigr)+\eps(B+2)/2$ for every
$\eps>0$.

(b) By Lemma~\ref{lem:joint} and Lemma~\ref{lem:stability}(b) applied
to $B_0:=\{(y,y'):y\le0<y'\}$, whose boundary lies in
$\{y=0\}\cup\{y'=0\}$ and is a null set for $\mathrm{law}(Y_M,Y_{M'})$
because both marginals are absolutely continuous
(Theorem~\ref{thm:WRinterp}), the logarithmic density of $\mathcal V$
exists and equals $\PP(Y_{M'}>0\ge Y_M)$. (Here
$\cE_{m,0}(x)$ has the sign of $D_m(x)$.)

On the event $\{Y_{M'}>0\ge Y_M\}$ we have $Y_M/M-Y_{M'}/M'<0$,
while by \eqref{eq:Ydiff} this difference is almost surely at least
$\kappa(1/(MM')-B)=\kappa\lambda_\chi(M,M')$. Hence the event is null
if $\lambda_\chi\ge0$; boundedness of $\mathcal V$ for
$\lambda_\chi>0$ is Theorem~\ref{thm:A}(ii).

Let now $\lambda:=\lambda_\chi(M,M')<0$. Write
$b_j=b_{\gamma_j}$, $\alpha_j=\alpha_{\gamma_j}$,
$c_j'=c_{\gamma_j}(M')$, $\beta_j:=\arg(M'+i\gamma_j)$, so that
$Y_{M'}=1+\sum_jc_j'\cos(\theta_j-\beta_j)$. Choose $J$ with
$\sum_{j>J}b_j\le|\lambda|/8$, and let $A_\tau$ be the event that
$|\theta_j-\alpha_j-\pi|<\tau$ modulo $2\pi$ for all $j\le J$. On
$A_\tau$, by \eqref{eq:Ydiff},
\[
\frac{Y_M}{M}-\frac{Y_{M'}}{M'}
\le\kappa\Bigl[\frac1{MM'}-\cos\tau\sum_{j\le J}b_j+\sum_{j>J}b_j\Bigr]
\le\kappa\Bigl[\lambda+(1-\cos\tau)B+\frac{|\lambda|}{4}\Bigr],
\]
and the head $H':=1+\sum_{j\le J}c_j'\cos(\theta_j-\beta_j)$ differs
from $h_0:=1-\sum_{j\le J}c_j'\cos(\alpha_j-\beta_j)$ by at most
$\tau\sum_{j\le J}c_j'$. Fix $\tau>0$ so small that
$(1-\cos\tau)B\le|\lambda|/4$ and
$\eps_1:=\tau\sum_{j\le J}c_j'\le\kappa|\lambda|M'/8$. Then
$\PP(A_\tau)>0$, and on $A_\tau$ we have
$Y_M/M-Y_{M'}/M'\le\kappa\lambda/2$ and $|H'-h_0|\le\eps_1$. The tail
$W':=\sum_{j>J}c_j'\cos(\theta_j-\beta_j)$ is independent of $A_\tau$
and has full support (Lemma~\ref{lem:fullsupport}), so
$B_1:=\{-h_0+\eps_1<W'<-h_0+2\eps_1\}$ has positive probability. On
$A_\tau\cap B_1$: \ $0<Y_{M'}=H'+W'<3\eps_1$, hence
\[
\frac{Y_M}{M}\le\frac{Y_{M'}}{M'}+\frac{\kappa\lambda}{2}
<\frac{3\eps_1}{M'}-\frac{\kappa|\lambda|}{2}<0 .
\]
Thus $\PP(Y_{M'}>0\ge Y_M)\ge\PP(A_\tau)\PP(B_1)>0$.
\end{proof}

\subsection{A Landau-type converse and an equivalence with GRH}

The next two results are unconditional; no hypothesis on zeros is in
force. Put $P(w,\chi):=\sum_p\chi(p)p^{-w}$ for $\re w>1$.

\begin{lemma}[Landau; {\cite[Lemma~15.1]{MV07}}]\label{lem:landau}
Let $A:[1,\infty)\to\mathbb R$ be bounded and Riemann integrable on
every finite interval, with $A(x)\ge0$ for all $x>X_0$. Let
$\sigma_c$ be the infimum of the real $\sigma$ for which
$\int_{X_0}^\infty A(x)x^{-\sigma-1}\,dx<\infty$. Then
$F(s):=\int_1^\infty A(x)x^{-s-1}\,dx$ is analytic in the half-plane
$\re s>\sigma_c$, and, if $\sigma_c>-\infty$, $F$ has no analytic
continuation to a neighbourhood of the point $s=\sigma_c$.
\end{lemma}

(In \cite{MV07} the kernel is $x^{-s}$; the statement above is the
same with $s$ replaced by $s+1$.)

\begin{lemma}\label{lem:mellin}
Let $-\tfrac12<m<m'$ and $I=I_{m,m'}$ as in \eqref{eq:defI}, extended
by $I(x)=0$ for $1\le x<2$. For $\re s>m+1$,
\[
\int_1^\infty I(x)\,x^{-s-1}\,dx
=\frac{\kappa}{s(s+\kappa)}\sum_p\chi^*(p)\,p^{m-s}
=-\frac{\kappa\,P(s-m,\chi)}{s(s+\kappa)} .
\]
Moreover, for $\re w>1$,
\begin{equation}\label{eq:Pstructure}
P(w,\chi)=\log L(w,\chi)-\tfrac12\log\zeta(2w)+H_\chi(w),
\end{equation}
where $\log L(w,\chi):=\sum_p\sum_{k\ge1}\chi(p)^kp^{-kw}/k$,
$\log\zeta$ is defined likewise, and $H_\chi$ is analytic in
$\re w>\tfrac13$.
\end{lemma}

\begin{proof}
By Proposition~\ref{prop:formula}(a),
$I(x)=\sum_{p\le x}\chi^*(p)p^m(1-p^\kappa x^{-\kappa})$. For real
$\sigma>m+1$ (so $\sigma>\tfrac12>0$),
$\sum_pp^m\int_p^\infty x^{-\sigma-1}\,dx=\sigma^{-1}\sum_pp^{m-\sigma}
<\infty$, so by Fubini's theorem, for $\re s>m+1$,
\[
\int_1^\infty I(x)x^{-s-1}\,dx
=\sum_p\chi^*(p)p^m\int_p^\infty\bigl(x^{-s-1}-p^\kappa x^{-s-\kappa-1}\bigr)dx
=\sum_p\chi^*(p)p^{m-s}\Bigl(\frac1s-\frac1{s+\kappa}\Bigr).
\]
For \eqref{eq:Pstructure}: for $\re w>1$,
$\log L(w,\chi)=P(w,\chi)+\tfrac12\sum_p\chi(p)^2p^{-2w}
+\sum_{k\ge3}\frac1k\sum_p\chi(p)^kp^{-kw}$ and
$\sum_p\chi(p)^2p^{-2w}=\log\zeta(2w)-\sum_{k\ge2}\frac1k\sum_pp^{-2kw}
-\sum_{p\mid q}p^{-2w}$. Hence \eqref{eq:Pstructure} holds with
\[
H_\chi(w)=\frac12\sum_{k\ge2}\frac1k\sum_pp^{-2kw}
+\frac12\sum_{p\mid q}p^{-2w}
-\sum_{k\ge3}\frac1k\sum_p\chi(p)^kp^{-kw},
\]
and each of these series converges absolutely and locally uniformly
in $\re w>\tfrac13$.
\end{proof}

\begin{corollary}[Landau-type converse]\label{cor:landau}
Let $\chi$ be a real primitive character and let
\[
\beta_0(\chi):=\max\Bigl(\{\tfrac12\}\cup
\{\beta\in(\tfrac12,1):L(\beta,\chi)=0\}\Bigr)
\]
(a maximum over a finite set). Suppose that for some
$-\tfrac12<m<m'$ the function $I_{m,m'}(x)
=\kappa\int_2^xt^{-\kappa-1}D_{m'}(t)\,dt$ is of constant sign for all
sufficiently large $x$ (that is, $I_{m,m'}(x)\ge0$ for all
$x\ge x_1$, or $I_{m,m'}(x)\le0$ for all $x\ge x_1$). Then
$L(s,\chi)\ne0$ for $\re s>\beta_0(\chi)$. In particular, if
$L(\sigma,\chi)\ne0$ for all real $\sigma\in(\tfrac12,1)$, then
GRH$(\chi)$ holds.
\end{corollary}

\begin{proof}
The real zeros of $L(\cdot,\chi)$ in $[\tfrac12,1]$ are finite in
number and $L(\sigma,\chi)\ne0$ for real $\sigma\ge1$, so
$\beta_0=\beta_0(\chi)\in[\tfrac12,1)$ is well defined and
$L(\sigma,\chi)\ne0$ for real $\sigma>\beta_0$. The zeros of
$L(\cdot,\chi)$ form a closed discrete set disjoint from the ray
$(\beta_0,\infty)$, so there is a continuous function
$r:(\beta_0,\infty)\to(0,\infty)$ such that $L(w,\chi)\ne0$ on
$U:=\{\sigma+it:\sigma>\beta_0,\ |t|<r(\sigma)\}$. The region $U$ is
simply connected and contains part of the half-plane $\re w>1$, and
its intersection with any half-plane $\re w>c$ is connected; on $U$,
$L(\cdot,\chi)$ has an analytic logarithm which is real on the real
ray and therefore extends the function $\log L(w,\chi)$ of
Lemma~\ref{lem:mellin} from $U\cap\{\re w>1\}$, $\zeta(2w)$ is analytic and zero-free (as
$\re(2w)>1$), and $H_\chi$ is analytic. By \eqref{eq:Pstructure},
$P(\cdot,\chi)$ extends analytically to $U$, and
\[
F_1(s):=-\frac{\kappa\,P(s-m,\chi)}{s(s+\kappa)}
\]
is analytic on $U+m$ (there $\re s>m+\tfrac12>0$, so $s(s+\kappa)\ne0$).

Replacing $I$ by $-I$ if necessary, assume $I(x)\ge0$ for $x\ge x_1$.
The function $I$ is a finite sum of continuous functions with jumps
at primes; it is bounded and Riemann integrable on finite intervals.
Let $F$ and $\sigma_c$ be as in Lemma~\ref{lem:landau} with $A=I$. By
Lemma~\ref{lem:mellin}, $\sigma_c\le m+1$ and $F=F_1$ on
$\re s>m+1$. Suppose $\sigma_c>m+\beta_0$. Then $F$ is analytic on
$\re s>\sigma_c$, $F_1$ is analytic on $U+m$, and they agree on the
connected set $(U+m)\cap\{\re s>\sigma_c\}$, because they agree on
its nonempty open subset $(U+m)\cap\{\re s>m+1\}$. Hence $F_1$ is an
analytic continuation of $F$ to a neighbourhood of $s=\sigma_c\in
U+m$, contradicting Lemma~\ref{lem:landau}. Therefore
$\sigma_c\le m+\beta_0$: $F$ is analytic in $\re s>m+\beta_0$ and
equals $F_1$ on $\re s>m+1$. Consequently
$P(w,\chi)=-\kappa^{-1}(w+m)(w+m+\kappa)F(w+m)$ extends analytically
to $\re w>\beta_0$, and so does
$G(w):=P(w,\chi)+\tfrac12\log\zeta(2w)-H_\chi(w)$ (recall
$\beta_0\ge\tfrac12$). Since $e^{G(w)}=L(w,\chi)$ for $\re w>1$, the
identity theorem gives $e^{G}=L(\cdot,\chi)$ on $\re w>\beta_0$, so
$L(w,\chi)\ne0$ there. If $\beta_0=\tfrac12$, all nontrivial zeros
have $\beta\le\tfrac12$, hence $\beta=\tfrac12$ by the functional
equation.
\end{proof}

\begin{corollary}[equivalence with GRH for $q\in\{3,4\}$]
\label{cor:equiv}
Let $\chi\in\{\chi_3,\chi_4\}$ and $-\tfrac12<m<m'$.
\begin{enumerate}
\item[(a)] If $I_{m,m'}(x)$ is of constant sign for all sufficiently
large $x$, then GRH$(\chi)$ holds.
\item[(b)] If GRH$(\chi)$ holds and $\lambda_\chi(M,M')>0$, then
$I_{m,m'}(x)>0$ for all sufficiently large $x$.
\item[(c)] Suppose that the (unconditionally defined) number
$V_\chi(M'):=4M'\,\frac{\Lambda'}{\Lambda}(M'+\tfrac12,\chi)$ satisfies
$V_\chi(M')\le2$. Then the following three statements are
equivalent: (1) GRH$(\chi)$; (2) $D_m(x)>x^{-(m'-m)}D_{m'}(x)$ for
all sufficiently large $x$; (3) $I_{m,m'}(x)$ is of constant sign
for all sufficiently large $x$.
\end{enumerate}
In particular, since $V_{\chi_4}(\tfrac32)=1.3686\ldots$ and
$V_{\chi_3}(\tfrac32)=1.0070\ldots$, for $(m,m')=(0,1)$:
\[
\text{GRH}(\chi_4)\iff
\sum_{p\le x}\bigl(-\chi_4(p)\bigr)\Bigl(1-\frac px\Bigr)>0
\ \text{ for all sufficiently large }x ,
\]
the sum being $x^{-1}\int_2^x\bigl(\pi(t;4,3)-\pi(t;4,1)\bigr)dt$; and
the same holds for $\chi_3$ with $\pi(t;3,2)-\pi(t;3,1)$.
\end{corollary}

\begin{proof}
(a) For both characters the Dirichlet series converges for real
$\sigma>0$ to $L(\sigma,\chi)$ (bounded partial sums of $\chi$;
\cite[Ch.~4]{Dav00}), and grouping consecutive terms of the
convergent series gives
$L(\sigma,\chi_4)=\sum_{k\ge0}[(4k+1)^{-\sigma}-(4k+3)^{-\sigma}]>0$
and
$L(\sigma,\chi_3)=\sum_{k\ge0}[(3k+1)^{-\sigma}-(3k+2)^{-\sigma}]>0$
for $0<\sigma<1$. So $\beta_0(\chi)=\tfrac12$ and
Corollary~\ref{cor:landau} applies. (b) is Theorem~\ref{thm:A}(ii).
(c) NCZ$(\chi)$ holds unconditionally for these two characters. If
GRH$(\chi)$ holds, then $\sigma_{M'}^2=V_\chi(M')\le2$ by
Proposition~\ref{prop:WRvar}, so $\lambda_\chi(M,M')>0$ by
Theorem~\ref{thm:A}(iii) and (b) gives the second statement; the
second statement trivially implies the third; and the third implies
GRH$(\chi)$ by (a). The value $V_{\chi_4}(\tfrac32)=\sigma^2$ at
$m=1$ is $1.368555\ldots$ \cite[Table~2]{WR}; the value for $\chi_3$
was computed here in multiprecision floating point and is enclosed in
ball arithmetic in the auxiliary file of the direct certificate
(Sections~\ref{subsec:confirm}(a) and~\ref{sec:numerics}); both are far below $2$. The
identity for the sum is Proposition~\ref{prop:formula}(a) with
$\kappa=1$: $I_{0,1}(x)=\int_2^xt^{-2}D_1(t)\,dt
=\sum_{p\le x}\chi^*(p)(1-p/x)=x^{-1}\int_2^xD_0(t)\,dt$.
\end{proof}

\begin{remark}[the exceptional-zero caveat]\label{rem:exceptional}
For a general real primitive character Corollary~\ref{cor:landau}
gives only the zero-free half-plane $\re s>\beta_0(\chi)$: a real
zero $\beta\in(\tfrac12,1)$ of $L(s,\chi)$ produces a genuine real
singularity of $P(s-m,\chi)$ at $s=m+\beta$, and Landau's theorem
cannot see past it. An equivalence with GRH$(\chi)$ therefore
requires the absence of real zeros in $(\tfrac12,1)$ as a separate
input, which for $q=3,4$ is elementary and in general is open. For
general $\chi$ the statement that survives is: \emph{if
$L(\sigma,\chi)\ne0$ on $(\tfrac12,1)$ and $\lambda_\chi(M,M')>0$
under GRH$(\chi)$, then GRH$(\chi)$ is equivalent to the eventual
positivity of $I_{m,m'}$.}
\end{remark}

\begin{remark}[relation to the literature]\label{rem:literature}
(1) By Proposition~\ref{prop:formula}(a), $I_{m,m'}(x)$ is the Riesz
typical mean of order one,
$\sum_{p\le x}a_p\bigl(1-(p/x)^\kappa\bigr)$ with $a_p=\chi^*(p)p^m$,
of the series defining $D_m$, and simultaneously a power-weighted
integral mean of the heavier race $D_{m'}$; for $(m,m')=(0,1)$ it is
the Ces\`aro mean of Chebyshev's race. Smoothed versions of
Chebyshev's bias that hold for all large $x$ are classical. The
closest in spirit are the following (we quote statements at the level
at which we were able to check them; see Section~\ref{sec:gaps}).
Hardy--Littlewood and Landau \cite{HL18,Lan18} showed that Chebyshev's
assertion $\sum_{p>2}(-1)^{(p-1)/2}e^{-p/x}\to-\infty$ is equivalent
to GRH$(\chi_4)$; Fujii \cite{Fuj88} and Platt--Trudgian \cite{PT19}
extended this to the kernels $e^{-(p/x)^\alpha}$ for
$0<\alpha<4.19$ resp.\ $0<\alpha<20.40442$. Johnston \cite{Joh23}
showed that RH is equivalent to $\int_2^x(\pi(t)-\mathrm{li}(t))\,dt<0$
for all $x>2$; Washington \cite[Theorem~3]{Was24} showed that for
$0<k\le10.32$, RH is equivalent to
$\int_1^x(\pi_k(t)-\pi(t^{k+1}))\,dt<0$ for all sufficiently large
$x$, where $\pi_k(t)=\sum_{p\le t}p^k$; and Zhao \cite{Zhao25} showed
that RH is equivalent to the positivity, for all $X>2$, of the
integrated error terms in Mertens' theorems, with analogues for
quadratic characters. In these results, as in Theorem~\ref{thm:A}
with Corollaries~\ref{cor:landau}--\ref{cor:equiv}, a bias is
compared with an absolutely convergent sum over zeros and a
Landau-type argument gives the converse. In the three parametric
ones --- Fujii, Platt--Trudgian and Washington --- the result is
established only for a limited range of the smoothing parameter
($\alpha$, resp.\ $k$), determined by a numerical inequality between
the bias and the zero sum; this is the analogue of the condition
$\lambda_\chi(M,M')>0$. The results of Johnston and Zhao involve no
such parameter. Suzuki \cite{Suz25} is of a different type: he proves
that sign constancy of certain weighted summatory functions of
$\Lambda$ implies RH or GRH (one direction only), and establishes
that sign constancy on average. Liu--Wang--Shi \cite{LWS26} announce
a classification of the residue-class pairs for which an integrated
prime-counting error is eventually of one sign. We therefore claim no
novelty for the method of Theorem~\ref{thm:A}, and none for the
identity of Proposition~\ref{prop:formula}(a) as such, which is one
line of partial summation. What we did not find in the literature
searched is the use of that identity as a link \emph{between two
races of different weights}, and the resulting inheritance of leads
and monotonicity of upper and lower densities
(Theorem~\ref{thm:A}(ii), Corollary~\ref{cor:densities}).

(2) Hayani \cite{Hay25} obtains, under GRH alone, density-one
statements at the single critical weight $p^{-1/2}$; the present
GRH-only statements concern pairs of weights in $m>-\tfrac12$ and are
of a different nature.
\end{remark}

\begin{remark}[on the status of GRH]\label{rem:GRHstatus}
Theorem~\ref{thm:A}, Corollary~\ref{cor:densities} and the inequality
of Corollary~\ref{cor:sharp}(a) need GRH$(\chi)$ only (the equality
clause of Corollary~\ref{cor:sharp}(a) needs LI$(\chi)$ as well). A
proof of GRH for Dirichlet $L$-functions was announced on 20 November
2026 \cite{Alp26}, accompanied by a Lean formalization that
\cite{Lean26} reports to have checked; no human-readable proof has
been refereed. We have not examined either, and do not use them.
Every result of this article remains stated conditionally. If that
proof is confirmed, Theorem~\ref{thm:A},
Corollary~\ref{cor:densities} and the inequality in
Corollary~\ref{cor:sharp}(a) become unconditional, and
Corollary~\ref{cor:equiv} becomes an unconditional positivity
statement. Theorems~\ref{thm:B} and~\ref{thm:C} and the results of
Sections~\ref{sec:envelope}--\ref{sec:nonmono} would still need
NCZ$(\chi)$ (automatic for $q=3,4$) and, for their statements about
races, LI$(\chi)$; the computer-assisted ones would also still depend
on their trusted inputs.
\end{remark}

\section{Logarithmic weights: proof of Theorem~\ref{thm:C}}
\label{sec:thmC}

Assume GRH$(\chi)$ and NCZ$(\chi)$ and fix $m>-\tfrac12$. We write
$E(x):=\cE_{m,1}(x)=2M\,x^{-M}R_m(x)$ for $x\ge1$ (so $E(x)=0$ for
$x<2$).

\begin{lemma}[local mean-square bound]\label{lem:localms}
Under GRH$(\chi)$ and NCZ$(\chi)$,
\[
K_m:=\sup_{U\ge0}\int_U^{U+1}\bigl|E(e^u)\bigr|^2\,du<\infty .
\]
\end{lemma}

\begin{proof}
On $0\le u\le3$ the function $E(e^u)$ is bounded, since
$|R_m(x)|\le\sum_{p\le e^3}p^m\log p$ there. Let $U\ge2$,
$u\in[U,U+1]$ and $T:=e^{U+1}$. By Lemma~\ref{lem:WRexp} with $b=1$,
\[
E(e^u)=1+\re S_T(u)+\cE^{(1)}(u,T),\qquad
S_T(u):=\sum_{0<\gamma\le T}r_\gamma(M)e^{i\gamma u},
\]
and $|\cE^{(1)}(u,T)|\ll(U+1)e^{(U+1)/2}(2U+2)^2e^{-(U+1)}
+u^2e^{-\kappa_mu}+ue^{-u/6}+ue^{-c\sqrt u}\ll1$. It remains to bound
$\int_U^{U+1}|S_T(u)|^2\,du$ uniformly in $U$ and $T$. Let
$\mathcal F(v):=\bigl(\sin(v/2)/(v/2)\bigr)^2\ge0$; then
$\mathcal F(v)\ge(\sin(\tfrac12)/\tfrac12)^2>0.9$ for $|v|\le1$, and
$\int_{\mathbb R}\mathcal F(v)e^{i\xi v}\,dv=2\pi\max(0,1-|\xi|)$
(the Fej\'er kernel; \cite[Ch.~XV, \S2]{Fel71}). Hence, with
$U':=U+\tfrac12$,
\begin{align*}
0.9\int_U^{U+1}|S_T(u)|^2\,du
&\le\int_{\mathbb R}|S_T(u)|^2\mathcal F(u-U')\,du
=2\pi\sum_{\gamma,\gamma'}r_\gamma\overline{r_{\gamma'}}\,
e^{i(\gamma-\gamma')U'}\max(0,1-|\gamma-\gamma'|)\\
&\le2\pi\sum_{\substack{\gamma,\gamma'\\|\gamma-\gamma'|<1}}
|r_\gamma|\,|r_{\gamma'}|
\;\le\;2\pi\sum_{\gamma>0}|r_\gamma|^2\,\#\{\gamma':|\gamma'-\gamma|<1\},
\end{align*}
the sums being over ordinates in $(0,T]$ and the last step using
$|r_\gamma||r_{\gamma'}|\le\tfrac12(|r_\gamma|^2+|r_{\gamma'}|^2)$ and
symmetry. Since $|r_\gamma|^2\le16M^2\gamma^{-2}$ and
$\#\{\gamma':|\gamma'-\gamma|<1\}\ll\log(\gamma+3)$ by
\eqref{eq:zerosum}, the right-hand side is at most
$C\sum_{\gamma>0}\gamma^{-2}\log(\gamma+3)<\infty$ by
\eqref{eq:gammasums}, uniformly in $U$ and $T$. Finally
$|E|^2\le3(1+|S_T|^2+|\cE^{(1)}|^2)$.
\end{proof}

\begin{proof}[Proof of Theorem~\ref{thm:C}]
(i) Since $(\log p)^b=(\log p)^{b-1}\log p$, Riemann--Stieltjes
partial summation gives, for $x\ge2$ (recall $R_m=0$ on $[1,2)$),
\[
D_{m,b}(x)=\int_{3/2}^x(\log t)^{b-1}\,dR_m(t)
=(\log x)^{b-1}R_m(x)-(b-1)\int_2^xR_m(t)(\log t)^{b-2}\,\frac{dt}{t}.
\]
Multiplying by $2Mx^{-M}(\log x)^{1-b}$ and substituting
$R_m(t)=t^ME(t)/(2M)$,
\begin{equation}\label{eq:Cdecomp}
\cE_{m,b}(x)=E(x)-(b-1)\,\mathcal R_b(x),
\qquad
\mathcal R_b(e^u)=u^{1-b}\int_{\log2}^ue^{-M(u-v)}E(e^v)\,v^{b-2}\,dv .
\end{equation}
Let $u\ge2$ and split the integral at $v=u/2$. For
$\log2\le v\le u/2$ we have $e^{-M(u-v)}\le e^{-Mu/2}$ and
$v^{b-2}\le(\log2)^{b-2}+(u/2)^{b-2}$, and by the Cauchy--Schwarz
inequality on unit intervals
$\int_0^{u/2}|E(e^v)|\,dv\le(u/2+1)K_m^{1/2}$; this part of
$\mathcal R_b(e^u)$ is therefore
$\ll u^{1-b}e^{-Mu/2}\bigl(1+u^{b-2}\bigr)(u+1)\ll1/u$. For
$u/2\le v\le u$ we have $v^{b-2}\le2^{|b-2|}u^{b-2}$, and covering
$[u/2,u]$ by the intervals $[u-k-1,u-k]$, $k\ge0$, on which
$e^{-M(u-v)}\le e^{-Mk}$,
\[
\int_{u/2}^ue^{-M(u-v)}|E(e^v)|\,dv
\le\sum_{k\ge0}e^{-Mk}\Bigl(\int_{u-k-1}^{u-k}|E(e^v)|^2\,dv\Bigr)^{1/2}
\le\frac{K_m^{1/2}}{1-e^{-M}}
\]
(with $E(e^v):=0$ for $v<0$). Hence this part is at most
$2^{|b-2|}K_m^{1/2}(1-e^{-M})^{-1}u^{-1}$, and altogether
$|\mathcal R_b(e^u)|\ll1/u$ for $u\ge2$. This proves (i) for
$x\ge e^2$; for $3\le x\le e^2$ both sides are bounded.

(ii) By Theorem~\ref{thm:WRinterp}, $u\mapsto E(e^u)$ has the limiting
distribution $\nu_m$; by (i) and Lemma~\ref{lem:stability}(a) so does
$u\mapsto\cE_{m,b}(e^u)$.

(iii) Under LI$(\chi)$, $\nu_m$ is absolutely continuous, so the
boundary $\{0\}$ of $B:=(0,\infty)$ is a $\nu_m$-null set, and
Lemma~\ref{lem:stability}(b) shows that the logarithmic density of
the set
\[
\{x>1:\cE_{m,b}(x)>0\}=\{x>1:D_{m,b}(x)>0\}
\]
exists and equals $\nu_m((0,\infty))=\PP(X_M>0)=\delta_\chi(m)$.
\end{proof}

\begin{remark}\label{rem:Cgeneral}
(1) The proof uses of the factor $(\log p)^{b-1}$ only that the
function $w(v)=v^{b-1}$ is $C^1$ and positive on $[\log2,\infty)$
with $|w'(v)|\ll w(v)/v$ and $w(v)\ll w(u)$ for $u/2\le v\le u$, and
grows or decays at most polynomially; the same argument therefore
applies to weights $p^m\,w(\log p)$ with such a $w$, the normalisation
being $2Mx^{-M}(\log x)D(x)/w(\log x)$. We do not pursue this.
(2) Theorem~\ref{thm:C}(i) is a pointwise statement: the races with
different $b$ differ, after normalisation, by $O(1/\log x)$. They do
\emph{not} have the same sign for all large $x$; only the densities
agree. (3) NCZ$(\chi)$ enters through Lemma~\ref{lem:WRexp}. With a
central zero of order $n_0$ the constant $1$ in the expansions becomes
$1+2n_0$ (cf.\ Theorem~\ref{thm:A}(i)) and we expect the same
conclusion; since \cite{WR} assumes NCZ we have not written this out.
\end{remark}

\section{The derivative of the density and an envelope criterion}
\label{sec:envelope}

This section and the next are transcribed from two working notes in
the repository: \S\S1--3A of \texttt{BRIDGE-ATTEMPT.md} and the note
\texttt{ANALYTIC-LEG.md} of the envelope certificate (directories
\texttt{monotonicity/} and \texttt{monotonicity/envelope-cert/}).
There
Lemma~\ref{lem:diff}, Lemmas~\ref{lem:J} and~\ref{lem:L},
Proposition~\ref{prop:C0}, Theorem~\ref{thm:envelope},
Corollary~\ref{cor:cell} and Theorem~\ref{thm:E} are called Lemma~D,
Lemmas~J and~L, Proposition~C$_0$, Theorem~C, Corollary~C$'$ and
Theorem~E. The proofs were checked by internal referees
(Appendix~\ref{sec:review}); we re-checked each step while
transcribing and found no error.

\subsection{Setting}\label{subsec:envsetting}
Nothing arithmetic enters this section. Let
$0<\gamma_1\le\gamma_2\le\cdots$ be a finite or infinite sequence of
positive real numbers with \emph{at least nine terms} and
$\sum_j\gamma_j^{-2}<\infty$. For $M>0$ put
\begin{gather*}
c_j=c_j(M):=\frac{4M}{\sqrt{M^2+\gamma_j^2}},\qquad
\eps_j=\eps_j(M):=\frac{M^2}{M^2+\gamma_j^2},\qquad
S=S_M:=\sum_jc_j\cos\theta_j,\\
\Phi(t)=\Phi_M(t):=\prod_jJ_0(c_jt),\qquad
\Phi^{(j)}(t):=\prod_{i\ne j}J_0(c_it),\\
\sigma^2=\sigma^2(M):=\sum_j\frac{c_j^2}{2},\qquad
S_4=S_4(M):=\sum_jc_j^4,
\end{gather*}
with independent uniform phases $\theta_j$, and
\[
\hd(M):=\PP(S_M>-1)=\PP(1+S_M>0).
\]
Thus $c_1=\max_jc_j$ belongs to the lowest $\gamma$, and
$\Phi_M$ is the characteristic function of $S_M$ (the product
converges because $0\le1-J_0(z)\le z^2/4$ and $\sum_jc_j^2<\infty$).
For the positive ordinates of $L(s,\chi)$, $\chi$ real primitive, the
hypotheses hold under GRH$(\chi)$ and NCZ$(\chi)$ (by
\eqref{eq:gammasums}; there are infinitely many ordinates), and then
$\hd=\hd_\chi$, $\sigma^2(M)=\sigma_M^2$ in the notation of
Section~\ref{subsec:setting}. A zero of multiplicity $k$ is $k$ equal
terms of the sequence, with $k$ independent phases; this is the model
of Section~\ref{subsec:setting}, and it is the race only under LI,
when all zeros are simple.

\subsection{Facts about \texorpdfstring{$J_0$}{J0}}\label{subsec:bessel}
We use the following, for real $z$.
\begin{itemize}
\item[(B1)] $|J_0(z)|\le1$, \ $0\le1-J_0(z)\le z^2/4$, \ $J_0'=-J_1$
and $|J_1(z)|\le|z|/2$. (All from Bessel's integral
$J_0(z)=\frac1{2\pi}\int_0^{2\pi}\cos(z\cos\theta)\,d\theta$ and
$J_1(z)=\frac1{2\pi}\int_0^{2\pi}\sin(z\cos\theta)\cos\theta\,d\theta$.)
\item[(B2)] $|J_0(z)|\le\frac8\pi z^{-1/2}$ for $z>0$
\cite[Lemma~4, (J2)]{WR}. (It also follows from the sharper classical
bound $|J_0(z)|\le\sqrt{2/(\pi z)}$ of \cite[\S13.74]{Wat44}.)
\item[(B3)] The first positive zero of $J_0$ satisfies $j_{0,1}>2.4$,
and for $|z|<j_{0,1}$
\[
-\log J_0(z)=\sum_{n\ge1}\frac{\lambda_n}{n}\,z^{2n}
\qquad\text{with all }\lambda_n>0,\ \lambda_1=\tfrac14
\]
\cite[Lemma~14(i), (iii)]{WR}; this is the Hadamard product of $J_0$,
and $\lambda_n=\sum_kj_{0,k}^{-2n}$.
\item[(B4)] $J_0$ is strictly decreasing on $[0,2\sqrt2\,]$, and
$\sup_{z\ge2\sqrt2}|J_0(z)|\le0.446$ \cite[Lemma~15(i), (ii)]{WR}.
\end{itemize}
We also use the elementary inequality
$|1-\prod_ka_k|\le\sum_k|1-a_k|$ for real numbers $|a_k|\le1$.

\subsection{The derivative identity}

\begin{lemma}[derivative identity at every weight]\label{lem:diff}
In the setting of Section~\ref{subsec:envsetting}, $\hd$ is
continuously differentiable on $(0,\infty)$, the law of $S_M$ has a
bounded continuous even density $f_M$, and for every $M>0$
\begin{equation}\label{eq:deriv}
\begin{aligned}
\hd'(M)&=\mathcal D(M):=\frac1\pi\int_0^\infty
\partial_M\Phi_M(t)\,\frac{\sin t}{t}\,dt,\\
\partial_M\Phi_M(t)&=-\sum_{j}c_j'(M)\,t\,J_1\bigl(c_j(M)t\bigr)\,
\Phi_M^{(j)}(t),\qquad c_j'(M)=\frac{4\gamma_j^2}{(M^2+\gamma_j^2)^{3/2}},
\end{aligned}
\end{equation}
and
\begin{equation}\label{eq:Did}
\begin{aligned}
-M\,\hd'(M)&=f_M(1)+R(M),\qquad
f_M(1)=\frac1\pi\int_0^\infty\Phi_M(t)\cos t\,dt,\\
R(M)&:=-\frac1\pi\sum_j\eps_jc_j\int_0^\infty J_1(c_jt)\,
\Phi_M^{(j)}(t)\,\sin t\,dt ,
\end{aligned}
\end{equation}
all integrals and series converging absolutely.
\end{lemma}

For the ordinates of $L(s,\chi)$ we write $\mathcal D_\chi$ for
$\mathcal D$. Since $\hd$ is $C^1$ on the open half-line, it has a
two-sided derivative, given by \eqref{eq:deriv}, at every point of
every closed interval $[M_a,M_b]\subset(0,\infty)$, endpoints
included. Probabilistically $R(M)=\EE[Z;S=-1]$ with
$Z=\sum_j\eps_jc_j\cos\theta_j$, so that \eqref{eq:Did} is the
rigorous form of the formal identity of
Remark~\ref{rem:infinitesimal}; the proof does not use this reading.
The number nine is a convenient sufficient number of terms, not a
sharp one.

\begin{proof}
\emph{Step 1 (the product is $C^1$ in $(M,t)$).} Fix $0<M_a\le M_b$
and $T>0$, and let $M\in[M_a,M_b]$, $|t|\le T$. We have
$c_j\le4M_b/\gamma_j$ and $c_j'=dc_j/dM=c_j(1-\eps_j)/M$, so
$0<c_j'\le c_j/M_a$. Put $A:=16M_b^2\sum_j\gamma_j^{-2}$ and
$A_N:=16M_b^2\sum_{j>N}\gamma_j^{-2}$; then $\sum_jc_j^2\le A$,
$\sum_{j>N}c_j^2\le A_N$, and $A_N\to0$. Let
$P_N(M,t):=\prod_{j\le N}J_0(c_jt)$, a $C^1$ function with
\[
\partial_MP_N=-\sum_{j\le N}c_j'\,t\,J_1(c_jt)\!\!\prod_{i\le N,\,i\ne j}\!\!J_0(c_it),
\qquad
\partial_tP_N=-\sum_{j\le N}c_j\,J_1(c_jt)\!\!\prod_{i\le N,\,i\ne j}\!\!J_0(c_it).
\]
By (B1) and the inequality for products, the infinite products
$\Phi_M(t)$ and $\Phi^{(j)}_M(t)$ converge and
\[
|P_N-\Phi_M(t)|\le\tfrac14T^2A_N,\qquad
\Bigl|\prod_{i\le N,\,i\ne j}J_0(c_it)-\Phi_M^{(j)}(t)\Bigr|\le\tfrac14T^2A_N
\quad(j\le N).
\]
By (B1), the $j$-th summand of $\partial_MP_N$ is at most
$c_jc_j't^2/2\le c_j^2T^2/(2M_a)$ in modulus and that of
$\partial_tP_N$ at most $c_j^2T/2$. Hence, with the absolutely
convergent series
\[
R_M:=-\sum_jc_j'tJ_1(c_jt)\Phi^{(j)}_M(t),\qquad
R_t:=-\sum_jc_jJ_1(c_jt)\Phi^{(j)}_M(t),
\]
we have
\[
|\partial_MP_N-R_M|\le\frac{T^2}{2M_a}\Bigl(\frac{T^2A_N}{4}\,A+A_N\Bigr),
\qquad
|\partial_tP_N-R_t|\le\frac{T}{2}\Bigl(\frac{T^2A_N}{4}\,A+A_N\Bigr),
\]
and both tend to $0$ uniformly on $[M_a,M_b]\times[-T,T]$. Since
$P_N\to\Phi$ and the gradients of the $C^1$ functions $P_N$ converge
locally uniformly, $\Phi$ is $C^1$ in $(M,t)$ on
$(0,\infty)\times\mathbb R$, with
\begin{equation}\label{eq:dPhi}
\partial_M\Phi_M(t)=-\sum_jc_j'\,t\,J_1(c_jt)\,\Phi_M^{(j)}(t),
\qquad
\Phi_M'(t)=-\sum_jc_j\,J_1(c_jt)\,\Phi_M^{(j)}(t).
\end{equation}
(For a finite sequence this step is the product rule.)

\emph{Step 2 (decay).} Put $\kappa:=c_9(M_a)>0$. For
$M\in[M_a,M_b]$ and $i\le9$, $c_i(M)\ge c_i(M_a)\ge\kappa$, because
$c_i(M)$ is increasing in $M$ and nonincreasing in $i$. In $\Phi$ and
in every $\Phi^{(j)}$ at least eight of the first nine factors are
present, and by (B1), (B2) each of them is at most
$\min(1,\frac8\pi(\kappa t)^{-1/2})$ in modulus. With
$C:=(8/\pi)^8\kappa^{-4}$, for $t>0$, all $j$ and all
$M\in[M_a,M_b]$,
\begin{equation}\label{eq:decay}
\begin{gathered}
|\Phi_M(t)|,\ |\Phi_M^{(j)}(t)|\le\min(1,Ct^{-4}),\\
|\partial_M\Phi_M(t)|\le\frac{t^2}{2M_a}\,A\,\min(1,Ct^{-4}),\qquad
|\Phi_M'(t)|\le\frac t2\,A\,\min(1,Ct^{-4}).
\end{gathered}
\end{equation}

\emph{Step 3 (differentiation under the integral).} By
\eqref{eq:decay}, $\Phi_M\in L^1(\mathbb R)$; it is real and even.
Hence $S_M$ has the bounded continuous even density
$f_M(x)=\frac1\pi\int_0^\infty\Phi_M(t)\cos(tx)\,dt$, and
\[
\hd(M)-\frac12=\frac1\pi\int_0^\infty\Phi_M(t)\,\frac{\sin t}{t}\,dt
\]
(\cite[Lemma~7(b)]{WR}; its proof uses only that $\Phi_M$ is an
integrable real even characteristic function). By \eqref{eq:decay}
and $|\sin t|/t\le\min(1,1/t)$, the function
$\frac{A}{2M_a}\min(t^2,Ct^{-3})$ is an integrable dominant for
$\partial_M\Phi_M(t)\sin t/t$, uniform on $[M_a,M_b]$. Therefore
$\hd$ is differentiable on $[M_a,M_b]$ with derivative
$\mathcal D(M)$, and $\mathcal D$ is continuous because
$\partial_M\Phi_M(t)$ is continuous in $M$ and dominated. As
$0<M_a\le M_b$ were arbitrary, $\hd\in C^1(0,\infty)$ and
\eqref{eq:deriv} holds.

\emph{Step 4 (algebra).} By \eqref{eq:dPhi} and
$c_j'=c_j(1-\eps_j)/M$,
\[
\partial_M\Phi_M(t)=\frac tM\,\Phi_M'(t)
+\frac tM\sum_j\eps_jc_j\,J_1(c_jt)\,\Phi_M^{(j)}(t),
\]
so
\[
M\,\hd'(M)=\frac1\pi\int_0^\infty\Phi_M'(t)\sin t\,dt
+\frac1\pi\int_0^\infty\sum_j\eps_jc_j\,J_1(c_jt)\,\Phi_M^{(j)}(t)\,\sin t\,dt .
\]
Integrating by parts, with vanishing boundary terms by
\eqref{eq:decay}, the first integral equals
\[
-\int_0^\infty\Phi_M(t)\cos t\,dt=-\pi f_M(1).
\]
In the second, sum and integral may be exchanged, because by (B1) and \eqref{eq:decay}
\[
\sum_j\eps_jc_j\int_0^\infty|J_1(c_jt)|\,|\Phi_M^{(j)}(t)|\,dt
\le\Bigl(\sum_j\frac{c_j^2}{2}\Bigr)\int_0^\infty t\min(1,Ct^{-4})\,dt<\infty .
\]
This gives \eqref{eq:Did}.
\end{proof}

\subsection{Envelopes for \texorpdfstring{$J_0$}{J0}}
Fix the constants
\[
\rho:=0.446,\qquad z_*:=2\sqrt{\log(1/\rho)}=1.79714\ldots,\qquad
z_{\mathrm{far}}:=\Bigl(\frac{8}{\pi\rho}\Bigr)^2=32.59\ldots,
\]
and define
\[
b(z):=\begin{cases}e^{-z^2/4}&(0\le z\le z_*),\\
\rho&(z_*\le z\le z_{\mathrm{far}}),\\
\frac8\pi z^{-1/2}&(z\ge z_{\mathrm{far}}),\end{cases}
\qquad
\beta(z_1):=\frac{-\log J_0(z_1)-z_1^2/4}{z_1^4}\quad(0<z_1<j_{0,1}).
\]
The function $b$ is continuous and nonincreasing, with $0<b\le1$
(the three pieces agree at $z_*$ and at $z_{\mathrm{far}}$), and
$z_*<2.4<j_{0,1}$.

\begin{lemma}\label{lem:J}
$|J_0(z)|\le b(z)$ for all $z\ge0$.
\end{lemma}

\begin{proof}
For $0\le z<j_{0,1}$, (B3) gives $0<J_0(z)\le e^{-z^2/4}$; this
covers $z\le z_*$. On $[z_*,2\sqrt2\,]$, $J_0$ is decreasing by (B4),
so $J_0(z)\le J_0(z_*)\le e^{-z_*^2/4}=\rho$ and
$J_0(z)\ge J_0(2\sqrt2)\ge-\rho$, the last by (B4). For $z\ge2\sqrt2$,
$|J_0(z)|\le0.446=\rho$ by (B4). Finally $|J_0(z)|\le\frac8\pi z^{-1/2}$
by (B2), and this is $\le\rho$ exactly when $z\ge z_{\mathrm{far}}$.
\end{proof}

\begin{lemma}\label{lem:L}
For $0\le z\le z_1<j_{0,1}$: \
$J_0(z)\ge\exp\bigl(-z^2/4-\beta(z_1)z^4\bigr)$, and $\beta(z_1)>0$.
\end{lemma}

\begin{proof}
By (B3), $-\log J_0(z)-z^2/4=z^4\sum_{n\ge2}\frac{\lambda_n}{n}z^{2n-4}
\le z^4\sum_{n\ge2}\frac{\lambda_n}{n}z_1^{2n-4}=\beta(z_1)z^4$, by the
positivity of the $\lambda_n$.
\end{proof}

\subsection{The criterion}

\begin{proposition}\label{prop:C0}
In the setting of Section~\ref{subsec:envsetting} let $M>0$,
$z_1\in(0,j_{0,1})$ and $0<t_1\le\min(\pi/2,\,z_1/c_1(M))$. Let
$\overline{\sigma}^2\ge\sigma^2(M)$ and $\overline S_4\ge S_4(M)$ be
real numbers, and let $W,W^1:[0,\infty)\to[0,\infty)$ be measurable
functions with
\[
|\Phi_M(t)|\le W(t)\quad(t\ge t_1),\qquad
|\Phi_M^{(j)}(t)|\le W^1(t)\quad(t\ge0,\ \text{every }j).
\]
Then
\[
-M\,\hd'(M)\;\ge\;\frac1\pi\int_0^{t_1}
\exp\bigl(-\tfrac12\overline{\sigma}^2t^2-\beta(z_1)\overline S_4t^4\bigr)\cos t\,dt
-\frac1\pi\int_{t_1}^\infty W(t)\,dt
-\frac{\overline S_4}{32\pi}\int_0^\infty t^2\,W^1(t)\,dt .
\]
\end{proposition}

\begin{proof}
By Lemma~\ref{lem:diff}, $-M\hd'(M)=f_M(1)+R(M)$.

\emph{The term $R$.} By (B1) and $|\sin t|\le t$,
\[
|R(M)|\le\frac1{2\pi}\sum_j\eps_jc_j^2\int_0^\infty t^2\,|\Phi_M^{(j)}(t)|\,dt
\le\frac1{2\pi}\Bigl(\sum_j\eps_jc_j^2\Bigr)\int_0^\infty t^2W^1(t)\,dt
\le\frac{\overline S_4}{32\pi}\int_0^\infty t^2W^1(t)\,dt,
\]
because $\eps_jc_j^2=16M^4/(M^2+\gamma_j^2)^2=c_j^4/16$, i.e.\
$\sum_j\eps_jc_j^2=S_4/16$.

\emph{The term $f_M(1)$.} For $0\le t\le t_1$ every argument
satisfies $c_jt\le c_1t_1\le z_1$, so by Lemma~\ref{lem:L}
\begin{align*}
\Phi_M(t)&\ge\exp\Bigl(-\sum_j\frac{c_j^2t^2}{4}-\beta(z_1)\sum_jc_j^4t^4\Bigr)
=\exp\bigl(-\tfrac12\sigma^2t^2-\beta(z_1)S_4t^4\bigr)\\
&\ge\exp\bigl(-\tfrac12\overline{\sigma}^2t^2-\beta(z_1)\overline S_4t^4\bigr),
\end{align*}
and $\cos t\ge0$ there. Hence
\begin{align*}
\pi f_M(1)&=\int_0^{t_1}\Phi_M\cos t\,dt+\int_{t_1}^\infty\Phi_M\cos t\,dt\\
&\ge\int_0^{t_1}\exp\bigl(-\tfrac12\overline{\sigma}^2t^2-\beta(z_1)\overline S_4t^4\bigr)\cos t\,dt
-\int_{t_1}^\infty W(t)\,dt .\qedhere
\end{align*}
\end{proof}

\begin{theorem}[envelope criterion]\label{thm:envelope}
In the setting of Section~\ref{subsec:envsetting} let $M>0$ and
$n\ge9$, and consider the $n$ lowest terms
$\gamma_1\le\cdots\le\gamma_n$ of the sequence. Let $v\ge0$ be either
\begin{itemize}
\item[(a)] $v=0$, or
\item[(b)] $v=v_n:=\sigma^2(M)-\sum_{i\le n}c_i^2/2
=\sum_{i>n}c_i^2/2$.
\end{itemize}
Put
\[
B(t):=\prod_{i=1}^nb(c_it)\cdot
\exp\Bigl(-\frac v2\,\min\bigl(t,\,z_*/c_n\bigr)^2\Bigr),
\qquad B^1(t):=\frac{B(t)}{b(c_1t)},
\]
choose $z_1\in(0,j_{0,1})$ and $0<t_1\le\min(\pi/2,z_1/c_1)$, and
define
\begin{align*}
\mathrm{MAIN}&:=\frac1\pi\int_0^{t_1}
\exp\bigl(-\tfrac12\sigma^2t^2-\beta(z_1)S_4t^4\bigr)\cos t\,dt,\\
\mathrm{TAILF}&:=\frac1\pi\int_{t_1}^\infty B(t)\,dt,\qquad
\mathrm{NUM}:=\frac{S_4}{32\pi}\int_0^\infty t^2B^1(t)\,dt .
\end{align*}
Then
\[
-M\,\hd'(M)\;\ge\;\mathrm{MAIN}-\mathrm{TAILF}-\mathrm{NUM}.
\]
In particular $\hd'(M)<0$ whenever
$\mathrm{MAIN}-\mathrm{TAILF}-\mathrm{NUM}>0$.
\end{theorem}

(In case (b) the sequence is ordered, so every term with index $>n$
has $c_i\le c_n$; for the ordinates of an $L$-function this means
that $\gamma_1,\dots,\gamma_n$ must be \emph{all} the ordinates in
$(0,\gamma_n]$, with multiplicity.)

\begin{proof}
We verify the hypotheses of Proposition~\ref{prop:C0} with $W=B$,
$W^1=B^1$, $\overline{\sigma}^2=\sigma^2$, $\overline S_4=S_4$.

\emph{$|\Phi|\le B$.} By Lemma~\ref{lem:J} and $|J_0|\le1$,
$|\Phi(t)|\le\prod_{i\le n}b(c_it)$, which is case (a). In case (b),
every $i>n$ has $c_i\le c_n$. For $t\le z_*/c_n$ then $c_it\le z_*$
and $b(c_it)=\exp(-c_i^2t^2/4)$, and the product of these over $i>n$
is $\exp(-v_nt^2/2)$. For $t>z_*/c_n$,
$b(c_it)\le b(c_iz_*/c_n)$ by monotonicity, with product
$\exp(-v_nz_*^2/(2c_n^2))$. So $|\Phi(t)|\le B(t)$ in both cases.

\emph{$|\Phi^{(j)}|\le B^1$ for every $j$.} For $j=1$ this is the same
computation with the first factor omitted. For $2\le j\le n$ the
bound obtained is $B(t)/b(c_jt)$, and $b(c_1t)\le b(c_jt)$ since
$c_1\ge c_j$. For $j>n$ the bound obtained is $B(t)$ in case (a), and
$B(t)/b\bigl(c_j\min(t,z_*/c_n)\bigr)$ in case (b) (the factor
$\exp(-c_j^2\min(t,z_*/c_n)^2/4)$ is absent from the product); and
$b(c_1t)\le b\bigl(c_j\min(t,z_*/c_n)\bigr)\le1$ because
$c_j\min(t,z_*/c_n)\le c_1t$.

Since $B^1$ contains at least eight factors $b(c_it)$, we have
$t^2B^1(t)=O(t^{-2})$ and all integrals converge.
Proposition~\ref{prop:C0} gives the claim.
\end{proof}

\begin{corollary}[cell form]\label{cor:cell}
Let $0<M_a\le M_b$. With the notation of Theorem~\ref{thm:envelope},
evaluate $\mathrm{MAIN}$ with $\sigma^2(M_b)$, $S_4(M_b)$ and
$t_1\le\min(\pi/2,z_1/c_1(M_b))$; evaluate $B$ and $B^1$ with
$c_i(M_a)$ $(i\le n)$ and, in case \textup{(b)}, with $v_n(M_a)$ and
$c_n(M_a)$; evaluate the prefactor of $\mathrm{NUM}$ with $S_4(M_b)$.
Call the results $\mathrm{MAIN}_{\mathrm{cell}}$,
$\mathrm{TAILF}_{\mathrm{cell}}$, $\mathrm{NUM}_{\mathrm{cell}}$,
$B_{\mathrm{cell}}$, $B^1_{\mathrm{cell}}$. Then for every
$M\in[M_a,M_b]$
\[
-M\,\hd'(M)\;\ge\;\mathrm{MAIN}_{\mathrm{cell}}
-\mathrm{TAILF}_{\mathrm{cell}}-\mathrm{NUM}_{\mathrm{cell}} ,
\]
so $\hd'<0$ on the whole closed cell if the right-hand side is
positive.
\end{corollary}

\begin{proof}
Fix $M\in[M_a,M_b]$; we verify the hypotheses of
Proposition~\ref{prop:C0} with $W=B_{\mathrm{cell}}$,
$W^1=B^1_{\mathrm{cell}}$, $\overline{\sigma}^2=\sigma^2(M_b)$,
$\overline S_4=S_4(M_b)$.

(i) \emph{Monotone quantities.} Each
$c_i(M)=4/\sqrt{1+\gamma_i^2/M^2}$ is increasing in $M$; hence so are
$\sigma^2(M)$ and $S_4(M)$, termwise. So
$\sigma^2(M)\le\sigma^2(M_b)$, $S_4(M)\le S_4(M_b)$,
$c_i(M)\ge c_i(M_a)$, and $t_1\le z_1/c_1(M_b)\le z_1/c_1(M)$: the
chosen $t_1$ is admissible for $M$.

(ii) \emph{Listed factors.} For $i\le n$:
$|J_0(c_i(M)t)|\le b(c_i(M)t)\le b(c_i(M_a)t)$, by Lemma~\ref{lem:J}
and because $b$ is nonincreasing.

(iii) \emph{Unlisted factors, case (b).} Put
$\mu:=\min(t,z_*/c_n(M_a))$. For any set $I$ of indices $>n$, using
that $b$ is nonincreasing (twice) and then
$c_i(M_a)\mu\le c_n(M_a)\mu\le z_*$,
\[
\prod_{i\in I}|J_0(c_i(M)t)|\le\prod_{i\in I}b(c_i(M_a)t)
\le\prod_{i\in I}b(c_i(M_a)\mu)
=\exp\Bigl(-\frac{\mu^2}{4}\sum_{i\in I}c_i(M_a)^2\Bigr).
\]
For $I=\{i>n\}$ the exponent is $-\frac{\mu^2}{2}v_n(M_a)$; for
$I=\{i>n\}\setminus\{j\}$ the bound is
\[
\exp\bigl(-\tfrac{\mu^2}{2}v_n(M_a)\bigr)\big/b(c_j(M_a)\mu).
\]
(In case (a) the unlisted factors are bounded by $1$.)

(iv) \emph{Assembly.} From (ii)--(iii),
$|\Phi_M(t)|\le B_{\mathrm{cell}}(t)$. For $j\le n$,
\[
|\Phi_M^{(j)}(t)|\le\frac{B_{\mathrm{cell}}(t)}{b(c_j(M_a)t)}
\le\frac{B_{\mathrm{cell}}(t)}{b(c_1(M_a)t)}=B^1_{\mathrm{cell}}(t),
\]
because $c_1(M_a)\ge c_j(M_a)$. For $j>n$,
\[
|\Phi_M^{(j)}(t)|\le\frac{B_{\mathrm{cell}}(t)}{b(c_j(M_a)\mu)}
\le B^1_{\mathrm{cell}}(t),
\]
because $c_j(M_a)\mu\le c_1(M_a)t$ (in case (a) simply
$|\Phi_M^{(j)}|\le B_{\mathrm{cell}}\le B^1_{\mathrm{cell}}$).
Proposition~\ref{prop:C0} now gives the claim.
\end{proof}

\begin{remark}\label{rem:K}
(1) Where $\mathrm{MAIN}>\mathrm{TAILF}$, the ratio
$K:=\mathrm{NUM}/(\mathrm{MAIN}-\mathrm{TAILF})$ measures the margin
of the criterion, which holds when $K<1$. (The inequality $K<1$ alone
proves nothing when $\mathrm{MAIN}\le\mathrm{TAILF}$.) In the Gaussian
regime $\mathrm{MAIN}\approx(2\pi\sigma^2)^{-1/2}$ and
$\mathrm{NUM}\approx\frac{S_4}{32\sigma^2}(2\pi\sigma^2)^{-1/2}$, so
$K\approx S_4/(32\sigma^2)=\sum_j\eps_jc_j^2/\sum_jc_j^2\le\tfrac12$
once $S_4\le16\sigma^2$: heuristically the criterion needs no better
than a factor $2$, which is why it has room where bounds that track
the oscillation of the inversion integral do not. This is a
heuristic; the statements used below are
Theorems~\ref{thm:envelope} and~\ref{thm:E} and
Corollary~\ref{cor:cell}.

(2) No cancellation is used anywhere: the proof bounds
$|J_1(z)|\le z/2$, $|\sin t|\le t$, $|J_0|\le b$, and uses the
positivity of the Taylor coefficients of $-\log J_0$.
\end{remark}

\subsection{Closed forms for the ordinates of an \texorpdfstring{$L$}{L}-function}

\begin{lemma}\label{lem:S4}
Let $\chi$ be a real primitive character mod $q$ with
$\chi(-1)=(-1)^a$, assume GRH$(\chi)$ and NCZ$(\chi)$, and for $M>0$
put $s=M+\tfrac12$ and
\[
h(M):=\frac{\Lambda'}{\Lambda}(s,\chi)
=\tfrac12\log\frac q\pi+\tfrac12\,\psi\Bigl(\frac{s+a}{2}\Bigr)+\frac{L'}{L}(s,\chi).
\]
Then, for the sequence of positive ordinates of $L(s,\chi)$,
\begin{gather*}
\sigma^2(M)=4M\,h(M),\qquad
S_4(M)=64M\,h(M)-64M^2\,h'(M),\\
h'(M)=\tfrac14\,\psi'\Bigl(\frac{s+a}{2}\Bigr)+\Bigl(\frac{L'}{L}\Bigr)'(s,\chi).
\end{gather*}
\end{lemma}

\begin{proof}
The first identity is Proposition~\ref{prop:WRvar}; it says
$\sum_\gamma(M^2+\gamma^2)^{-1}=h(M)/(2M)$. The series
$\sum_\gamma2M(M^2+\gamma^2)^{-2}$ of termwise derivatives converges
locally uniformly on $(0,\infty)$, so
$-2M\sum_\gamma(M^2+\gamma^2)^{-2}=h'(M)/(2M)-h(M)/(2M^2)$, i.e.\
$\sum_\gamma(M^2+\gamma^2)^{-2}=h/(4M^3)-h'/(4M^2)$, and
$S_4=256M^4\sum_\gamma(M^2+\gamma^2)^{-2}$.
\end{proof}

\section{A uniform large-weight theorem, and the General Theorem}
\label{sec:largeweight}

Throughout this section $\chi$ is a real primitive Dirichlet character
mod $q\ge3$ with $\chi(-1)=(-1)^a$, GRH$(\chi)$ and NCZ$(\chi)$ are
assumed, the sequence of Section~\ref{subsec:envsetting} is the
sequence of positive ordinates of $L(s,\chi)$ (with multiplicity), and
$h$ is the function of Lemma~\ref{lem:S4}. Put
\begin{gather*}
\tau:=\frac{z_*}{4}=0.44928\ldots,\qquad
c_\kappa:=\frac{4}{\sqrt{26}},\qquad
T_c:=\Bigl(\frac8\pi\Bigr)^2\frac1{c_\kappa}=8.2662\ldots,\\
\beta_*:=\beta(z_*)=\frac{\log(\rho/J_0(z_*))}{z_*^4},
\end{gather*}
so that $\rho=e^{-z_*^2/4}$ and $e^{-x\tau^2/2}=\rho^{x/8}$, and for
$x>8$
\begin{multline}\label{eq:defG}
G(x):=1-\frac{\frac12+48\beta_*}{x}-\frac12\Bigl(1-\frac8x\Bigr)^{-3/2}
-\sqrt{2\pi x}\,\Bigl[\frac{\rho^{x/8}}{\pi x\tau}
+\frac{9}{7\pi}\,T_c\,\rho^{(x-72)/8}\Bigr]\\
-\frac{x^{3/2}}{\sqrt{2\pi}}\cdot\frac43\,T_c^3\,\rho^{(x-80)/8}.
\end{multline}

\begin{theorem}[uniform large-weight theorem]\label{thm:E}
Let $\chi$ be a real primitive Dirichlet character and assume
GRH$(\chi)$ and $L(\tfrac12,\chi)\ne0$. Then for every $M\ge40$
\[
-M\,\hd_\chi'(M)\;\ge\;\frac{G(\sigma_M^2)}{\sqrt{2\pi}\,\sigma_M}
\;\ge\;\frac{0.3219}{\sqrt{2\pi}\,\sigma_M}\;>\;0 .
\]
In particular $\hd_\chi$ is strictly decreasing on $[40,\infty)$.
\end{theorem}

No ordinate of any zero is used. The proof reduces the theorem to the
four numerical inequalities \textup{(E2)--(E5)} below, between values
of elementary functions, the digamma function and $J_0$ at explicit
points. These were evaluated in ball arithmetic (program
\texttt{analytic\_leg.py} in \texttt{monotonicity/envelope-cert/},
$128$-bit Arb balls, lower endpoints rounded down and upper endpoints
rounded up) and recomputed independently at $30$ digits by an internal
referee; the certified values are
\begin{equation}\label{eq:Evalues}
\begin{gathered}
\textup{(E2)}\ \ \frac1{2(s+1)}-Z_2(s)\ \ge\ 0.0120481\ \ (s=40.5);\qquad
\textup{(E3)}\ \ \sigma_0^2\ \ge\ 236.97110;\\
\textup{(E4)}\ \ N_\chi(200)\ \ge\ 11.3751;\qquad
\textup{(E5)}\ \ G(\sigma_0^2)\ \ge\ 0.3219715,\\
\text{with }\beta_*\le0.025553,\ T_c\le8.266220\text{ in (E5)},
\end{gathered}
\end{equation}
in the notation introduced in the proof.

\begin{proof}
Write $s=M+\tfrac12$ and $L_2:=\log2$. For real $s\ge4$,
\begin{align*}
\Bigl|\frac{L'}{L}(s,\chi)\Bigr|&\le\sum_{n\ge2}\frac{\log n}{n^{s}}
\le Z_1(s):=L_2\,2^{-s}+2^{1-s}\Bigl(\frac{L_2}{s-1}+\frac1{(s-1)^2}\Bigr),\\
\Bigl|\Bigl(\frac{L'}{L}\Bigr)'(s,\chi)\Bigr|&\le\sum_{n\ge2}\frac{\log^2n}{n^{s}}
\le Z_2(s):=L_2^2\,2^{-s}+2^{1-s}\Bigl(\frac{L_2^2}{s-1}+\frac{2L_2}{(s-1)^2}
+\frac{2}{(s-1)^3}\Bigr)
\end{align*}
(the first term plus the integral over $[2,\infty)$; the integrands
$x^{-s}\log x$ and $x^{-s}\log^2x$ are decreasing on $[2,\infty)$
because $\log2>2/s$). $Z_1$ and $Z_2$ are decreasing in $s$.

\emph{Step 0.} Every $c_j=4/\sqrt{1+\gamma_j^2/M^2}$ is $<4$, so
$c_j^2/2\le8$.

\emph{Step 1: $S_4\le16\sigma^2$ for $M\ge40$.} By
Lemma~\ref{lem:S4}, $S_4=16\sigma^2-64M^2h'$. Since
$\psi'(x)=\sum_{k\ge0}(x+k)^{-2}\ge\int_x^\infty u^{-2}du=1/x$ and
$\psi'$ is decreasing,
\[
h'(M)\ge\tfrac14\psi'\Bigl(\frac{s+1}{2}\Bigr)-Z_2(s)
\ge\frac{1}{2(s+1)}-Z_2(s).
\]
The function $(s+1)Z_2(s)$ is decreasing for $s\ge4$ (a product of
positive decreasing functions), so $\frac1{2(s+1)}>Z_2(s)$ for all
$s\ge40.5$ as soon as it holds at $s=40.5$; that is (E2). Hence
$h'(M)>0$ and $S_4(M)\le16\sigma^2(M)$ for all $M\ge40$.

\emph{Step 2: $\sigma^2(M)\ge\sigma_0^2$ for $M\ge40$.} Each term
$8M^2/(M^2+\gamma^2)$ of $\sigma^2(M)$ is increasing in $M$, so
$\sigma^2(M)\ge\sigma^2(40)=160\,h(40)$, and by Lemma~\ref{lem:S4}
\[
h(40)\ge\tfrac12\log\frac q\pi+\tfrac12\psi\Bigl(\frac{40.5+a}{2}\Bigr)-Z_1(40.5).
\]
Odd real primitive characters have $q\ge3$ and even ones $q\ge5$.
Let $\sigma_0^2$ be the minimum, over the two pairs
$(a,q_{\min})\in\{(1,3),(0,5)\}$, of
$160\,[\tfrac12\log(q_{\min}/\pi)+\tfrac12\psi((40.5+a)/2)-Z_1(40.5)]$.
Then $\sigma^2(M)\ge\sigma_0^2$ for every $\chi$ and every $M\ge40$,
and (E3) is the evaluation of $\sigma_0^2$ (the minimum is attained
for $(a,q)=(1,3)$).

\emph{Step 3: at least nine ordinates are $\le5M$, for $M\ge40$.} Let
$N_\chi(\Gamma)$ be the number of ordinates in $(0,\Gamma]$. We claim
that for all $Y,\Gamma>0$
\begin{equation}\label{eq:countlemma}
N_\chi(\Gamma)\;\ge\;\frac Y2\Bigl[h(Y)-\frac{2Y}{\Gamma}\,h(\Gamma)\Bigr].
\end{equation}
Indeed, in $\sigma^2(Y)=\sum_\gamma8Y^2/(Y^2+\gamma^2)$ each ordinate
$\gamma\le\Gamma$ contributes at most $8$, and for $\gamma>\Gamma$
\[
\frac{8Y^2}{Y^2+\gamma^2}\le\frac{2Y^2}{\Gamma^2}\cdot
\frac{8\Gamma^2}{\Gamma^2+\gamma^2},
\]
because this is equivalent to $\Gamma^2+\gamma^2\le2Y^2+2\gamma^2$.
Summing, $\sigma^2(Y)\le8N_\chi(\Gamma)+(2Y^2/\Gamma^2)\sigma^2(\Gamma)$,
and \eqref{eq:countlemma} follows from $\sigma^2(Y)=4Yh(Y)$,
$\sigma^2(\Gamma)=4\Gamma h(\Gamma)$. Take $Y=40$, $\Gamma=200$,
$w:=2Y/\Gamma=\tfrac25$. By Lemma~\ref{lem:S4},
\[
h(Y)-w\,h(\Gamma)\ge(1-w)\,\tfrac12\log\frac q\pi
+\tfrac12\psi\Bigl(\frac{Y+\frac12+a}{2}\Bigr)
-\frac w2\,\psi\Bigl(\frac{\Gamma+\frac12+a}{2}\Bigr)
-Z_1(Y+\tfrac12)-w\,Z_1(\Gamma+\tfrac12),
\]
and since $1-w\ge0$ the right-hand side is smallest at $q=q_{\min}$.
(E4) states that for both pairs $(a,q_{\min})$ the resulting lower
bound for $\frac Y2[\,\cdots]$ is at least $11.3751$. Hence
$N_\chi(200)\ge12\ge9$ for every $\chi$, and for $M\ge40$ the nine
lowest ordinates satisfy $\gamma_i\le200\le5M$, i.e.
\begin{equation}\label{eq:ninelow}
c_i(M)\ge c_\kappa=\frac4{\sqrt{26}}\qquad(i=1,\dots,9).
\end{equation}

\emph{Step 4: envelopes.} Let $\mathcal L:=\{1,\dots,9\}$ and
$M\ge40$.

For $0\le t\le\tau$: $c_jt\le4\tau=z_*<j_{0,1}$ for all $j$, and
$0<J_0(c_jt)\le e^{-c_j^2t^2/4}$ by (B3). Hence
\[
|\Phi(t)|\le e^{-\sigma^2t^2/2},\qquad
|\Phi^{(j)}(t)|\le e^{-(\sigma^2-c_j^2/2)t^2/2}\le e^{-(\sigma^2-8)t^2/2}
\quad(\text{every }j).
\]
For $t\ge\tau$ and $j\notin\mathcal L$:
$|J_0(c_jt)|\le b(c_jt)\le b(c_j\tau)=e^{-c_j^2\tau^2/4}$, by
Lemma~\ref{lem:J}, because $b$ is nonincreasing and $c_j\tau\le z_*$.
For $i\in\mathcal L$: $|J_0(c_it)|\le p(t):=\min\bigl(1,\frac8\pi(c_\kappa t)^{-1/2}\bigr)$
by (B1), (B2) and \eqref{eq:ninelow}; thus $p(t)=1$ for $t\le T_c$ and
$p(t)=(T_c/t)^{1/2}$ for $t\ge T_c$. The terms outside $\mathcal L$
carry the variance $\sigma^2-\sum_{i\in\mathcal L}c_i^2/2\ge\sigma^2-72$.
Therefore, for $t\ge\tau$,
\[
|\Phi(t)|\le\rho^{(\sigma^2-72)/8}\,p(t)^9,\qquad
|\Phi^{(j)}(t)|\le\rho^{(\sigma^2-80)/8}\,p(t)^8\quad(\text{every }j):
\]
for $j\in\mathcal L$ eight factors $p$ remain and the other terms give
$\rho^{(\sigma^2-72)/8}\le\rho^{(\sigma^2-80)/8}$; for
$j\notin\mathcal L$ nine factors $p\le1$ remain and the terms outside
$\mathcal L\cup\{j\}$ have variance at least $\sigma^2-72-8$. Define
\[
W(t):=\rho^{(\sigma^2-72)/8}p(t)^9\ \ (t\ge\tau),\qquad
W^1(t):=\begin{cases}e^{-(\sigma^2-8)t^2/2}&(t\le\tau),\\
\rho^{(\sigma^2-80)/8}p(t)^8&(t>\tau).\end{cases}
\]

\emph{Step 5: Proposition~\ref{prop:C0}.} Apply it with $z_1=z_*$
($<2.4<j_{0,1}$), $t_1=\tau$ ($\le\pi/2$, and $c_1t_1\le4\tau=z_1$),
$\overline{\sigma}^2=\sigma^2$, $\overline S_4=16\sigma^2$ (Step~1),
and $W$, $W^1$ as above (the sequence has at least nine terms by
Step~3):
\begin{equation}\label{eq:Ethree}
-M\hd'(M)\ge\frac1\pi\int_0^\tau e^{-\sigma^2t^2/2-16\beta_*\sigma^2t^4}\cos t\,dt
-\frac1\pi\int_\tau^\infty W\,dt-\frac{\sigma^2}{2\pi}\int_0^\infty t^2W^1\,dt .
\end{equation}
\emph{First integral.} For $x\ge0$ and $0\le t\le\tau$,
$e^{-x}\cos t\ge1-x-t^2/2$ (if $1-x\ge0$ multiply the bounds
$e^{-x}\ge1-x\ge0$ and $\cos t\ge1-t^2/2>0$; otherwise the left side
is $\ge0\ge1-x$). With $x=16\beta_*\sigma^2t^4$, the Gaussian moments
$\int_0^\infty t^2e^{-\sigma^2t^2/2}dt=\sqrt{\pi/2}\,\sigma^{-3}$,
$\int_0^\infty t^4e^{-\sigma^2t^2/2}dt=3\sqrt{\pi/2}\,\sigma^{-5}$, and
$\int_\tau^\infty e^{-\sigma^2t^2/2}dt\le e^{-\sigma^2\tau^2/2}/(\sigma^2\tau)
=\rho^{\sigma^2/8}/(\sigma^2\tau)$,
\[
\frac1\pi\int_0^\tau e^{-\sigma^2t^2/2-16\beta_*\sigma^2t^4}\cos t\,dt
\ge\frac{1}{\sqrt{2\pi}\,\sigma}\Bigl[1-\frac1{2\sigma^2}-\frac{48\beta_*}{\sigma^2}\Bigr]
-\frac{\rho^{\sigma^2/8}}{\pi\sigma^2\tau}.
\]
\emph{Second integral.} $\int_\tau^\infty p^9\le T_c+\int_{T_c}^\infty(T_c/t)^{9/2}dt
=\frac97T_c$, so
$\frac1\pi\int_\tau^\infty W\le\frac{9}{7\pi}T_c\,\rho^{(\sigma^2-72)/8}$.

\noindent\emph{Third integral.}
$\int_0^\tau t^2e^{-(\sigma^2-8)t^2/2}dt\le\sqrt{\pi/2}\,(\sigma^2-8)^{-3/2}$
and $\int_\tau^\infty t^2p^8\le\frac13T_c^3+\int_{T_c}^\infty T_c^4t^{-2}dt
=\frac43T_c^3$, so
\[
\frac{\sigma^2}{2\pi}\int_0^\infty t^2W^1\,dt\le
\frac{1}{\sqrt{2\pi}\,\sigma}\cdot\frac12\Bigl(1-\frac8{\sigma^2}\Bigr)^{-3/2}
+\frac{\sigma^2}{2\pi}\cdot\frac43\,T_c^3\,\rho^{(\sigma^2-80)/8}.
\]
Multiplying \eqref{eq:Ethree} by $\sqrt{2\pi}\,\sigma$ and collecting
the terms gives $\sqrt{2\pi}\,\sigma\cdot(-M\hd'(M))\ge G(\sigma^2)$,
with $G$ as in \eqref{eq:defG}.

\emph{Step 6: $G(\sigma^2)\ge G(\sigma_0^2)>0$.} For $x\ge15$ every
subtracted term of $G$ is decreasing in $x$: $1/x$;
$(1-8/x)^{-3/2}$; $x^{-1/2}\rho^{x/8}$; $x^{1/2}\rho^{x/8}$ (for
$x>4/\log(1/\rho)=4.954\ldots$); and $x^{3/2}\rho^{x/8}$ (for
$x>12/\log(1/\rho)=14.862\ldots$). So $G$ is increasing on $[15,\infty)$, and
by Step~2, $G(\sigma^2(M))\ge G(\sigma_0^2)$; by (E5) this is
$\ge0.3219715>0.3219$.
\end{proof}

\begin{remark}\label{rem:Enotes}
(1) \emph{Multiplicities.} The theorem concerns the model of
Section~\ref{subsec:setting}, in which a zero of multiplicity $k$
contributes $k$ independent amplitudes $c\le4$; the step
$c_j\tau\le z_*$ would fail if such a zero were modelled as one
amplitude $kc$. Under LI all zeros are simple and the two readings
coincide.

(2) \emph{What the bound is worth.} For $q=3,4$ and
$M\in\{40,50,100,400,2000\}$ the bound of Theorem~\ref{thm:E} is
between $0.50$ and $0.88$ of the certified lower bound of
Section~\ref{sec:main} for the same $M$, and never above it (for
instance $q=4$, $M=50$: $0.010220$ against $0.013341$). The constant
$0.3219$ is not optimised; the same proof with other auxiliary
parameters gives $G\ge0.475$ for $M\ge50$, and on the grid of
parameters that was searched the smallest threshold reached is
$M\ge38.5$. The theorem does not reproduce the two-sided asymptotic
for $\hd_\chi'$ of \cite[Theorem~4]{WR}.

(3) \emph{Relation to \cite{WR}.} The argument is independent of the
large-weight results of the companion paper: it uses none of its zero
counts or tail bounds, only the variance identity and the facts (B1)--(B4).
\end{remark}

\begin{theorem}[General Theorem, precise form]\label{thm:general}
Let $\chi$ be a real primitive Dirichlet character and assume
GRH$(\chi)$ and $L(\tfrac12,\chi)\ne0$. Then:
\begin{enumerate}
\item[(i)] $\hd_\chi\in C^1(0,\infty)$, and $\hd_\chi'$ is given by
\eqref{eq:deriv} and \eqref{eq:Did};
\item[(ii)] $\hd_\chi$ is strictly decreasing on $(0,M_0(\chi)]$;
\item[(iii)] for every $M\ge40$,
$-M\hd_\chi'(M)\ge G(\sigma_M^2)/(\sqrt{2\pi}\,\sigma_M)
\ge0.3219/(\sqrt{2\pi}\,\sigma_M)>0$; in particular $\hd_\chi$ is
strictly decreasing on $[40,\infty)$;
\item[(iv)] $Q_\chi(40)>38$, hence $M_0(\chi)<40$.
\end{enumerate}
Under LI$(\chi)$ in addition, $\hd_\chi(M)$ is the logarithmic density
of $\{x\ge2:D_{m,b}(x)>0\}$, $m=M-\tfrac12$, for every real $b$.
\end{theorem}

\begin{proof}
(i) is Lemma~\ref{lem:diff}, (ii) is Theorem~\ref{thm:B}(ii), (iii) is
Theorem~\ref{thm:E}, and the last sentence is
Theorem~\ref{thm:B}(iii). For (iv): since $c_j<4$,
$Q_\chi(M)=\sum_jc_j^3/16\ge\sum_jc_j^4/64=S_4(M)/64$. At $M=40$,
$s=40.5$, Lemma~\ref{lem:S4} gives
$S_4=16\sigma^2-64\cdot1600\cdot h'(40)$, and
$h'(40)\le\frac14\psi'(20.25)+Z_2(40.5)\le0.01296+10^{-10}$, using
$\psi'(x)=\sum_{k\ge0}(x+k)^{-2}\le x^{-2}+x^{-1}$ at $x=20.25$ (the
smaller of the two possible arguments $(s+a)/2$) and
$Z_2(40.5)<2^{-39}$. With $\sigma^2(40)\ge236.97$ (Step~2 of the proof
of Theorem~\ref{thm:E}),
$S_4(40)\ge16\cdot236.97-102400\cdot0.01297>2463$, so
$Q_\chi(40)>38$. As $Q_\chi$ is increasing and $Q_\chi(M_0)=1$, this
gives $M_0(\chi)<40$.
\end{proof}

Thus for every real primitive $\chi$ (under GRH$(\chi)$ and
NCZ$(\chi)$) the half-line splits into a range $(0,M_0(\chi)]$ of
monotonicity near the critical weight, a range $[40,\infty)$ of
monotonicity for large weights, and a nonempty intermediate interval
$(M_0(\chi),40)$ on which Theorem~\ref{thm:general} asserts nothing.
Sections~\ref{sec:main} and~\ref{sec:nonmono} show that both
behaviours occur there.

\section{Full monotonicity for \texorpdfstring{$q=3,4$}{q=3,4}: a
computer-assisted theorem}\label{sec:main}

Throughout this section $q\in\{3,4\}$ and $\chi=\chi_q$; both
characters are odd, NCZ$(\chi)$ holds unconditionally, and
GRH$(\chi)$ is assumed. The sequence of
Section~\ref{subsec:envsetting} is the sequence of positive ordinates
of $L(s,\chi_q)$, and $\hd_q$, $\sigma^2$, $S_4$, $h$ are as in
Sections~\ref{sec:envelope}--\ref{sec:largeweight}.

\subsection{The envelope certificate}\label{subsec:envcert}

The certificate (directory \texttt{monotonicity/envelope-cert/} of the
repository; program \texttt{certify\_envelope.py}, method note
\texttt{METHOD.md}) evaluates the hypothesis of
Corollary~\ref{cor:cell}, case~(b), with the $n=12$ lowest ordinates,
in Arb ball arithmetic at $128$ bits. For a cell $[M_a,M_b]$ with
exact rational endpoints it computes a number
$\mathrm{LB}$ and reports the cell as certified only if it has verified
that
\[
\mathrm{MAIN}_{\mathrm{cell}}-\mathrm{TAILF}_{\mathrm{cell}}
-\mathrm{NUM}_{\mathrm{cell}}\;\ge\;\mathrm{LB}\;>\;0 ,
\]
the last comparison being made on the lower endpoint of a ball. By
Corollary~\ref{cor:cell} this gives $-M\hd_q'(M)\ge\mathrm{LB}$, hence
$\hd_q'(M)\le-\mathrm{LB}/M_b<0$, for every real $M$ in the closed
cell. In outline:
\begin{itemize}
\item \emph{Ordinates.} The twelve lowest ordinates enter as balls
$d\pm10^{-20}$, $d$ the listed decimal. Completeness of the list is
used to identify $v_{12}(M_a)$ with the contribution of all other
zeros.
\item \emph{Closed forms.} $\sigma^2=4Mh$ and $S_4=64Mh-64M^2h'$
(Lemma~\ref{lem:S4}) are evaluated at the exact endpoints. For
$s=M+\tfrac12\ge4$ this is done from the Dirichlet series of $L'/L$ and
$(L'/L)'$ over prime powers up to $20000$ with integral-test bounds
for the remainder, and $\psi$, $\psi'$ in Arb. For the $127$ cells
with $M_a<3.5$, $h$ is evaluated by a complex-step formula with a
Cauchy remainder, and $S_4$ is enclosed without $h'$ from the full
list of ordinates:
$\sum_{i\le N}c_i^4\le S_4\le\sum_{i\le N}c_i^4+c_N^2\cdot2V$, where
$V=\sigma^2-\sum_{i\le N}c_i^2/2$ is the variance of the unlisted
zeros; this step uses that the list is complete and the listed zeros
simple.
\item \emph{Integrals.} $\mathrm{MAIN}$ is a lower Riemann sum of a
positive decreasing integrand; $\mathrm{TAILF}$ and $\mathrm{NUM}$ are
upper Riemann sums (the envelopes $B$, $B^1$ are nonincreasing) on an
exact dyadic grid, plus a closed-form bound for the part beyond the
last grid point, where nine factors are in the range
$z\ge z_{\mathrm{far}}$. For $\beta(z_1)$ the function $J_0$ is
evaluated at one of $20$ exact rational points $z_1<2.4$. Floating
point is used only to place grid points and to print; every quantity
chosen with its help is an exact rational that is re-verified in ball
arithmetic.
\end{itemize}

\noindent\textbf{Recorded result} (from the summary blocks of the
files \texttt{results\_q4.txt} and \texttt{results\_q3.txt} of that
directory, which also record the SHA-256 digests of the program, of
the shared module and of the zero file used):
\[
\begin{array}{c|c|c|c|c|c}
q & \text{certified range of }M & \text{cells} & \text{failed} &
\text{smallest }\mathrm{LB} & \text{weakest bound on }\hd_q'\\\hline
4 & [1,\,2000] & 765 & 0 & 1.318804\times10^{-3} & \hd_4'\le-6.616829\times10^{-7}\\
3 & [1,\,2000] & 765 & 0 & 1.353296\times10^{-3} & \hd_3'\le-6.789885\times10^{-7}
\end{array}
\]
The cells are contiguous, consecutive endpoints having ratio $1.01$,
and tile $[1,2000]$ exactly. On every cell
$\mathrm{MAIN}>\mathrm{TAILF}$ and
$\mathrm{NUM}/(\mathrm{MAIN}-\mathrm{TAILF})\le0.44115$ ($q=4$), resp.\
$\le0.43819$ ($q=3$), the largest values occurring at the top of the
range; so the criterion never has less than a factor $2.26$ to spare.
The first cell $[1,1.01]$ carries the certified upper bounds
\begin{equation}\label{eq:sig101}
\sigma^2(1.01)\le0.6292\quad(q=4),\qquad
\sigma^2(1.01)\le0.4599\quad(q=3).
\end{equation}
We re-parsed the two result files with the script of
Section~\ref{sec:numerics}(5): $765$ certified cells each, contiguous,
covering $[1,2000]$, with the values of the table.

\begin{proposition}[content of the envelope certificate]
\label{prop:envcert}
Let $q\in\{3,4\}$ and assume GRH$(\chi_q)$. Granted the inputs
\textup{(T1)--(T3)} of Section~\ref{subsec:inputs},
$\hd_q'(M)<0$ for every $M\in[1,2000]$, with the bounds of the table
above, and \eqref{eq:sig101} holds.
\end{proposition}

\begin{proof}
Corollary~\ref{cor:cell}, case (b), $n=12$, on each cell; its
hypothesis is what the certificate evaluates. The identification of
the computed numbers with $\sigma^2$, $S_4$, $v_{12}$ uses
Lemma~\ref{lem:S4} (i.e.\ (T3)) and the completeness and simplicity of
the lists (T1).
\end{proof}

\subsection{The trusted inputs}\label{subsec:inputs}

\begin{itemize}
\item[\textup{(T1)}] \emph{The zero files.} The two files
\texttt{chi4\_zeros.txt} and \texttt{chi3\_zeros.txt} of the
directory \texttt{prototype/explore/} ($511$ ordinates of
$L(s,\chi_4)$ and $537$ ordinates of $L(s,\chi_3)$), with SHA-256
digests
\begin{center}
\small
\begin{tabular}{ll}
$\chi_4$: & \texttt{0e73e00a6c1609bdb4eec62c13292435}\\
          & \texttt{d0beee55750735a51f0fab84024ddfbc}\\
$\chi_3$: & \texttt{e3075235a847f43eb35057d3998a853f}\\
          & \texttt{9383e27e958cc5fd3c4115dad7eb8aaa}
\end{tabular}
\end{center}
together with the certificate, produced by the separate program
\texttt{certify\_arb.py} of \cite[\S7]{WR} in ball arithmetic,
that Hardy's function changes sign in each window $d\pm10^{-20}$ and
that the winding number of $L(s,\chi)$ around a rectangle reaching
height $640.165707$ (resp.\ $700.024747$) equals the number of
windows. This certificate implies that every listed value is within
$10^{-20}$ of an ordinate, that the listed zeros are simple, that
there is no other zero with ordinate in $(0,\gamma_N]$, and that
$L(\tfrac12,\chi)\ne0$. It was re-run by two internal auditors with
the same outcome. The envelope certificate uses of it: the twelve
lowest ordinates, the absence of other zeros below $\gamma_{12}$, and,
for the $127$ cells with $M_a<3.5$, the whole list.
\item[\textup{(T2)}] \emph{Software.} The correctness of FLINT/Arb
(FLINT~3.6.0 through python-flint~0.9.0: ball arithmetic, elementary
functions, $\psi$, the Hurwitz zeta function, $J_0$ at exact points,
and Dirichlet $L$-functions at complex arguments for $M<3.5$), and of
the certifying programs \texttt{certify\_envelope.py} (with the module
it shares with the direct certificate of
Section~\ref{subsec:confirm}) and \texttt{analytic\_leg.py}.
\item[\textup{(T3)}] \emph{The variance identity} of the companion
paper, $\sigma_M^2=4M\frac{\Lambda'}{\Lambda}(M+\tfrac12,\chi)$
\cite[eq.~(8)]{WR} (Proposition~\ref{prop:WRvar}); its $M$-derivative
is Lemma~\ref{lem:S4}.
\end{itemize}
The only unproved mathematical hypothesis is GRH$(\chi_q)$. Besides
(T1)--(T3) the argument uses the proofs of this article and the facts
(B1)--(B4) of Section~\ref{subsec:bessel} and
\cite[Lemma~7(b)]{WR}, which are ordinary mathematics and can be
refereed as such.

\subsection{The theorem}

\begin{theorem}[computer-assisted]\label{thm:q34}
Let $q\in\{3,4\}$ and assume GRH$(\chi_q)$. Granted the inputs
\textup{(T1)--(T3)} of Section~\ref{subsec:inputs}, the function
$\hd_q$ is strictly decreasing on $(0,\infty)$; equivalently
$m\mapsto\delta_q(m)$ is strictly decreasing on $(-\tfrac12,\infty)$,
which is \cite[Conjecture~1]{WR} for this $q$. More precisely,
$\hd_q$ is strictly decreasing on $(0,M_0(\chi_q)]$, where
$M_0(\chi_q)>1.01$, and $\hd_q'(M)<0$ for every $M\ge1$. Under
LI$(\chi_q)$ in addition, the logarithmic density of
$\{x\ge2:D_{m,b}(x)>0\}$ exists for every $m>-\tfrac12$ and every real
$b$ and is a strictly decreasing function of $m$.
\end{theorem}

\begin{proof}
By Theorem~\ref{thm:B}(ii), $\hd_q$ is strictly decreasing on
$(0,M_0(\chi_q)]$. By Proposition~\ref{prop:M0}(c) and
\eqref{eq:sig101}, $Q_{\chi_q}(1.01)<\sigma^2(1.01)/2<1$, so
$M_0(\chi_q)>1.01$. By Proposition~\ref{prop:envcert},
$\hd_q'<0$ on $[1,2000]$, and by Theorem~\ref{thm:E}, $\hd_q'<0$ on
$[40,\infty)$; the two ranges overlap. Thus $\hd_q$ is strictly
decreasing on $(0,1.01]$ and on $[1,\infty)$, hence on $(0,\infty)$.
The last assertion is Theorem~\ref{thm:C}(iii).
\end{proof}

Theorem~\ref{thm:E} is needed here only for $M>2000$. Neither
\cite[Theorem~4]{WR} nor \cite[Corollary~7]{WR} (with its
Proposition~5 and Lemma~24) is used.

\subsection{Independent confirmations}\label{subsec:confirm}

The following are not links of the chain above. They confirm its
conclusion on parts of the range by other means.

\subsubsection*{(a) The direct certificate}\label{subsec:cert}
A second certificate (the program \texttt{certify\_derivative.py}
in the directory \texttt{monotonicity/cert/}) encloses the
oscillatory integral
$\mathcal D_\chi(M)$ of \eqref{eq:deriv} itself, simultaneously for
all $M$ in a cell: interval Riemann sums of the integrand on
$[0,T_1]$, in which all listed zeros ($511$ resp.\ $537$) enter
through evaluated Bessel functions and the unlisted ones through
exact tail sums obtained from the variance identity and its
derivative, and an explicit bound for $t\ge T_1$. It certifies
two-sided bounds $\mathrm{lo}\le\hd_q'(M)\le\mathrm{hi}<0$ on each
cell:
\[
\begin{array}{c|l|c|c|c|l}
q & \text{file} & \text{range of }M & \text{cells} &
\text{largest }\mathrm{hi} & \text{status}\\\hline
4 & \texttt{results\_q4.txt} & [3.6,\,17.4] & 317 & -1.675\times10^{-3}
& \text{audited, re-proved}\\
4 & \texttt{results\_q4\_stretch.txt} & [17.4,\,320] & 293 & -1.224\times10^{-5}
& \text{not audited}\\
3 & \texttt{results\_q3.txt} & [4.4,\,19.2] & 296 & -1.549\times10^{-3}
& \text{audited, re-proved}\\
3 & \texttt{results\_q3\_stretch.txt} & [19.2,\,350] & 292 & -1.098\times10^{-5}
& \text{not audited}
\end{array}
\]
Its auxiliary file reports, under the same hypotheses,
\begin{equation}\label{eq:M0cert}
3.688854<M_0(\chi_4)<3.689059,\qquad
4.483812<M_0(\chi_3)<4.484161 ,
\end{equation}
(these lie inside the cruder floating-point brackets of
Remark~\ref{rem:M0values}), as well as enclosures of the root of
$\sigma_M^2=2$ and of $\lambda_\chi(M,M')$ for the pairs quoted in
Section~\ref{sec:numerics}.

\emph{Outcome of the independent audit} (report
\texttt{monotonicity/cert-audit/AUDIT.md}; internal). No break of
rigor was found: every inequality of the method was re-derived, is
used in the correct direction and inside its range of validity, and is
implemented as stated. Three gaps were identified. (G1) The identity
between the enclosed integral and $\hd_q'$ was not covered by
\cite[Lemma~21]{WR} on most of the range; it is Lemma~\ref{lem:diff}
of this article, which holds at every $M>0$, endpoints of closed cells
included. (G2) The certificate is tied to its zero files only by
convention, and the program accepts a wrong list silently (with the
lowest ordinate deleted it still reports a cell as certified, with a
false enclosure); this is why the digests of the files are part of
(T1). (G3) The result files had been produced by an earlier version of
the program; the audit recomputed the two audited ranges from scratch
with the final version and reproduced all $613$ lines field for field.
Eleven minor items are listed in the report. Beyond the audit of the
code (including mutation tests), the statement was \emph{re-proved by a
second, independently written rigorous program} on every one of the
$317+296$ cells of $[3.6,17.4]$ and $[4.4,19.2]$: every independent
enclosure is negative and lies strictly inside the audited one, giving
$\hd_4'\le-1.706259\times10^{-3}$ on $[3.6,17.4]$ and
$\hd_3'\le-1.574568\times10^{-3}$ on $[4.4,19.2]$. The values
\eqref{eq:M0cert} were re-certified independently as well. The
stretch files up to $320$ and $350$ were produced by the same program
and are pinned by their digests in the audit's list, but they were not
audited; we re-parsed them (contiguous, all cells certified) and
nothing more.

\subsubsection*{(b) Agreement of the two certificates}
On the common range, whenever a cell of the envelope certificate and a
cell of the direct certificate overlap in $[x,y]$, both statements
hold there, so necessarily $\mathrm{LB}\le x\,|\mathrm{lo}|$. This
holds for all $1061$ ($q=4$) and $1028$ ($q=3$) overlapping pairs, the
largest value of $\mathrm{LB}/(x|\mathrm{lo}|)$ being $0.945$. (This
comparison uses printed decimals and is a consistency check, not a
proof; as an internal referee observed, it would not detect an error
in $\mathrm{TAILF}$, which was tested separately by an independent
evaluation of the cell quantity.) The two certificates share the
hypotheses (T1)--(T3), the zero files and the evaluation of $h$ and
$h'$; they share no integrand, no quadrature, and no Bessel evaluation
on wide balls.

\subsubsection*{(c) The analytic large-weight lemma of the companion paper}
For $M\ge300$, \cite[Lemma~24]{WR} asserts $\hd_q'(M)<0$ with no zero
data. In \cite{WR} the extension from the single point $M=300$ to all
$M\ge300$ was described as checked in floating point only.
\emph{Outcome of the audit of that leg} (report
\texttt{monotonicity/ANALYTIC-LEG-AUDIT.md}; internal): the chain of
lemmas of \cite[\S6]{WR} leading to Lemma~24 was re-derived from first
principles and the following was established.
\begin{quote}
Assume GRH$(\chi_q)$, $q\in\{3,4\}$. Then $\hd_q$ is differentiable on
$(299,\infty)$ and $\hd_q'(M)<0$ for all $M\ge300$.
\end{quote}
No fatal finding. The extension step is a correct hand proof once
three monotonicity facts, left implicit in \cite{WR}, are stated; the
report writes it out in full and certifies every numerical fact it
uses in ball arithmetic ($90$ checks, all passing), and confirms it
independently on $[300,10^7]$ by a cell sweep. The report lists three
repairable items in the text of \cite{WR}, none of which affects the
statement above; they are recorded there for a revision of \cite{WR}.
That leg uses neither zero lists nor LI.

\medskip
In summary: on $[3.6,17.4]$ resp.\ $[4.4,19.2]$ the conclusion
$\hd_q'<0$ is established three times (envelope certificate; direct
certificate; independent re-proof of the latter), on the rest of
$[3.6,320]$ resp.\ $[4.4,350]$ twice (the second time unaudited), and
for $M\ge300$ three times (envelope certificate up to $2000$;
Theorem~\ref{thm:E}; the audited lemma of \cite{WR}). On
$1\le M<3.6$ (resp.\ $<4.4$) it rests on the envelope certificate
alone, and there Theorem~\ref{thm:B} applies independently up to
$M_0(\chi_q)$, so that the only part of the half-line resting on a
single argument is $0<M<1$ (Theorem~\ref{thm:B}).

\begin{remark}[what is taken on trust]\label{rem:mainstatus}
Theorem~\ref{thm:q34} is a computer-assisted theorem, conditional on
GRH$(\chi_q)$. Its inputs (T1)--(T3) are not mathematical hypotheses
in the usual sense. (T1) and (T2) are statements about computations:
that a certified list of zeros is the one that was used, and that a
library and three programs do what they are meant to do. The programs
were read line by line, re-run from scratch and subjected to mutation
tests by internal referees, and the envelope certificate was compared
with an independent evaluation at $13$ cells per modulus; none of
this is a formal verification. (T3) is a theorem of an unrefereed
companion preprint; it was re-derived in three of the internal
reports. All of
this refereeing was internal and AI-conducted
(Section~\ref{sec:gaps}).
\end{remark}

\section{Monotonicity is not a general fact}\label{sec:nonmono}

The monotonicity established for $q=3,4$ is not a consequence of the
shape of the random variables \eqref{eq:XM} alone: it depends on
where the zeros are. This section gives an elementary model in which
$\hd$ increases, and then a computer-assisted theorem showing that the
same happens for two actual characters.

\subsection{A two-zero model}

Here $\gamma_1<\gamma_2$ are any two positive numbers, not ordinates
of an $L$-function, and Lemma~\ref{lem:diff} (which needs nine terms)
is not used.

\begin{proposition}[a two-zero model in which the density increases]
\label{prop:twozero}
Let $0<\gamma_1<\gamma_2$, let $\theta_1,\theta_2$ be independent and
uniform on $[0,2\pi)$, and for $M>0$ put
$c_j=4M/\sqrt{M^2+\gamma_j^2}$,
$c_j'=4\gamma_j^2(M^2+\gamma_j^2)^{-3/2}$ and
\[
\delta^{(2)}(M):=\PP\bigl(1+c_1\cos\theta_1+c_2\cos\theta_2>0\bigr).
\]
If, at some $M$, $c_2<1$, $c_1>1+c_2$ and
\begin{equation}\label{eq:condN}
\frac{c_2\,c_2'}{2c_1^3}\;>\;
\frac{c_1'\,(1+c_2)}{c_1\sqrt{c_1^2-(1+c_2)^2}} ,
\end{equation}
then $\delta^{(2)}$ is differentiable at $M$ with
$\frac{d}{dM}\delta^{(2)}(M)>0$. For fixed $\gamma_2$ and fixed
$M<\gamma_2/\sqrt{15}$ these conditions hold for all sufficiently
small $\gamma_1$. For example, for
$(\gamma_1,\gamma_2,M)=(0.03,10,1)$ the two sides of \eqref{eq:condN}
are $1.23\times10^{-3}$ and $3.4\times10^{-4}$.
\end{proposition}

\begin{proof}
For $-1<\alpha<1$, $\PP(\cos\theta_1>-\alpha)=\tfrac12+\tfrac1\pi\arcsin\alpha$.
The conditions $c_2<1$, $c_1>1+c_2$ are open in $M$, and under them
$0<1-c_2\le1+c_2x\le1+c_2<c_1$ for all $x\in[-1,1]$. Conditioning on
$x=\cos\theta_2$,
\[
\delta^{(2)}=\frac12+\frac1\pi\,G_2(c_1,c_2),\qquad
G_2(c_1,c_2):=\EE\Bigl[\arcsin\frac{1+c_2x}{c_1}\Bigr],
\]
and $G_2$ is continuously differentiable on the open set
$\{c_2<1,\ c_1>1+c_2\}$, with derivatives obtained by differentiating
under the expectation (the integrands are bounded there). With
$g(u):=(c_1^2-u^2)^{-1/2}$ and the symmetry of $x$,
\[
\frac{\partial G_2}{\partial c_2}
=\EE\Bigl[\frac{x}{\sqrt{c_1^2-(1+c_2x)^2}}\Bigr]
=\frac12\,\EE\Bigl[|x|\bigl(g(1+c_2|x|)-g(1-c_2|x|)\bigr)\Bigr].
\]
Since $g'(u)=u(c_1^2-u^2)^{-3/2}\ge u\,c_1^{-3}$ on $[0,c_1)$, we have
$g(1+y)-g(1-y)\ge c_1^{-3}\int_{1-y}^{1+y}u\,du=2y\,c_1^{-3}$ for
$0\le y\le c_2$, and therefore
$\partial G_2/\partial c_2\ge c_2\,\EE[x^2]\,c_1^{-3}=c_2/(2c_1^3)$. On
the other hand, $u\mapsto u(c_1^2-u^2)^{-1/2}$ is increasing on
$[0,c_1)$, so
\[
\Bigl|\frac{\partial G_2}{\partial c_1}\Bigr|
=\EE\Bigl[\frac{1+c_2x}{c_1\sqrt{c_1^2-(1+c_2x)^2}}\Bigr]
\le\frac{1+c_2}{c_1\sqrt{c_1^2-(1+c_2)^2}} .
\]
By the chain rule
$\pi\,\frac{d}{dM}\delta^{(2)}=\frac{\partial G_2}{\partial c_1}c_1'
+\frac{\partial G_2}{\partial c_2}c_2'$, with $c_1',c_2'>0$, and
\eqref{eq:condN} makes the right-hand side positive. For the second
assertion: $c_2<1$ is equivalent to $M<\gamma_2/\sqrt{15}$, and as
$\gamma_1\to0$ with $M$ and $\gamma_2$ fixed, $c_1\to4>1+c_2$ and
$c_1'\to0$, so the right-hand side of \eqref{eq:condN} tends to $0$
while the left-hand side tends to $c_2c_2'/128>0$. The numerical
example is a direct evaluation.
\end{proof}

The mechanism is that a saturated arcsine summand (here $c_1\approx4$,
i.e.\ $M\gg\gamma_1$) has a density that is \emph{smallest} at the
centre of its support, so that adding independent symmetric noise
raises $\PP(|c_1\cos\theta_1+\,\cdot\,|<1)$ and hence $\delta^{(2)}$.
There is no conflict with Lemma~\ref{lem:det}, which holds for finite
sequences as well: in the example $Q(u)=1$ already at
$u\approx0.811\,\gamma_1=0.024$, far below $M=1$.

\subsection{Two real characters}

Let $\chi_{-163}=(\frac{-163}{\cdot})$ and
$\chi_{-232}=(\frac{-232}{\cdot})$ (Kronecker symbols): the real
primitive characters of conductor $163$ and $232$; both are odd. Their
first two ordinates are $\gamma_1=0.2029013\ldots$,
$\gamma_2=2.3685339\ldots$ for $\chi_{-163}$ and
$\gamma_1=0.3106761\ldots$, $\gamma_2=2.4922\ldots$ for $\chi_{-232}$,
so the first zero is very low compared with the second
($\gamma_1/\gamma_2\approx0.086$ and $0.125$, against $0.59$ and $0.71$
for $\chi_4$ and $\chi_3$).

\begin{theorem}[non-monotonicity; computer-assisted]\label{thm:nonmono}
\textup{(i)} Assume GRH$(\chi_{-163})$ and write
$\hd=\hd_{\chi_{-163}}$. Then
\[
\begin{array}{c|c|c}
M & m=M-\tfrac12 & \text{certified enclosure of }\hd(M)\\\hline
3/10 & -1/5 & [0.599269990566,\ 0.599270019978]\\
57/100 & 7/100 & [0.590250631768,\ 0.590250669413]\\
3/5 & 1/10 & [0.590328985850,\ 0.590329024804]\\
117/100 & 67/100 & [0.596052042574,\ 0.596052089770]\\
6/5 & 7/10 & [0.596034510065,\ 0.596034556476]\\
2 & 3/2 & [0.586642408419,\ 0.586642448867]
\end{array}
\]
Hence
$\hd(\tfrac3{10})>\hd(\tfrac{57}{100})<\hd(\tfrac35)<\hd(\tfrac{117}{100})
>\hd(\tfrac65)>\hd(2)$, with
\[
\hd(\tfrac65)-\hd(\tfrac35)\in[0.005705485,\,0.005705571],\qquad
\hd(\tfrac{117}{100})-\hd(\tfrac{57}{100})\in[0.005801373,\,0.005801459].
\]
\textup{(ii)} Assume GRH$(\chi_{-232})$ and write
$\hd=\hd_{\chi_{-232}}$. Then
\[
\begin{array}{c|c|c}
M & m=M-\tfrac12 & \text{certified enclosure of }\hd(M)\\\hline
2/5 & -1/10 & [0.606642556775,\ 0.606642613806]\\
71/100 & 21/100 & [0.597435124908,\ 0.597435193395]\\
18/25 & 11/50 & [0.597438040240,\ 0.597438108755]\\
1 & 1/2 & [0.598231098965,\ 0.598231174771]\\
101/100 & 51/100 & [0.598233617676,\ 0.598233693573]\\
3/2 & 1 & [0.594070789210,\ 0.594070852885]
\end{array}
\]
Hence
$\hd(\tfrac25)>\hd(\tfrac{71}{100})<\hd(\tfrac{18}{25})<\hd(1)
<\hd(\tfrac{101}{100})>\hd(\tfrac32)$, with
\[
\hd(1)-\hd(\tfrac{18}{25})\in[0.000792990,\,0.000793135],\qquad
\hd(\tfrac{101}{100})-\hd(\tfrac{71}{100})\in[0.000798424,\,0.000798569].
\]
\textup{(iii)} Under GRH,
\[
0.16441017<M_0(\chi_{-163})<0.16441035,\qquad
0.2511058<M_0(\chi_{-232})<0.2511076 .
\]
These statements hold provided the hypotheses and trusted inputs
\textup{(N1)--(N5)} listed below are correct.
\end{theorem}

\begin{corollary}[race form]\label{cor:nonmono}
Let $\chi\in\{\chi_{-163},\chi_{-232}\}$ and assume GRH$(\chi)$ and
LI$(\chi)$. Granted \textup{(N1)--(N5)}, the logarithmic density
$\delta_\chi(m)$ of $\{x\ge2:D_{m,b}(x)>0\}$, which exists for every
$m>-\tfrac12$ and every real $b$, is not a monotone function of $m$ on
$(-\tfrac12,\infty)$. For $\chi_{-163}$:
$\delta_\chi(\tfrac1{10})<\delta_\chi(\tfrac7{10})$, while
$\delta_\chi(-\tfrac15)>\delta_\chi(\tfrac1{10})$ and
$\delta_\chi(\tfrac7{10})>\delta_\chi(\tfrac32)$. For $\chi_{-232}$:
$\delta_\chi(\tfrac{11}{50})<\delta_\chi(\tfrac12)$, while
$\delta_\chi(-\tfrac1{10})>\delta_\chi(\tfrac{21}{100})$ and
$\delta_\chi(\tfrac{51}{100})>\delta_\chi(1)$.
\end{corollary}

\begin{proof}
$L(\tfrac12,\chi)\ne0$ is part of what the certificate establishes
(N3). By Theorems~\ref{thm:WRinterp} and~\ref{thm:C}(iii),
$\delta_\chi(m)=\hd_\chi(m+\tfrac12)$ for every real $b$, and the
inequalities are those of Theorem~\ref{thm:nonmono}.
\end{proof}

\noindent\textbf{Hypotheses and trusted inputs} (as stated in the note
\texttt{NONMONOTONE.md} of the directory
\texttt{monotonicity/nonmonotone/}, which contains the programs and the
file \texttt{results.txt} from which the tables are copied, and in the
audit report quoted below).
\begin{itemize}
\item[\textup{(N1)}] \emph{Assumed:} GRH$(\chi)$. It makes the
ordinates real and gives the variance identity. LI$(\chi)$ is assumed
in Corollary~\ref{cor:nonmono} only.
\item[\textup{(N2)}] \emph{Certified by the certificate itself,
unconditionally:} the listed zeros ($133$ ordinates up to height
$7617/64=119.015625$ for $\chi_{-163}$, $114$ up to
$6351/64=99.234375$ for $\chi_{-232}$) are all the zeros of
$\Lambda(s,\chi)$ with $|\im s|$ below that height, each simple and
within $10^{-20}$ of the listed value. (Sign changes of Hardy's
function in each window, and the argument principle on a rectangle,
the count being exactly twice the number of windows.)
\item[\textup{(N3)}] \emph{Certified:} $L(\tfrac12,\chi)\ne0$
($L(\tfrac12,\chi_{-163})=0.0685317\ldots$,
$L(\tfrac12,\chi_{-232})=0.2462991\ldots$; it also follows from the
count in (N2)).
\item[\textup{(N4)}] \emph{Cited, not re-proved there:} the classical
facts on the completed $L$-function (entire of order one, root number
$+1$, zeros in the critical strip); $|J_0(x)|\le x^{-1/2}$
\cite[\S13.74]{Wat44}; $0<J_0(x)\le e^{-x^2/4}$ and
$-\log J_0(x)-x^2/4\le0.018\,x^4$ on $[0,1]$; the remainder of
Gauss--Legendre quadrature; the inversion formula
\cite[Lemma~7(b)]{WR}. \emph{Imported from \cite{WR}:} the variance
identity (Proposition~\ref{prop:WRvar}) and, for the corollary, the
limiting-distribution theorem (Theorem~\ref{thm:WRinterp}).
\item[\textup{(N5)}] \emph{Trusted:} FLINT/Arb~3.6.0 through
python-flint~0.9.0 (ball arithmetic, Dirichlet $L$-functions, Hardy's
function, the power-series jet of $L$, $\log\Gamma$, $\psi$,
$J_0$) and the certifying programs; and that the zero files used
are the certified ones, namely those with SHA-256 digests
\begin{center}
\small
\begin{tabular}{ll}
\texttt{zeros\_q163.txt}: & \texttt{b9d38aba6df821dea401035ea6ae24aa}\\
                          & \texttt{045305e757f81ea1a6233b632ad25b9d}\\
\texttt{zeros\_q232.txt}: & \texttt{42f40b29d79cc5228478acfd0480f239}\\
                          & \texttt{39c6f3b5d0479acaf178c661fa0af378}
\end{tabular}
\end{center}
\end{itemize}

\noindent\textbf{Method, in brief.} For a rational $M$ the density is
the inversion integral
$\hd(M)-\tfrac12=\frac1\pi\int_0^\infty\Phi_M(t)\frac{\sin t}{t}\,dt$.
The listed zeros enter through evaluated factors $J_0(c_jt)$; the
unlisted ones through the exact variance
$V=\sigma^2(M)-\sum_{\mathrm{listed}}c_j^2/2$ and the two-sided bound
$e^{-Vt^2/2-c_bt^4}\le\prod_{\mathrm{unlisted}}J_0(c_\gamma t)\le e^{-Vt^2/2}$
on the range where all unlisted arguments are at most $1$, with no
statistical assumption on the unlisted zeros. The integral over
$[0,T]$ is computed by composite Gauss--Legendre quadrature in
$192$-bit balls with a rigorous remainder, the model error by upper
Riemann sums, and the truncation error from
$|J_0(x)|\le x^{-1/2}$. The width of each enclosure (at most
$7.6\times10^{-8}$) is essentially the model error, i.e.\ the price of
knowing the zeros only up to height about $100$; the gaps
$5.7\times10^{-3}$ (conductor $163$) and $7.9\times10^{-4}$
(conductor $232$) exceed it by a factor of more than $10^4$. The
final comparisons are made on the printed, outward-rounded endpoints
in exact rational arithmetic.

\medskip
\noindent\textbf{Outcome of the independent audit of this
certificate.} \NonmonoAuditOutcome

\begin{remark}[what the theorem shows, and what it does not]
\label{rem:nonmono}
(1) The intermediate range of Theorem~\ref{thm:general} genuinely
matters: for $\chi_{-163}$ the density is larger at $M=1.17$ than at
$M=0.57$, and for $\chi_{-232}$ larger at $M=1.01$ than at $M=0.71$,
and these points lie well inside $(M_0(\chi),40)$. So Theorem~\ref{thm:q34} is not an instance
of a general fact about real primitive characters; the proof for
$q=3,4$ uses quantitative information on the lowest ordinates (twelve
of them on most of the range), and some such information is unavoidable.

(2) Theorem~\ref{thm:B} is not contradicted: by
Theorem~\ref{thm:nonmono}(iii) both increases lie above $M_0(\chi)$.
But its range cannot be extended, uniformly over real primitive
characters, by more than a bounded factor: for $\chi_{-163}$,
$\hd(0.57)<\hd(0.6)$ although $0.6\approx3.65\,M_0(\chi_{-163})$.
Theorem~\ref{thm:E} is not contradicted either; for $M\ge40$ the
character $\chi_{-163}$ satisfies its conclusion like every other.

(3) Nothing is claimed about the location of the local extrema, about
the exact intervals on which $\hd'>0$ (in floating point, about
$[0.57,1.17]$ and $[0.71,1.01]$), or about the sign of $\hd'$ at any
point: the theorem asserts only inequalities between values of $\hd$
at the listed points. Nothing is claimed for other characters. An
exploratory floating-point survey of the $224$ real primitive
characters of conductor at most $420$, reported in the repository
notes, found these two cases and one too marginal to call
(conductor $403$); it is not part of this article.
\end{remark}

\section{Numerical sanity checks}\label{sec:numerics}

The following checks were run for this article. They are
floating-point computations (double precision or \texttt{mpmath} at
$20$--$30$ digits), not certified, and they are not part of any proof.
The scripts are in \texttt{monotonicity/checks/} of the repository;
the zero lists are those of \cite[\S7]{WR}, i.e.\ the files of (T1).

\medskip\noindent\textbf{(1) Constants.} From the $511$ resp.\ $537$
ordinates, the closed form $C(\chi)=(\log\Lambda)''(\tfrac12,\chi)$,
and the enclosure \eqref{eq:Qenclose}: the bounds of
Remark~\ref{rem:M0values} and
\[
M_0(\chi_4)\in[3.68848,\,3.68931],\qquad
M_0(\chi_3)\in[4.48317,\,4.48459],
\]
which contain the ball-arithmetic enclosures \eqref{eq:M0cert} of the
direct certificate (those use a sharper tail bound). The roots of
$\sigma_M^2=2$ were found from the closed form of
Proposition~\ref{prop:WRvar}: $M_\sigma(\chi_4)=1.823915$ and
$M_\sigma(\chi_3)=2.127751$; the auxiliary file of the direct
certificate encloses them in the intervals $[1.82391477,1.82391478)$ and
$[2.12775095,2.12775096)$. The margins $\lambda_\chi(M,M')$ of
Remark~\ref{rem:Arange}(2) were computed from the listed zeros with
the tail bound stated there. For the pairs $(m,m')$ with
$m,m'\in\{0,1,2,3\}$ the same file gives
ball-arithmetic enclosures, for example
\begin{gather*}
\lambda_{\chi_4}(\tfrac12,\tfrac32)\in[1.0257439,1.0257440],\qquad
\lambda_{\chi_4}(\tfrac32,\tfrac52)\in[-0.0311712,-0.0311711],\\
\lambda_{\chi_3}(\tfrac32,\tfrac52)\in[0.0454118,0.0454119],
\end{gather*}
and $\sigma^2=1.0070099\ldots$ for $\chi_3$ at $m=1$.

\medskip\noindent\textbf{(2) The formula of Theorem~\ref{thm:A}(i).}
For $q=4$, $(m,m')=(0,1)$, we sieved the primes up to $10^7$ and
computed $I(x)=D_0(x)-D_1(x)/x$ directly (in double precision). We
compared it with (a) the right-hand side of \eqref{eq:formula} without
the $O$-term, truncated to the $511$ listed zeros, and (b) the
quantity
\[
I_{\mathrm{pred}}(x):=2\,\re\sum_{0<\gamma\le\Gamma}\int_2^xh(t)\,
t^{-1/2+i\gamma}\,dt+\sum_{k\ge2}\frac1k\sum_{p^k\le x}\chi(p)^kg(p^k),
\]
in the notation of the proof of Proposition~\ref{prop:formula}, the
integrals being evaluated in closed form by exponential integrals.
(b) is the formal one-integration-by-parts version of the explicit
formula, whose zero sum is only conditionally convergent; it differs
from $I(x)$ by boundary terms that are $O(1)$, and it tests every sign
and constant of Section~\ref{sec:formula} more sharply than (a).
Results:
\[
\begin{array}{r|rrr|r}
x & I(x) & I_{\mathrm{pred}}(x) & I-I_{\mathrm{pred}} &
I(x)\log x/(\kappa x^M)\\\hline
10^3 & 4.131 & 3.698 & 0.433 & 0.902\\
10^4 & 9.017 & 8.581 & 0.435 & 0.831\\
10^5 & 22.667 & 22.232 & 0.436 & 0.825\\
10^6 & 63.588 & 63.151 & 0.437 & 0.879\\
9.99\cdot10^6 & 170.097 & 169.643 & 0.454 & 0.867
\end{array}
\]
The difference $I-I_{\mathrm{pred}}$ is constant to within $0.02$
over four decades, as it should be. For (a), the discrepancy
$\bigl(I(x)-\text{main term}\bigr)\log x/(\kappa x^M)$ equals
$0.20$, $0.24$, $0.21$, $0.18$, $0.15$ at the five points, i.e.\
about $2.4/\log x$, consistent with the error term
$O(x^M/\log^2x)$. On a grid of $2\times10^4$ points in
$[10^4,10^7]$ the normalised quantity $I(x)\log x/(\kappa x^M)$
stays in $[0.753,0.961]$; the asymptotic range allowed by
\eqref{eq:bracketbounds} is
$[\lambda/2,\,1/(MM')-\lambda/2]=[0.513,0.821]$, and the excess over
$0.821$ is accounted for by the positive $O(1/\log x)$ correction. In
particular $I(x)>0$ throughout this range. Two further runs, $q=4$
with $(m,m')=(1,2)$ and $q=3$ with $(m,m')=(-\tfrac14,\tfrac12)$,
behaved in the same way (in the first, double-precision rounding of
sums of size $10^{20}$ limits the comparison to relative accuracy
$10^{-3}$; in the second, $M=\tfrac14$ is small and the $O(1/\log x)$
correction is correspondingly large, about $19/\log x$ in the
normalised quantity at $x=10^7$).

\medskip\noindent\textbf{(3) The coupling.} With the $511$ listed
zeros of $\chi_4$ (no tail) and $4\times10^5$ samples of the phase
vector we evaluated $X_M$ for several $M$ simultaneously. For the
pairs $(M,M')=(0.5,1.5)$, $(1.5,2.5)$, $(2.5,3.5)$, $(3,3.68)$,
$(1,5)$, $(2,5)$, for which the truncated $\eta$ is positive
($1.306$, $0.217$, $0.049$, $0.014$, $0.552$, $0.091$), there was no
sample with $X_{M'}>0\ge X_M$, and the sample minimum of
$X_M/M-X_{M'}/M'$ exceeded $\eta$ in every case, as
Lemma~\ref{lem:det} requires. For $(5,8)$ and $(8,12)$
($\eta=-0.174$, $-0.273$) there were $1416$ and $3658$ such samples,
in accordance with Proposition~\ref{prop:sharpB}; for $(3.68,4.5)$
and $(3.6,5)$ ($\eta=-0.013$, $-0.031$) the sample minimum of
$X_M/M-X_{M'}/M'$ was negative ($-0.007$, $-0.019$) but no violation
of the nesting occurred among the $4\times10^5$ samples, the
probability of a violation being positive but small.

\medskip\noindent\textbf{(4) Theorem~\ref{thm:C}(i).} For $q=4$,
$m\in\{0,1\}$ and $b\in\{-1,0,2,3\}$ we computed
$(\cE_{m,b}(x)-\cE_{m,1}(x))\log x$ from the same sieve at the five
values of $x$ of the table above. The values stay bounded, as
Theorem~\ref{thm:C}(i) requires; they lie in the following intervals.
\[
\begin{array}{c|cccc}
 & b=-1 & b=0 & b=2 & b=3\\\hline
m=1 & [1.25,1.87] & [0.54,0.88] & [-0.82,-0.42] & [-1.57,-0.74]\\
m=0 & [8.1,12.6] & [2.84,3.01] & [-2.01,-1.29] & [-3.54,-2.07]
\end{array}
\]
(For $m=0$ the weight $M=\tfrac12$ is small and the values for
$b=-1,2,3$ are still drifting slowly at $x=10^7$, as lower-order
terms in $1/\log x$ would produce.) The signs agree with
\eqref{eq:Cdecomp}.

\medskip\noindent\textbf{(5) Result files of the certificates.} The
script \texttt{parse\_cert.py} reads a result file of either
certificate with exact rational endpoints and checks that every cell
is marked as certified with a bound of the right sign, that
consecutive cells share their endpoints, and that the file covers the
stated range; it reports the number of cells and the weakest bound.
Run on the two files of the envelope certificate it reproduces the
table of Section~\ref{subsec:envcert} ($765$ cells each, contiguous
from $1$ to $2000$, smallest lower bounds $1.318804\times10^{-3}$ and
$1.353296\times10^{-3}$); run on the four files of the direct
certificate it reproduces the table of
Section~\ref{subsec:confirm}(a) ($317$, $293$, $296$, $292$ cells;
ranges $[3.6,17.4]$, $[17.4,320]$, $[4.4,19.2]$, $[19.2,350]$). This
checks what the files say, not that the programs which wrote them are
correct.

\medskip\noindent\textbf{(6) Constants of Theorems~\ref{thm:E}
and~\ref{thm:general}.} The script \texttt{thmE\_constants.py}
recomputes, at $30$ digits, the quantities of \eqref{eq:Evalues} and
of the proof of Theorem~\ref{thm:general}(iv):
$z_*=1.7971492\ldots$, $\tau=0.4492873\ldots$,
$T_c=8.2662190\ldots$, $\beta_*=0.0255529\ldots$; the left-hand side
of (E2) is $0.01204819\ldots$; the two candidates for $\sigma_0^2$ are
$236.971101\ldots$ for $(a,q)=(1,3)$ and $275.837\ldots$ for $(0,5)$;
the two bounds of (E4) are $11.37517\ldots$ and $14.21016\ldots$;
$G(\sigma_0^2)=0.32197151\ldots$; and the lower bound for
$Q_\chi(40)$ is $38.51\ldots$. All agree with the certified values.

\section{Gaps and caveats}\label{sec:gaps}

\begin{enumerate}
\item \emph{All refereeing so far is internal and AI-conducted.} This
article, its proofs, the certifying programs, the audits of those
programs and the referee reports summarised in
Appendix~\ref{sec:review} were all produced by AI agents under the
author's direction. No part of it has been examined by an external
referee. The reports found errors of presentation and of rounding
and hypotheses that had to be made explicit (all repaired or stated
here), and found no error in a proof; this is evidence of care, not a
substitute for independent human review.
\item \emph{Nothing is formally verified beyond the Lean core.} The
directory \texttt{monotonicity/lean/} of the repository contains a
Lean~4 (Mathlib) formalization, without \texttt{sorry}, of the
following and nothing else: the deterministic coupling inequality of
Lemma~\ref{lem:det}; the criterion ``$Q\le1$ on $[M,M']$ implies
$\eta(M,M')\ge0$'', for finite sums and for series; and the
measure-theoretic step from the pathwise inequality to the inequality
between probabilities used in Theorem~\ref{thm:B}. It contains
nothing about $L$-functions, zeros, GRH, LI, limiting distributions,
Bessel functions, Lemma~\ref{lem:diff}, Theorems~\ref{thm:envelope}
and~\ref{thm:E}, or any certificate. Every other statement of this
article rests on the written proofs and, where indicated, on
computation.
\item \emph{The certificates certify soundness given the listed
inputs, and no more.} Theorem~\ref{thm:q34} holds provided
(T1)--(T3) of Section~\ref{subsec:inputs} are correct, and
Theorem~\ref{thm:nonmono} provided (N1)--(N5) of
Section~\ref{sec:nonmono} are. In both cases this includes the
correctness of FLINT/Arb, the correctness of the certifying programs
(Python programs, audited internally but not formally verified), and the identity of the zero files:
both audits showed by mutation that a zero file with an ordinate
removed produced a false ``certified'' output without any error, which
is why the digests of the files are part of the hypotheses (the
programs of the non-monotonicity certificate have since been changed
to refuse such a file). A
certificate of this kind excludes rounding error; it does not exclude
an error in the mathematics that the program implements, which is why
that mathematics is written out in
Sections~\ref{sec:envelope}--\ref{sec:largeweight} and was re-derived
by the internal referees. Theorem~\ref{thm:E}, and with it parts
(iii) and (iv) of the General Theorem, is not free of computation
either: its proof ends with the four numerical inequalities
\eqref{eq:Evalues} between explicit values of elementary functions,
$\psi$ and $J_0$. They were evaluated in ball arithmetic and
recomputed in floating point at $30$ digits
(Section~\ref{sec:numerics}(6)); they involve no zero data and could
be checked by hand, which we have not done.
\item \emph{What was audited and what was not.} Audited: the direct
certificate on $[3.6,17.4]$ and $[4.4,19.2]$ (and independently
re-proved there); the analytic leg of \cite{WR} on $M\ge300$; the
envelope certificate and the proof of Theorem~\ref{thm:E} (by a
referee who re-derived the estimates, recomputed the constants, re-ran
the certificate in full and evaluated the exact cell quantity
independently at thirteen cells per modulus, on the version of the
program preceding the last one; see Appendix~\ref{sec:review}); the
non-monotonicity certificate. Not separately audited: the stretch
files of the direct certificate ($[17.4,320]$ and $[19.2,350]$), which
is one reason why the direct certificate is presented as a
confirmation and not as a link in the proof; and the changes made to
the certifying programs after their audits in response to the audit
findings.
\item \emph{The variance identity for general $\chi$.}
Proposition~\ref{prop:WRvar} is proved in \cite{WR} for $\chi_4$, with
the statement that the proof applies verbatim to every real primitive
$\chi$ with $L(\tfrac12,\chi)\ne0$. The ingredients are that the
completed $L$-function is entire of order one, real on the real axis,
with root number $+1$, together with GRH and NCZ; three of the
internal reports re-derived the identity, two of them for characters
other than $\chi_3$, $\chi_4$. Theorems~\ref{thm:E}, \ref{thm:general}
and~\ref{thm:nonmono} use it for characters other than $\chi_3$,
$\chi_4$.
\item \emph{Quoted limiting-distribution theory.}
Theorems~\ref{thm:B}(iii) and~\ref{thm:C}(ii)--(iii),
Lemma~\ref{lem:joint}, Corollaries~\ref{cor:sharp}(b)
and~\ref{cor:nonmono} and the race forms of the three headline
theorems rest on \cite[Theorem~1, Proposition~1]{WR} and, through
them and directly in Lemma~\ref{lem:joint}, on
\cite[Corollary~1.3(a), Theorem~1.9]{ANS14}. We did not re-derive
these. When the article was written we relied on the statements of
\cite{ANS14} as reproduced in \cite{WR}; an internal referee has since
checked them against the arXiv version of \cite{ANS14} and found them
to match. The two-dimensional statement of Lemma~\ref{lem:joint} is
obtained from the one-dimensional results by the Cram\'er--Wold device
(it could also be quoted from \cite[Theorems~1.4, 1.9]{ANS14}).
\item \emph{What $\delta_\chi$ means without LI.} Under GRH$(\chi)$
and NCZ$(\chi)$ alone, $\hd_\chi(M)=\PP(X_M>0)$ is a quantity of the
random model with independent phases; Theorems~\ref{thm:B}(i), (ii),
(iv), Lemma~\ref{lem:diff}, Theorems~\ref{thm:envelope},
\ref{thm:E}, \ref{thm:general}, \ref{thm:q34} and~\ref{thm:nonmono}
are statements about it. Its identification with the logarithmic
density of the race needs LI$(\chi)$. Moreover Theorem~\ref{thm:B}
orders densities; it does not nest the lead sets $\{x:D_m(x)>0\}$
(Remark~\ref{rem:Arange}(3)).
\item \emph{What is not proved about the intermediate range.} For a
general real primitive $\chi$ nothing is proved on
$(M_0(\chi),40)$, and by Theorem~\ref{thm:nonmono} nothing of the
form ``$\hd_\chi$ is decreasing there'' can be. The constant $40$ of
Theorem~\ref{thm:E} is not optimised beyond what is said in
Remark~\ref{rem:Enotes}, and we do not know the least $M_1$ such that
$\hd_\chi$ is decreasing on $[M_1,\infty)$ for every $\chi$; the
example of conductor $232$ shows $M_1>1$, since
$\hd(1)<\hd(\tfrac{101}{100})$ there. The non-monotonicity
theorem concerns two characters only; it asserts inequalities between
values at six points each and nothing about the derivative or the
extrema.
\item \emph{Numerical values.} The values of $M_0$, $M_\sigma$ and
$\lambda_\chi$ in Remarks~\ref{rem:M0values}, \ref{rem:Arange} and
Section~\ref{sec:numerics} are floating-point evaluations with
rigorous tail bounds but without interval arithmetic, except where a
ball-arithmetic value from a certificate's file is quoted. They enter
no proof: Theorems~\ref{thm:A}--\ref{thm:C} are stated in terms of
$M_0(\chi)$, $\lambda_\chi$, $\sigma_M$. Corollary~\ref{cor:equiv}
for the specific pair $(0,1)$ uses $\sigma^2<2$ at $m=1$, where the
margin is large. The proof of Theorem~\ref{thm:q34} uses the single
certified inequality \eqref{eq:sig101} to place $M_0(\chi_q)$ above
$1.01$; it does not use the enclosures \eqref{eq:M0cert}.
\item \emph{The size of $M_0(\chi)$ for large $q$} is not determined
(Remark~\ref{rem:corner}). The proved bounds of this article are
those of Proposition~\ref{prop:M0} and $M_0(\chi)<40$
(Theorem~\ref{thm:general}); the sharper bounds and the numerical
values in Remark~\ref{rem:corner} are quoted from a repository note
and are not claimed as results of this article.
\item \emph{Effectivity.} The threshold $x_0$ of
Theorem~\ref{thm:A}(ii) is effective in principle but was not made
explicit; the numerical evidence of Section~\ref{sec:numerics}
suggests it is small for $q=4$, $(m,m')=(0,1)$, but we prove nothing
about its size.
\item \emph{Scope of Theorem~\ref{thm:A}.} It is a statement about
pairs of weights with $\lambda_\chi(M,M')>0$; for $q=4$ this excludes
e.g.\ $(m,m')=(1,2)$. Without LI it gives monotonicity of upper and
lower densities only on $M'\le M_\sigma(\chi)$, and only in the
non-strict sense. All results other than Theorem~\ref{thm:A} and
Corollaries~\ref{cor:densities}, \ref{cor:sharp}(a),
\ref{cor:landau} assume NCZ$(\chi)$.
\item \emph{Corollary~\ref{cor:landau}} gives GRH$(\chi)$ only in the
absence of real zeros in $(\tfrac12,1)$
(Remark~\ref{rem:exceptional}).
\item \emph{Literature and bibliography.} The works cited in
Remark~\ref{rem:literature} were checked at different levels. For
\cite{PT19,Was24,BBLMR24} the relevant statements were read in the
arXiv texts; for \cite{Joh23,Zhao25,Suz25,Hay25} the statements were
confirmed at the level of the arXiv abstracts (by an internal
referee); \cite{LWS26} was seen at abstract level only;
\cite{HL18,Lan18,Fuj88} are cited as they appear in \cite{PT19}, and
their originals were not consulted; \cite{GW18} is cited as it is
described in \cite{Was24}. Entries of the reference list whose
bibliographic details have not been verified against the published
source carry a visible \textsf{TODO-cite} mark. The entries
\cite{AK23,Sheth24,SK25,Aym22,BMS26,BMOR21,Wat44,Bir48,Pro65,KS99} are
reused from the reference list of \cite{WR} and were not re-checked
for this article. The literature search behind the novelty remarks of
Section~\ref{sec:intro} was limited and is recorded in the repository;
``not found'' means not found by that search. In particular we did
not search specifically for earlier examples of non-monotone
dependence of a bias density on a weight parameter.
\item \emph{The announcement mentioned in
Remark~\ref{rem:GRHstatus}} is not used anywhere, and we examined
neither the note nor the formalization.
\item \emph{No claim is made} for $m\le-\tfrac12$, for non-real
characters, or for races with more than two competitors.
\item \emph{Differences from the working notes that preceded this
article}: the oscillating sum in \eqref{eq:formula-normalised} enters
with a plus sign (the notes had a minus sign; immaterial for the
results, which use only its absolute value);
Lemma~\ref{lem:diff} needs nine ordinates and no hypothesis on the
weight, and replaces the use of \cite[Lemma~21]{WR}; in the proof of
Lemma~\ref{lem:J} the value of $J_0$ at $2\sqrt2$ is bounded by (B4)
instead of numerically; and the bound for $h'(40)$ in the proof of
Theorem~\ref{thm:general}(iv) is ours ($0.01297$, giving
$Q_\chi(40)>38$ as in the notes).
\end{enumerate}

\appendix
\section{Review history}\label{sec:review}

All reviewing recorded here was done by AI agents instructed to act as
adversarial referees or auditors; the reports are in the directory
\texttt{monotonicity/} of the repository under the names given.
``R'' items were classified by the referees as repairable, ``M'' items
as minor, ``G'' items as gaps (a hypothesis that had to be made
explicit).

\subsection*{Round 1: the first complete draft (reports
\texttt{REFEREE-1.md}, \texttt{REFEREE-2.md})}
The first draft contained Sections~\ref{sec:prelim}--\ref{sec:thmC}
essentially as they are now and, in place of
Sections~\ref{sec:envelope}--\ref{sec:nonmono}, a reduction of full
monotonicity for $q=3,4$ to three hypotheses, to be supplied by the
direct certificate and by \cite{WR}. Neither referee found a fatal
error; all displayed identities, signs and constants of Sections~3--7
were re-derived by them.

\emph{Referee 1} (coupling, Theorem~\ref{thm:B}, the reduction).
R1: the reduction was a conditional statement and the title, abstract
and introduction overstated it; resolved at the time by separating a
pure implication from a computer-assisted statement with an explicit
list of what is taken on trust (this structure survives in
Theorem~\ref{thm:q34} and Section~\ref{subsec:inputs}). R2: the
description of the status of \cite[Corollary~7]{WR} was inaccurate,
and the direct certificate on the large range was not used; resolved
at the time by adopting the stretch certificates, and now moot, since
neither is used in a proof. M1: stale enclosures of $M_0$ and a
circular ``agreement'' sentence; resolved (\eqref{eq:M0cert} carries
the values of the auxiliary file). M2: ``comonotone coupling'' was the
wrong term; now ``equal-phase coupling'' throughout, with the
referee's sentence on the quantile coupling after
Proposition~\ref{prop:sharpB}. M3: the differentiability lemma cited
the proof of a lemma whose hypothesis it removes; the proof of
Lemma~\ref{lem:diff} is self-contained. M4: the companion's existence
theorem also assumes $L(\tfrac12,\chi)\ne0$; stated. M5: the meaning
of $\delta_\chi$ without LI; Section~\ref{subsec:setting} and item~7
of Section~\ref{sec:gaps}. M6: Theorem~\ref{thm:B}(iii) could be
misread as nesting of the lead sets; wording changed, and the example
in Remark~\ref{rem:Arange}(3) added. M7: small corrections in
Remark~\ref{rem:M0values} and in the wording ``nested almost
surely''.

\emph{Referee 2} (Theorems~\ref{thm:A}, \ref{thm:C}, the Landau
converse, bibliography). R1: a wrong citation of Fujii, the two parts
of Landau's paper, and the date of Hardy--Littlewood; corrected in the
reference list following the reference list of \cite{PT19}, the
entries remaining marked because the originals were not consulted.
R2: the remark on the status of GRH asserted a machine-checking as a
fact and listed two inaccurate consequences;
Remark~\ref{rem:GRHstatus} uses the referee's wording. R3: the
abstract stated the equivalence with GRH under a hypothesis defined
only under GRH; Corollary~\ref{cor:equiv}(c) is stated with the
unconditional quantity $4M'\frac{\Lambda'}{\Lambda}(M'+\frac12,\chi)$.
R4: a sentence on ``a limited range of the smoothing parameter''
overgeneralised; restricted (Remark~\ref{rem:literature}). Minor
items: Lemma~\ref{lem:joint} can be quoted directly from
\cite{ANS14} (noted); the statements quoted from \cite{ANS14} were
checked by the referee against the arXiv version; the novelty sentence
on Proposition~\ref{prop:formula}(a) was toned down; unverified
bibliographic details are marked.

At the same time the heuristic $M_0(\chi)\approx\pi/(2\log q)$ of the
first draft was withdrawn (Remark~\ref{rem:corner}), and the two-zero
model (Proposition~\ref{prop:twozero}) was added.

\subsection*{Round 2: the derivative, the envelope criterion, the
certificates}

\emph{Referee 3} (report \texttt{REFEREE-3.md}, on the note
\texttt{BRIDGE-ATTEMPT.md}, which is the source of
Lemma~\ref{lem:diff}, Lemmas~\ref{lem:J} and~\ref{lem:L},
Proposition~\ref{prop:C0}, Theorem~\ref{thm:envelope},
Corollary~\ref{cor:cell} and Proposition~\ref{prop:twozero}). The
referee re-derived these statements, re-implemented the criterion
independently, and found no fatal error. Two repairable items. R1:
the note's corollary ``the derivative is negative whenever
$K:=\mathrm{NUM}/(\mathrm{MAIN}-\mathrm{TAILF})<1$'' lacked the sign
condition $\mathrm{MAIN}>\mathrm{TAILF}$; here the criterion is
stated as $\mathrm{MAIN}-\mathrm{TAILF}-\mathrm{NUM}>0$ and $K$
appears only in Remark~\ref{rem:K}, with the sign condition. R2
(presentation): a table of point values in the note had been computed
with quantities other than those of the theorem; no such table is
reproduced here. The referee also reproduced, in floating point, the
numerical evidence for conductors $163$ and $232$ that preceded the
certificate of Section~\ref{sec:nonmono}, noting that the case of
conductor $232$ is sensitive to how the unlisted zeros are modelled;
the certificate bounds that model error rigorously.

\emph{Audit of the direct certificate} (report
\texttt{cert-audit/AUDIT.md}). No break of rigour. Three gaps:
G1, differentiability of the density at weights below those covered
by \cite[Lemma~21]{WR} (supplied by Lemma~\ref{lem:diff}); G2, the
certifying program accepts any zero file (met by the digests in (T1));
G3, provenance of the result files (closed by the audit's replay,
which reproduced the $613$ result lines). Eleven minor items. An
independent program re-proved the negativity of the derivative on all
$613$ cells of $[3.6,17.4]$ and $[4.4,19.2]$. The stretch files were
not audited. See Section~\ref{subsec:confirm}(a).

\emph{Audit of the analytic leg of \cite{WR}} (report
\texttt{ANALYTIC-LEG-AUDIT.md}). No fatal error; the statement
``under GRH$(\chi_q)$, $q\in\{3,4\}$, the density is differentiable
on $(299,\infty)$ with negative derivative for all $M\ge300$'' was
confirmed, with $90$ ball-arithmetic checks. Three repairable items
concern the text of \cite{WR}, not this article. See
Section~\ref{subsec:confirm}(c).

\emph{Referee 5} (report \texttt{REFEREE-5.md}) examined the envelope
certificate and the note that is the source of
Theorem~\ref{thm:E}. Every step of the proof of
Theorem~\ref{thm:E} was re-derived, the constants were recomputed
independently, the certifying program was read line by line and
subjected to mutation tests, and a full fresh run reproduced both
result files. No fatal error. R1: the program for the constants of
Theorem~\ref{thm:E} printed its bounds rounded to nearest, so that
several displayed inequalities of the note were false in the last
digit (for example ``$T_c\le8.2662$''); the decisions of the program,
made on the balls, and the constant $0.3219$ were not affected. This
article uses the values with directed rounding, \eqref{eq:Evalues}.
Thirteen minor items, among them: sub-range maxima of the ratio $K$
quoted in the notes omitted one cell (not quoted here); a statement
on the loss of the bound held only on part of the range (likewise not
quoted here); ``the smallest threshold $38.5$'' is the best
on a searched grid only, and the convention on multiple zeros should
be stated (Remark~\ref{rem:Enotes}); the result files should record
digests (done in the re-run). The referee also observed that
the proof of Theorem~\ref{thm:E} gives $M_0(\chi)<40$, which is
Theorem~\ref{thm:general}(iv). After this report the envelope
certificate was re-run with a revised program that records the
digests of the program and of the zero file in the result files; the
result lines quoted in Section~\ref{subsec:envcert} are those of the
re-run and agree with the earlier ones in every figure quoted here.
The revised program itself was not re-refereed.

\emph{Audit of the non-monotonicity certificate} (report
\texttt{nonmonotone-audit/AUDIT.md}). No break of rigour; one gap
(the identity of the zero files, met by the digest clause of (N5));
six minor items; all twelve enclosures confirmed by independent
rigorous enclosures. See the end of Section~\ref{sec:nonmono}. After
the audit the certifying programs of that directory were revised as
it recommended (binding of the zero files by digest) and the result
file was regenerated, with identical enclosures; the revised programs
were not re-audited.

\subsection*{What changed in this article in round 2}
The reduction of the first draft and its hypotheses were removed. The
article now contains the derivative lemma and the envelope criterion
with proofs (Section~\ref{sec:envelope}), the uniform large-weight
theorem and the General Theorem (Section~\ref{sec:largeweight}), the
theorem for $q=3,4$, whose computational links are the envelope
certificate and the four evaluated inequalities of
Theorem~\ref{thm:E}, with the direct certificate and the analytic leg
of \cite{WR} as confirmations (Section~\ref{sec:main}), and the
non-monotonicity theorem (Section~\ref{sec:nonmono}). The title,
abstract and introduction were rewritten accordingly.

\begin{thebibliography}{99}
\bibitem{WR} R. Sharma, \emph{Rubinstein--Sarnak densities for
weighted prime number races}, preprint, August 2026,
\href{https://doi.org/10.5281/zenodo.21947101}{doi:10.5281/zenodo.21947101};
code and data: \url{https://github.com/raghavs1729/weighted-prime-races}.
\bibitem{RS94} M. Rubinstein, P. Sarnak, \emph{Chebyshev's bias},
Experiment.\ Math.\ \textbf{3} (1994), 173--197.
\bibitem{ANS14} A. Akbary, N. Ng, M. Shahabi, \emph{Limiting
distributions of the classical error terms of prime number theory},
Q. J. Math.\ \textbf{65} (2014), 743--780. arXiv:1306.1657.
\todocite{statements checked against the arXiv version; journal pages
not checked}
\bibitem{AK23} M. Aoki, S. Koyama, \emph{Chebyshev's bias against
splitting and principal primes in global fields}, J. Number Theory
\textbf{245} (2023), 233--262. arXiv:2203.12266.
\todocite{entry reused from \cite{WR}; not re-checked}
\bibitem{Sheth24} A. Sheth, \emph{Euler products at the centre and
applications to Chebyshev's bias}, Math.\ Proc.\ Cambridge Philos.\
Soc.\ \textbf{179} (2025), 331--349. arXiv:2405.01512.
\todocite{entry reused from \cite{WR}; not re-checked}
\bibitem{SK25} K. Shimada, S. Koyama, \emph{Weighted prime number
theorem on arithmetic progressions with refinements}, Mathematics
\textbf{13} (2025), no.\ 16, 2564.
\todocite{entry reused from \cite{WR}; not re-checked}
\bibitem{Hay25} M. Hayani, \emph{Chebyshev's bias without linear
independence}, arXiv:2512.23302 (2025).
\bibitem{Aym22} M. Aymone, \emph{A note on prime number races and zero
free regions for $L$ functions}, Int.\ J. Number Theory \textbf{18}
(2022). arXiv:2001.00764.
\todocite{entry reused from \cite{WR}; not re-checked}
\bibitem{BMS26} S. Bhattacharya, G. Martin, R.~M. Simpson,
\emph{Correlations of error terms for weighted prime counting
functions}, arXiv:2507.13504, to appear in Acta Arith.
\todocite{entry reused from \cite{WR}; not re-checked}
\bibitem{GW18} J. Gerard, L.~C. Washington, \emph{Sums of powers of
primes}, Ramanujan J. \textbf{45} (2018), 171--180.
\todocite{not verified; cited as described in \cite{Was24}}
\bibitem{Was24} L.~C. Washington, \emph{Sums of powers of primes II},
Ramanujan J. \textbf{65} (2024), 783--795. arXiv:2209.12845.
\todocite{Theorem~3 read in the arXiv text; journal reference not
verified}
\bibitem{BBLMR24} M. Boran, J. Byun, Z. Li, S.~J. Miller, S. Reyes,
\emph{Sums of powers of primes in arithmetic progression}, PUMP J.
Undergrad.\ Res.\ \textbf{7} (2024), 29--50. arXiv:2309.16007.
\bibitem{HL18} G.~H. Hardy, J.~E. Littlewood, \emph{Contributions to
the theory of the Riemann zeta-function and the theory of the
distribution of primes}, Acta Math.\ \textbf{41} (1917), 119--196.
\todocite{details as in the reference list of \cite{PT19}, which
gives the year 1917 (the volume is also cited as 1916 or 1918);
original not consulted}
\bibitem{Lan18} E. Landau, \emph{\"Uber einige \"altere Vermutungen
und Behauptungen in der Primzahltheorie}, I and II, Math.\ Z.
\textbf{1} (1918), 1--24 and 213--219. \todocite{details as in the
reference list of \cite{PT19}; originals not consulted}
\bibitem{Fuj88} A. Fujii, \emph{Some generalizations of Chebyshev's
conjecture}, Proc.\ Japan Acad.\ Ser.~A Math.\ Sci.\ \textbf{64}
(1988), no.~7, 260--263. \todocite{details as in the reference list
of \cite{PT19}; original not consulted}
\bibitem{PT19} D.~J. Platt, T.~S. Trudgian, \emph{Fujii's development
on Chebyshev's conjecture}, Int.\ J. Number Theory \textbf{15}
(2019), 639--644. arXiv:1802.00089. \todocite{arXiv text read; journal
pages not checked}
\bibitem{Joh23} D.~R. Johnston, \emph{On the average value of
$\pi(t)-\mathrm{li}(t)$}, Canad.\ Math.\ Bull.\ \textbf{66} (2023),
185--195. arXiv:2201.06184. \todocite{statement confirmed from the
arXiv abstract; journal reference not verified}
\bibitem{Zhao25} T. Zhao, \emph{On the mean values of the error terms
in Mertens' theorems}, Res.\ Number Theory \textbf{11} (2025),
paper~62. arXiv:2411.18903.
\bibitem{Suz25} M. Suzuki, \emph{On variants of Chebyshev's
conjecture}, Ramanujan J. \textbf{68} (2025), article~95.
arXiv:2411.07436.
\bibitem{LWS26} Liu, Wang, Shi, \emph{Two problems in comparative
prime number theory: the five-way race modulo $11$, and the sign of
the integrated prime counting error}, Zenodo,
doi:10.5281/zenodo.22208026 (2026). \todocite{authors' initials;
seen at abstract level only}
\bibitem{Dav00} H. Davenport, \emph{Multiplicative Number Theory},
3rd ed., revised by H.~L. Montgomery, Graduate Texts in Mathematics
\textbf{74}, Springer, New York, 2000.
\bibitem{MV07} H.~L. Montgomery, R.~C. Vaughan, \emph{Multiplicative
Number Theory I: Classical Theory}, Cambridge Studies in Advanced
Mathematics 97, Cambridge University Press, Cambridge, 2007.
\bibitem{BMOR21} M.~A. Bennett, G. Martin, K. O'Bryant,
A. Rechnitzer, \emph{Counting zeros of Dirichlet $L$-functions},
Math.\ Comp.\ \textbf{90} (2021), 1455--1482.
\bibitem{HW} G.~H. Hardy, E.~M. Wright, \emph{An Introduction to the
Theory of Numbers}, 5th ed., Oxford University Press, Oxford, 1979.
\bibitem{Fel71} W. Feller, \emph{An Introduction to Probability Theory
and Its Applications}, Vol.~II, 2nd ed., Wiley, New York, 1971.
\bibitem{Bil99} P. Billingsley, \emph{Convergence of Probability
Measures}, 2nd ed., Wiley, New York, 1999.
\bibitem{Wat44} G.~N. Watson, \emph{A Treatise on the Theory of Bessel
Functions}, 2nd ed., Cambridge University Press, Cambridge, 1944.
\bibitem{Bir48} Z.~W. Birnbaum, \emph{On random variables with
comparable peakedness}, Ann.\ Math.\ Statist.\ \textbf{19} (1948),
76--81.
\bibitem{Pro65} F. Proschan, \emph{Peakedness of distributions of
convex combinations}, Ann.\ Math.\ Statist.\ \textbf{36} (1965),
1703--1706.
\bibitem{KS99} N.~M. Katz, P. Sarnak, \emph{Zeroes of zeta functions
and symmetry}, Bull.\ Amer.\ Math.\ Soc.\ \textbf{36} (1999), 1--26.
\bibitem{Alp26} L. Alp\"oge, note dated 20 November 2026,
\url{https://alpo.ge/remark.pdf}. \todocite{not examined by the
author of this article; title to be added; re-check before
submission}
\bibitem{Lean26} Lean FRO and Mathlib Initiative, joint statement of
21 November 2026,
\url{https://leanprover-community.github.io/blog/posts/riemann-hypothesis/}.
\todocite{not examined by the author of this article; re-check
before submission}
\end{thebibliography}

\end{document}
